\documentclass[11pt]{article}
\usepackage[margin=1in]{geometry}
\usepackage{amsmath,amssymb,amsthm,mathtools,bm}
\usepackage[numbers]{natbib}
\usepackage{hyperref,booktabs,enumitem}
\newtheorem{theorem}{Theorem}
\newtheorem{lemma}[theorem]{Lemma}
\newtheorem{proposition}[theorem]{Proposition}
\newtheorem{corollary}[theorem]{Corollary}
\newtheorem{conjecture}[theorem]{Conjecture}
\theoremstyle{definition}
\newtheorem{assumption}{Assumption}
\newtheorem{definition}[theorem]{Definition}
\newtheorem{remark}[theorem]{Remark}
\newtheorem{example}[theorem]{Example}
\title{Theory: where the input should enter a learned KKL observer,\\ why a short input window suffices, and what learned observers can guarantee}
\date{draft, generated from theory/m1--m3}
% numbers produced by the mechanical results layer (analysis.paper_results); fallbacks keep the standalone build working
\IfFileExists{../paper/generated/numbers.tex}{% generated by analysis/paper_results.py -- every number quoted in the text. Do not edit.
% generated 2026-10-26 17:43

\providecommand{\ConeConcatBeatsAllId}{}\renewcommand{\ConeConcatBeatsAllId}{0/5}
\providecommand{\ConeConcatBeatsAllSeq}{}\renewcommand{\ConeConcatBeatsAllSeq}{0/5}
\providecommand{\ConeConcatBeatsAllShape}{}\renewcommand{\ConeConcatBeatsAllShape}{0/5}
\providecommand{\ConeConcatBeatsLatentId}{}\renewcommand{\ConeConcatBeatsLatentId}{1/5}
\providecommand{\ConeConcatBeatsLatentSeq}{}\renewcommand{\ConeConcatBeatsLatentSeq}{1/5}
\providecommand{\ConeConcatBeatsLatentShape}{}\renewcommand{\ConeConcatBeatsLatentShape}{1/5}
\providecommand{\ConeConcatLosesLatentId}{}\renewcommand{\ConeConcatLosesLatentId}{3/5}
\providecommand{\ConeConcatLosesLatentSeq}{}\renewcommand{\ConeConcatLosesLatentSeq}{3/5}
\providecommand{\ConeConcatLosesLatentShape}{}\renewcommand{\ConeConcatLosesLatentShape}{3/5}
\providecommand{\ConeConcatVsLatentRangeId}{}\renewcommand{\ConeConcatVsLatentRangeId}{-11 to -11\%}
\providecommand{\ConeConcatVsLatentRangeSeq}{}\renewcommand{\ConeConcatVsLatentRangeSeq}{-8 to -8\%}
\providecommand{\ConeConcatVsLatentRangeShape}{}\renewcommand{\ConeConcatVsLatentRangeShape}{-10 to -10\%}
\providecommand{\ConeDenseAblationBeatsHyperKKLId}{}\renewcommand{\ConeDenseAblationBeatsHyperKKLId}{1/5}
\providecommand{\ConeDenseAblationBeatsHyperKKLSeq}{}\renewcommand{\ConeDenseAblationBeatsHyperKKLSeq}{3/5}
\providecommand{\ConeDenseAblationBeatsHyperKKLShape}{}\renewcommand{\ConeDenseAblationBeatsHyperKKLShape}{1/5}
\providecommand{\ConeDenseAblationBeatsLatentId}{}\renewcommand{\ConeDenseAblationBeatsLatentId}{5/5}
\providecommand{\ConeDenseAblationBeatsLatentSeq}{}\renewcommand{\ConeDenseAblationBeatsLatentSeq}{3/5}
\providecommand{\ConeDenseAblationBeatsLatentShape}{}\renewcommand{\ConeDenseAblationBeatsLatentShape}{4/5}
\providecommand{\ConeDenseAblationLosesLatentId}{}\renewcommand{\ConeDenseAblationLosesLatentId}{0/5}
\providecommand{\ConeDenseAblationLosesLatentSeq}{}\renewcommand{\ConeDenseAblationLosesLatentSeq}{1/5}
\providecommand{\ConeDenseAblationLosesLatentShape}{}\renewcommand{\ConeDenseAblationLosesLatentShape}{0/5}
\providecommand{\ConeDenseAblationVsLatentRangeId}{}\renewcommand{\ConeDenseAblationVsLatentRangeId}{-13 to -67\%}
\providecommand{\ConeDenseAblationVsLatentRangeSeq}{}\renewcommand{\ConeDenseAblationVsLatentRangeSeq}{-34 to -62\%}
\providecommand{\ConeDenseAblationVsLatentRangeShape}{}\renewcommand{\ConeDenseAblationVsLatentRangeShape}{-24 to -69\%}
\providecommand{\ConeFilmBeatsAllId}{}\renewcommand{\ConeFilmBeatsAllId}{0/5}
\providecommand{\ConeFilmBeatsAllSeq}{}\renewcommand{\ConeFilmBeatsAllSeq}{0/5}
\providecommand{\ConeFilmBeatsAllShape}{}\renewcommand{\ConeFilmBeatsAllShape}{0/5}
\providecommand{\ConeFilmBeatsLatentId}{}\renewcommand{\ConeFilmBeatsLatentId}{3/5}
\providecommand{\ConeFilmBeatsLatentSeq}{}\renewcommand{\ConeFilmBeatsLatentSeq}{2/5}
\providecommand{\ConeFilmBeatsLatentShape}{}\renewcommand{\ConeFilmBeatsLatentShape}{3/5}
\providecommand{\ConeFilmLosesLatentId}{}\renewcommand{\ConeFilmLosesLatentId}{2/5}
\providecommand{\ConeFilmLosesLatentSeq}{}\renewcommand{\ConeFilmLosesLatentSeq}{3/5}
\providecommand{\ConeFilmLosesLatentShape}{}\renewcommand{\ConeFilmLosesLatentShape}{2/5}
\providecommand{\ConeFilmVsLatentRangeId}{}\renewcommand{\ConeFilmVsLatentRangeId}{-44 to -59\%}
\providecommand{\ConeFilmVsLatentRangeSeq}{}\renewcommand{\ConeFilmVsLatentRangeSeq}{-24 to -50\%}
\providecommand{\ConeFilmVsLatentRangeShape}{}\renewcommand{\ConeFilmVsLatentRangeShape}{-35 to -60\%}
\providecommand{\ConeHyperKKLBeatsAllId}{}\renewcommand{\ConeHyperKKLBeatsAllId}{4/5}
\providecommand{\ConeHyperKKLBeatsAllSeq}{}\renewcommand{\ConeHyperKKLBeatsAllSeq}{2/5}
\providecommand{\ConeHyperKKLBeatsAllShape}{}\renewcommand{\ConeHyperKKLBeatsAllShape}{2/5}
\providecommand{\ConeHyperKKLBeatsDenseAblationId}{}\renewcommand{\ConeHyperKKLBeatsDenseAblationId}{1/5}
\providecommand{\ConeHyperKKLBeatsDenseAblationSeq}{}\renewcommand{\ConeHyperKKLBeatsDenseAblationSeq}{1/5}
\providecommand{\ConeHyperKKLBeatsDenseAblationShape}{}\renewcommand{\ConeHyperKKLBeatsDenseAblationShape}{1/5}
\providecommand{\ConeHyperKKLBeatsLatConcId}{}\renewcommand{\ConeHyperKKLBeatsLatConcId}{5/5}  % HyperKKL (r=4) beats latent AND concat, CI
\providecommand{\ConeHyperKKLBeatsLatConcSeq}{}\renewcommand{\ConeHyperKKLBeatsLatConcSeq}{3/5}  % HyperKKL (r=4) beats latent AND concat, CI
\providecommand{\ConeHyperKKLBeatsLatConcShape}{}\renewcommand{\ConeHyperKKLBeatsLatConcShape}{4/5}  % HyperKKL (r=4) beats latent AND concat, CI
\providecommand{\ConeHyperKKLBeatsLatentId}{}\renewcommand{\ConeHyperKKLBeatsLatentId}{5/5}
\providecommand{\ConeHyperKKLBeatsLatentSeq}{}\renewcommand{\ConeHyperKKLBeatsLatentSeq}{3/5}
\providecommand{\ConeHyperKKLBeatsLatentShape}{}\renewcommand{\ConeHyperKKLBeatsLatentShape}{4/5}
\providecommand{\ConeHyperKKLLosesLatentId}{}\renewcommand{\ConeHyperKKLLosesLatentId}{0/5}
\providecommand{\ConeHyperKKLLosesLatentSeq}{}\renewcommand{\ConeHyperKKLLosesLatentSeq}{1/5}
\providecommand{\ConeHyperKKLLosesLatentShape}{}\renewcommand{\ConeHyperKKLLosesLatentShape}{0/5}
\providecommand{\ConeHyperKKLOverEkfMedianSeq}{}\renewcommand{\ConeHyperKKLOverEkfMedianSeq}{5.1}  % HyperKKL NRMSE / EKF NRMSE, exact model, median over systems
\providecommand{\ConeHyperKKLOverEkfRangeSeq}{}\renewcommand{\ConeHyperKKLOverEkfRangeSeq}{3.3--14.2}
\providecommand{\ConeHyperKKLOverEkfbatterySeq}{}\renewcommand{\ConeHyperKKLOverEkfbatterySeq}{14.2}
\providecommand{\ConeHyperKKLOverEkfbioreactorSeq}{}\renewcommand{\ConeHyperKKLOverEkfbioreactorSeq}{10.7}
\providecommand{\ConeHyperKKLOverEkflorenzSeq}{}\renewcommand{\ConeHyperKKLOverEkflorenzSeq}{4.0}
\providecommand{\ConeHyperKKLOverEkfquadTankSeq}{}\renewcommand{\ConeHyperKKLOverEkfquadTankSeq}{5.1}
\providecommand{\ConeHyperKKLOverEkfvdpSeq}{}\renewcommand{\ConeHyperKKLOverEkfvdpSeq}{3.3}
\providecommand{\ConeHyperKKLVsLatentRangeId}{}\renewcommand{\ConeHyperKKLVsLatentRangeId}{-10 to -68\%}
\providecommand{\ConeHyperKKLVsLatentRangeSeq}{}\renewcommand{\ConeHyperKKLVsLatentRangeSeq}{-30 to -58\%}
\providecommand{\ConeHyperKKLVsLatentRangeShape}{}\renewcommand{\ConeHyperKKLVsLatentRangeShape}{-32 to -69\%}
\providecommand{\ConeLatentBeatsAllId}{}\renewcommand{\ConeLatentBeatsAllId}{0/5}
\providecommand{\ConeLatentBeatsAllSeq}{}\renewcommand{\ConeLatentBeatsAllSeq}{1/5}
\providecommand{\ConeLatentBeatsAllShape}{}\renewcommand{\ConeLatentBeatsAllShape}{0/5}
\providecommand{\ConePreferredLocation}{}\renewcommand{\ConePreferredLocation}{none}
\providecommand{\CtestConcatBeatsAllId}{}\renewcommand{\CtestConcatBeatsAllId}{0/6}
\providecommand{\CtestConcatBeatsAllSeq}{}\renewcommand{\CtestConcatBeatsAllSeq}{0/6}
\providecommand{\CtestConcatBeatsAllShape}{}\renewcommand{\CtestConcatBeatsAllShape}{0/6}
\providecommand{\CtestConcatBeatsLatentId}{}\renewcommand{\CtestConcatBeatsLatentId}{2/6}
\providecommand{\CtestConcatBeatsLatentSeq}{}\renewcommand{\CtestConcatBeatsLatentSeq}{2/6}
\providecommand{\CtestConcatBeatsLatentShape}{}\renewcommand{\CtestConcatBeatsLatentShape}{2/6}
\providecommand{\CtestConcatLosesLatentId}{}\renewcommand{\CtestConcatLosesLatentId}{3/6}
\providecommand{\CtestConcatLosesLatentSeq}{}\renewcommand{\CtestConcatLosesLatentSeq}{3/6}
\providecommand{\CtestConcatLosesLatentShape}{}\renewcommand{\CtestConcatLosesLatentShape}{3/6}
\providecommand{\CtestConcatVsLatentRangeId}{}\renewcommand{\CtestConcatVsLatentRangeId}{-42 to -47\%}
\providecommand{\CtestConcatVsLatentRangeSeq}{}\renewcommand{\CtestConcatVsLatentRangeSeq}{-45 to -47\%}
\providecommand{\CtestConcatVsLatentRangeShape}{}\renewcommand{\CtestConcatVsLatentRangeShape}{-44 to -47\%}
\providecommand{\CtestDenseAblationBeatsHyperKKLId}{}\renewcommand{\CtestDenseAblationBeatsHyperKKLId}{1/6}
\providecommand{\CtestDenseAblationBeatsHyperKKLSeq}{}\renewcommand{\CtestDenseAblationBeatsHyperKKLSeq}{1/6}
\providecommand{\CtestDenseAblationBeatsHyperKKLShape}{}\renewcommand{\CtestDenseAblationBeatsHyperKKLShape}{1/6}
\providecommand{\CtestDenseAblationBeatsLatentId}{}\renewcommand{\CtestDenseAblationBeatsLatentId}{5/6}
\providecommand{\CtestDenseAblationBeatsLatentSeq}{}\renewcommand{\CtestDenseAblationBeatsLatentSeq}{4/6}
\providecommand{\CtestDenseAblationBeatsLatentShape}{}\renewcommand{\CtestDenseAblationBeatsLatentShape}{5/6}
\providecommand{\CtestDenseAblationLosesLatentId}{}\renewcommand{\CtestDenseAblationLosesLatentId}{0/6}
\providecommand{\CtestDenseAblationLosesLatentSeq}{}\renewcommand{\CtestDenseAblationLosesLatentSeq}{0/6}
\providecommand{\CtestDenseAblationLosesLatentShape}{}\renewcommand{\CtestDenseAblationLosesLatentShape}{0/6}
\providecommand{\CtestDenseAblationVsLatentRangeId}{}\renewcommand{\CtestDenseAblationVsLatentRangeId}{-37 to -92\%}
\providecommand{\CtestDenseAblationVsLatentRangeSeq}{}\renewcommand{\CtestDenseAblationVsLatentRangeSeq}{-46 to -89\%}
\providecommand{\CtestDenseAblationVsLatentRangeShape}{}\renewcommand{\CtestDenseAblationVsLatentRangeShape}{-35 to -91\%}
\providecommand{\CtestFilmBeatsAllId}{}\renewcommand{\CtestFilmBeatsAllId}{0/6}
\providecommand{\CtestFilmBeatsAllSeq}{}\renewcommand{\CtestFilmBeatsAllSeq}{0/6}
\providecommand{\CtestFilmBeatsAllShape}{}\renewcommand{\CtestFilmBeatsAllShape}{0/6}
\providecommand{\CtestFilmBeatsLatentId}{}\renewcommand{\CtestFilmBeatsLatentId}{5/6}
\providecommand{\CtestFilmBeatsLatentSeq}{}\renewcommand{\CtestFilmBeatsLatentSeq}{4/6}
\providecommand{\CtestFilmBeatsLatentShape}{}\renewcommand{\CtestFilmBeatsLatentShape}{5/6}
\providecommand{\CtestFilmLosesLatentId}{}\renewcommand{\CtestFilmLosesLatentId}{0/6}
\providecommand{\CtestFilmLosesLatentSeq}{}\renewcommand{\CtestFilmLosesLatentSeq}{0/6}
\providecommand{\CtestFilmLosesLatentShape}{}\renewcommand{\CtestFilmLosesLatentShape}{0/6}
\providecommand{\CtestFilmVsLatentRangeId}{}\renewcommand{\CtestFilmVsLatentRangeId}{-15 to -86\%}
\providecommand{\CtestFilmVsLatentRangeSeq}{}\renewcommand{\CtestFilmVsLatentRangeSeq}{-32 to -86\%}
\providecommand{\CtestFilmVsLatentRangeShape}{}\renewcommand{\CtestFilmVsLatentRangeShape}{-4 to -86\%}
\providecommand{\CtestHyperKKLBeatsAllId}{}\renewcommand{\CtestHyperKKLBeatsAllId}{5/6}
\providecommand{\CtestHyperKKLBeatsAllSeq}{}\renewcommand{\CtestHyperKKLBeatsAllSeq}{4/6}
\providecommand{\CtestHyperKKLBeatsAllShape}{}\renewcommand{\CtestHyperKKLBeatsAllShape}{4/6}
\providecommand{\CtestHyperKKLBeatsDenseAblationId}{}\renewcommand{\CtestHyperKKLBeatsDenseAblationId}{1/6}
\providecommand{\CtestHyperKKLBeatsDenseAblationSeq}{}\renewcommand{\CtestHyperKKLBeatsDenseAblationSeq}{1/6}
\providecommand{\CtestHyperKKLBeatsDenseAblationShape}{}\renewcommand{\CtestHyperKKLBeatsDenseAblationShape}{1/6}
\providecommand{\CtestHyperKKLBeatsLatConcId}{}\renewcommand{\CtestHyperKKLBeatsLatConcId}{5/6}  % HyperKKL (r=4) beats latent AND concat, CI
\providecommand{\CtestHyperKKLBeatsLatConcSeq}{}\renewcommand{\CtestHyperKKLBeatsLatConcSeq}{5/6}  % HyperKKL (r=4) beats latent AND concat, CI
\providecommand{\CtestHyperKKLBeatsLatConcShape}{}\renewcommand{\CtestHyperKKLBeatsLatConcShape}{5/6}  % HyperKKL (r=4) beats latent AND concat, CI
\providecommand{\CtestHyperKKLBeatsLatentId}{}\renewcommand{\CtestHyperKKLBeatsLatentId}{5/6}
\providecommand{\CtestHyperKKLBeatsLatentSeq}{}\renewcommand{\CtestHyperKKLBeatsLatentSeq}{5/6}
\providecommand{\CtestHyperKKLBeatsLatentShape}{}\renewcommand{\CtestHyperKKLBeatsLatentShape}{5/6}
\providecommand{\CtestHyperKKLLosesLatentId}{}\renewcommand{\CtestHyperKKLLosesLatentId}{0/6}
\providecommand{\CtestHyperKKLLosesLatentSeq}{}\renewcommand{\CtestHyperKKLLosesLatentSeq}{0/6}
\providecommand{\CtestHyperKKLLosesLatentShape}{}\renewcommand{\CtestHyperKKLLosesLatentShape}{0/6}
\providecommand{\CtestHyperKKLOverEkfMedianSeq}{}\renewcommand{\CtestHyperKKLOverEkfMedianSeq}{6.8}  % HyperKKL NRMSE / EKF NRMSE, exact model, median over systems
\providecommand{\CtestHyperKKLOverEkfRangeSeq}{}\renewcommand{\CtestHyperKKLOverEkfRangeSeq}{2.2--31.7}
\providecommand{\CtestHyperKKLOverEkfcartpoleSeq}{}\renewcommand{\CtestHyperKKLOverEkfcartpoleSeq}{31.7}
\providecommand{\CtestHyperKKLOverEkffhnSeq}{}\renewcommand{\CtestHyperKKLOverEkffhnSeq}{7.1}
\providecommand{\CtestHyperKKLOverEkfhighwayTrafficSeq}{}\renewcommand{\CtestHyperKKLOverEkfhighwayTrafficSeq}{6.5}
\providecommand{\CtestHyperKKLOverEkfnetworkSeq}{}\renewcommand{\CtestHyperKKLOverEkfnetworkSeq}{7.9}
\providecommand{\CtestHyperKKLOverEkfrosslerSeq}{}\renewcommand{\CtestHyperKKLOverEkfrosslerSeq}{2.2}
\providecommand{\CtestHyperKKLOverEkftwoLinkArmSeq}{}\renewcommand{\CtestHyperKKLOverEkftwoLinkArmSeq}{2.2}
\providecommand{\CtestHyperKKLVsLatentRangeId}{}\renewcommand{\CtestHyperKKLVsLatentRangeId}{-44 to -90\%}
\providecommand{\CtestHyperKKLVsLatentRangeSeq}{}\renewcommand{\CtestHyperKKLVsLatentRangeSeq}{-29 to -88\%}
\providecommand{\CtestHyperKKLVsLatentRangeShape}{}\renewcommand{\CtestHyperKKLVsLatentRangeShape}{-43 to -89\%}
\providecommand{\CtestLatentBeatsAllId}{}\renewcommand{\CtestLatentBeatsAllId}{0/6}
\providecommand{\CtestLatentBeatsAllSeq}{}\renewcommand{\CtestLatentBeatsAllSeq}{0/6}
\providecommand{\CtestLatentBeatsAllShape}{}\renewcommand{\CtestLatentBeatsAllShape}{0/6}
\providecommand{\CtestPreferredLocation}{}\renewcommand{\CtestPreferredLocation}{HyperKKL}
\providecommand{\CthreeFaultsBurstOutageRegime}{}\renewcommand{\CthreeFaultsBurstOutageRegime}{FAIL}
\providecommand{\CthreeFaultsCleanRegime}{}\renewcommand{\CthreeFaultsCleanRegime}{FAIL}
\providecommand{\CthreeFaultsColoredRegime}{}\renewcommand{\CthreeFaultsColoredRegime}{FAIL}
\providecommand{\CthreeFaultsColoredTRegime}{}\renewcommand{\CthreeFaultsColoredTRegime}{FAIL}
\providecommand{\CthreeFaultsDropoutRegime}{}\renewcommand{\CthreeFaultsDropoutRegime}{FAIL}
\providecommand{\CthreeFaultsInputBiasRegime}{}\renewcommand{\CthreeFaultsInputBiasRegime}{FAIL}
\providecommand{\CthreeFaultsInputDelayRegime}{}\renewcommand{\CthreeFaultsInputDelayRegime}{FAIL}
\providecommand{\CthreeFaultsInputNoiseRegime}{}\renewcommand{\CthreeFaultsInputNoiseRegime}{FAIL}
\providecommand{\CthreeFaultsNoiseTenRegime}{}\renewcommand{\CthreeFaultsNoiseTenRegime}{FAIL}
\providecommand{\CthreeFaultsOutliersRegime}{}\renewcommand{\CthreeFaultsOutliersRegime}{FAIL}
\providecommand{\CthreeFaultsPlantDriftTenRegime}{}\renewcommand{\CthreeFaultsPlantDriftTenRegime}{FAIL}
\providecommand{\CthreeFaultsPlantDriftThirtyRegime}{}\renewcommand{\CthreeFaultsPlantDriftThirtyRegime}{FAIL}
\providecommand{\CthreeFaultsPoorInitRegime}{}\renewcommand{\CthreeFaultsPoorInitRegime}{FAIL}
\providecommand{\CthreeFaultsQuantizedRegime}{}\renewcommand{\CthreeFaultsQuantizedRegime}{FAIL}
\providecommand{\CthreeFaultsRegimes}{}\renewcommand{\CthreeFaultsRegimes}{none of 15 conditions}
\providecommand{\CthreeFaultsSensorDriftRegime}{}\renewcommand{\CthreeFaultsSensorDriftRegime}{FAIL}
\providecommand{\CthreeMaxModelErrNoWin}{}\renewcommand{\CthreeMaxModelErrNoWin}{0.152}  % partial/quad_tank
\providecommand{\CthreeMinModelErrWin}{}\renewcommand{\CthreeMinModelErrWin}{0.064}  % mismatch20/battery
\providecommand{\CthreeMismatchFiftyIdBatteryBestDataDriven}{}\renewcommand{\CthreeMismatchFiftyIdBatteryBestDataDriven}{0.0185}
\providecommand{\CthreeMismatchFiftyIdBatteryBestDataDrivenMethod}{}\renewcommand{\CthreeMismatchFiftyIdBatteryBestDataDrivenMethod}{greybox\_ukf}
\providecommand{\CthreeMismatchFiftyIdBatteryBestLearned}{}\renewcommand{\CthreeMismatchFiftyIdBatteryBestLearned}{0.0381}
\providecommand{\CthreeMismatchFiftyIdBatteryBestWorking}{}\renewcommand{\CthreeMismatchFiftyIdBatteryBestWorking}{0.307}
\providecommand{\CthreeMismatchFiftyIdBatteryClimatology}{}\renewcommand{\CthreeMismatchFiftyIdBatteryClimatology}{0.736}
\providecommand{\CthreeMismatchFiftyIdBatteryFailedFilters}{}\renewcommand{\CthreeMismatchFiftyIdBatteryFailedFilters}{1/8}
\providecommand{\CthreeMismatchFiftyIdBatteryVsDataDriven}{}\renewcommand{\CthreeMismatchFiftyIdBatteryVsDataDriven}{+105\%}
\providecommand{\CthreeMismatchFiftyIdBatteryVsWorking}{}\renewcommand{\CthreeMismatchFiftyIdBatteryVsWorking}{-88\%}
\providecommand{\CthreeMismatchFiftyIdBioreactorBestDataDriven}{}\renewcommand{\CthreeMismatchFiftyIdBioreactorBestDataDriven}{0.0136}
\providecommand{\CthreeMismatchFiftyIdBioreactorBestDataDrivenMethod}{}\renewcommand{\CthreeMismatchFiftyIdBioreactorBestDataDrivenMethod}{nlox}
\providecommand{\CthreeMismatchFiftyIdBioreactorBestLearned}{}\renewcommand{\CthreeMismatchFiftyIdBioreactorBestLearned}{0.0501}
\providecommand{\CthreeMismatchFiftyIdBioreactorClimatology}{}\renewcommand{\CthreeMismatchFiftyIdBioreactorClimatology}{0.561}
\providecommand{\CthreeMismatchFiftyIdBioreactorFailedFilters}{}\renewcommand{\CthreeMismatchFiftyIdBioreactorFailedFilters}{8/8}
\providecommand{\CthreeMismatchFiftyIdBioreactorVsDataDriven}{}\renewcommand{\CthreeMismatchFiftyIdBioreactorVsDataDriven}{+285\%}
\providecommand{\CthreeMismatchFiftyIdLorenzBestDataDriven}{}\renewcommand{\CthreeMismatchFiftyIdLorenzBestDataDriven}{0.00386}
\providecommand{\CthreeMismatchFiftyIdLorenzBestDataDrivenMethod}{}\renewcommand{\CthreeMismatchFiftyIdLorenzBestDataDrivenMethod}{nlox}
\providecommand{\CthreeMismatchFiftyIdLorenzBestLearned}{}\renewcommand{\CthreeMismatchFiftyIdLorenzBestLearned}{0.00968}
\providecommand{\CthreeMismatchFiftyIdLorenzBestWorking}{}\renewcommand{\CthreeMismatchFiftyIdLorenzBestWorking}{0.0647}
\providecommand{\CthreeMismatchFiftyIdLorenzClimatology}{}\renewcommand{\CthreeMismatchFiftyIdLorenzClimatology}{0.409}
\providecommand{\CthreeMismatchFiftyIdLorenzFailedFilters}{}\renewcommand{\CthreeMismatchFiftyIdLorenzFailedFilters}{6/8}
\providecommand{\CthreeMismatchFiftyIdLorenzVsDataDriven}{}\renewcommand{\CthreeMismatchFiftyIdLorenzVsDataDriven}{+152\%}
\providecommand{\CthreeMismatchFiftyIdLorenzVsWorking}{}\renewcommand{\CthreeMismatchFiftyIdLorenzVsWorking}{-85\%}
\providecommand{\CthreeMismatchFiftyIdQuadTankBestDataDriven}{}\renewcommand{\CthreeMismatchFiftyIdQuadTankBestDataDriven}{0.0133}
\providecommand{\CthreeMismatchFiftyIdQuadTankBestDataDrivenMethod}{}\renewcommand{\CthreeMismatchFiftyIdQuadTankBestDataDrivenMethod}{greybox\_ekf}
\providecommand{\CthreeMismatchFiftyIdQuadTankBestLearned}{}\renewcommand{\CthreeMismatchFiftyIdQuadTankBestLearned}{0.0246}
\providecommand{\CthreeMismatchFiftyIdQuadTankClimatology}{}\renewcommand{\CthreeMismatchFiftyIdQuadTankClimatology}{0.549}
\providecommand{\CthreeMismatchFiftyIdQuadTankFailedFilters}{}\renewcommand{\CthreeMismatchFiftyIdQuadTankFailedFilters}{8/8}
\providecommand{\CthreeMismatchFiftyIdQuadTankVsDataDriven}{}\renewcommand{\CthreeMismatchFiftyIdQuadTankVsDataDriven}{+84\%}
\providecommand{\CthreeMismatchFiftyIdVdpBestDataDriven}{}\renewcommand{\CthreeMismatchFiftyIdVdpBestDataDriven}{0.0132}
\providecommand{\CthreeMismatchFiftyIdVdpBestDataDrivenMethod}{}\renewcommand{\CthreeMismatchFiftyIdVdpBestDataDrivenMethod}{greybox\_ekf}
\providecommand{\CthreeMismatchFiftyIdVdpBestLearned}{}\renewcommand{\CthreeMismatchFiftyIdVdpBestLearned}{0.0185}
\providecommand{\CthreeMismatchFiftyIdVdpBestWorking}{}\renewcommand{\CthreeMismatchFiftyIdVdpBestWorking}{0.0637}
\providecommand{\CthreeMismatchFiftyIdVdpClimatology}{}\renewcommand{\CthreeMismatchFiftyIdVdpClimatology}{1.2}
\providecommand{\CthreeMismatchFiftyIdVdpFailedFilters}{}\renewcommand{\CthreeMismatchFiftyIdVdpFailedFilters}{0/8}
\providecommand{\CthreeMismatchFiftyIdVdpVsDataDriven}{}\renewcommand{\CthreeMismatchFiftyIdVdpVsDataDriven}{+40\%}
\providecommand{\CthreeMismatchFiftyIdVdpVsWorking}{}\renewcommand{\CthreeMismatchFiftyIdVdpVsWorking}{-71\%}
\providecommand{\CthreeMismatchFiftyIdVsAny}{}\renewcommand{\CthreeMismatchFiftyIdVsAny}{5/5}
\providecommand{\CthreeMismatchFiftyIdVsDataDriven}{}\renewcommand{\CthreeMismatchFiftyIdVsDataDriven}{0/5}
\providecommand{\CthreeMismatchFiftyIdVsNominal}{}\renewcommand{\CthreeMismatchFiftyIdVsNominal}{5/5}
\providecommand{\CthreeMismatchFiftyIdVsRobust}{}\renewcommand{\CthreeMismatchFiftyIdVsRobust}{5/5}
\providecommand{\CthreeMismatchFiftyIdVsWorking}{}\renewcommand{\CthreeMismatchFiftyIdVsWorking}{5/5}
\providecommand{\CthreeMismatchFiftySeqBatteryBestDataDriven}{}\renewcommand{\CthreeMismatchFiftySeqBatteryBestDataDriven}{0.0222}
\providecommand{\CthreeMismatchFiftySeqBatteryBestDataDrivenMethod}{}\renewcommand{\CthreeMismatchFiftySeqBatteryBestDataDrivenMethod}{greybox\_ukf}
\providecommand{\CthreeMismatchFiftySeqBatteryBestLearned}{}\renewcommand{\CthreeMismatchFiftySeqBatteryBestLearned}{0.0503}
\providecommand{\CthreeMismatchFiftySeqBatteryBestWorking}{}\renewcommand{\CthreeMismatchFiftySeqBatteryBestWorking}{0.284}
\providecommand{\CthreeMismatchFiftySeqBatteryClimatology}{}\renewcommand{\CthreeMismatchFiftySeqBatteryClimatology}{0.756}
\providecommand{\CthreeMismatchFiftySeqBatteryFailedFilters}{}\renewcommand{\CthreeMismatchFiftySeqBatteryFailedFilters}{1/8}
\providecommand{\CthreeMismatchFiftySeqBatteryVsDataDriven}{}\renewcommand{\CthreeMismatchFiftySeqBatteryVsDataDriven}{+126\%}
\providecommand{\CthreeMismatchFiftySeqBatteryVsWorking}{}\renewcommand{\CthreeMismatchFiftySeqBatteryVsWorking}{-82\%}
\providecommand{\CthreeMismatchFiftySeqBioreactorBestDataDriven}{}\renewcommand{\CthreeMismatchFiftySeqBioreactorBestDataDriven}{0.0341}
\providecommand{\CthreeMismatchFiftySeqBioreactorBestDataDrivenMethod}{}\renewcommand{\CthreeMismatchFiftySeqBioreactorBestDataDrivenMethod}{nlox}
\providecommand{\CthreeMismatchFiftySeqBioreactorBestLearned}{}\renewcommand{\CthreeMismatchFiftySeqBioreactorBestLearned}{0.135}
\providecommand{\CthreeMismatchFiftySeqBioreactorBestWorking}{}\renewcommand{\CthreeMismatchFiftySeqBioreactorBestWorking}{0.353}
\providecommand{\CthreeMismatchFiftySeqBioreactorClimatology}{}\renewcommand{\CthreeMismatchFiftySeqBioreactorClimatology}{0.614}
\providecommand{\CthreeMismatchFiftySeqBioreactorFailedFilters}{}\renewcommand{\CthreeMismatchFiftySeqBioreactorFailedFilters}{6/8}
\providecommand{\CthreeMismatchFiftySeqBioreactorVsDataDriven}{}\renewcommand{\CthreeMismatchFiftySeqBioreactorVsDataDriven}{+327\%}
\providecommand{\CthreeMismatchFiftySeqBioreactorVsWorking}{}\renewcommand{\CthreeMismatchFiftySeqBioreactorVsWorking}{-62\%}
\providecommand{\CthreeMismatchFiftySeqLorenzBestDataDriven}{}\renewcommand{\CthreeMismatchFiftySeqLorenzBestDataDriven}{0.00447}
\providecommand{\CthreeMismatchFiftySeqLorenzBestDataDrivenMethod}{}\renewcommand{\CthreeMismatchFiftySeqLorenzBestDataDrivenMethod}{nlox}
\providecommand{\CthreeMismatchFiftySeqLorenzBestLearned}{}\renewcommand{\CthreeMismatchFiftySeqLorenzBestLearned}{0.0107}
\providecommand{\CthreeMismatchFiftySeqLorenzBestWorking}{}\renewcommand{\CthreeMismatchFiftySeqLorenzBestWorking}{0.0708}
\providecommand{\CthreeMismatchFiftySeqLorenzClimatology}{}\renewcommand{\CthreeMismatchFiftySeqLorenzClimatology}{0.814}
\providecommand{\CthreeMismatchFiftySeqLorenzFailedFilters}{}\renewcommand{\CthreeMismatchFiftySeqLorenzFailedFilters}{1/8}
\providecommand{\CthreeMismatchFiftySeqLorenzVsDataDriven}{}\renewcommand{\CthreeMismatchFiftySeqLorenzVsDataDriven}{+142\%}
\providecommand{\CthreeMismatchFiftySeqLorenzVsWorking}{}\renewcommand{\CthreeMismatchFiftySeqLorenzVsWorking}{-85\%}
\providecommand{\CthreeMismatchFiftySeqQuadTankBestDataDriven}{}\renewcommand{\CthreeMismatchFiftySeqQuadTankBestDataDriven}{0.0131}
\providecommand{\CthreeMismatchFiftySeqQuadTankBestDataDrivenMethod}{}\renewcommand{\CthreeMismatchFiftySeqQuadTankBestDataDrivenMethod}{greybox\_ekf}
\providecommand{\CthreeMismatchFiftySeqQuadTankBestLearned}{}\renewcommand{\CthreeMismatchFiftySeqQuadTankBestLearned}{0.0442}
\providecommand{\CthreeMismatchFiftySeqQuadTankClimatology}{}\renewcommand{\CthreeMismatchFiftySeqQuadTankClimatology}{0.6}
\providecommand{\CthreeMismatchFiftySeqQuadTankFailedFilters}{}\renewcommand{\CthreeMismatchFiftySeqQuadTankFailedFilters}{8/8}
\providecommand{\CthreeMismatchFiftySeqQuadTankVsDataDriven}{}\renewcommand{\CthreeMismatchFiftySeqQuadTankVsDataDriven}{+238\%}
\providecommand{\CthreeMismatchFiftySeqVdpBestDataDriven}{}\renewcommand{\CthreeMismatchFiftySeqVdpBestDataDriven}{0.015}
\providecommand{\CthreeMismatchFiftySeqVdpBestDataDrivenMethod}{}\renewcommand{\CthreeMismatchFiftySeqVdpBestDataDrivenMethod}{greybox\_ekf}
\providecommand{\CthreeMismatchFiftySeqVdpBestLearned}{}\renewcommand{\CthreeMismatchFiftySeqVdpBestLearned}{0.0203}
\providecommand{\CthreeMismatchFiftySeqVdpBestWorking}{}\renewcommand{\CthreeMismatchFiftySeqVdpBestWorking}{0.0777}
\providecommand{\CthreeMismatchFiftySeqVdpClimatology}{}\renewcommand{\CthreeMismatchFiftySeqVdpClimatology}{1.23}
\providecommand{\CthreeMismatchFiftySeqVdpFailedFilters}{}\renewcommand{\CthreeMismatchFiftySeqVdpFailedFilters}{0/8}
\providecommand{\CthreeMismatchFiftySeqVdpVsDataDriven}{}\renewcommand{\CthreeMismatchFiftySeqVdpVsDataDriven}{+36\%}
\providecommand{\CthreeMismatchFiftySeqVdpVsWorking}{}\renewcommand{\CthreeMismatchFiftySeqVdpVsWorking}{-74\%}
\providecommand{\CthreeMismatchFiftySeqVsAny}{}\renewcommand{\CthreeMismatchFiftySeqVsAny}{5/5}
\providecommand{\CthreeMismatchFiftySeqVsDataDriven}{}\renewcommand{\CthreeMismatchFiftySeqVsDataDriven}{0/5}
\providecommand{\CthreeMismatchFiftySeqVsNominal}{}\renewcommand{\CthreeMismatchFiftySeqVsNominal}{5/5}
\providecommand{\CthreeMismatchFiftySeqVsRobust}{}\renewcommand{\CthreeMismatchFiftySeqVsRobust}{5/5}
\providecommand{\CthreeMismatchFiftySeqVsWorking}{}\renewcommand{\CthreeMismatchFiftySeqVsWorking}{5/5}
\providecommand{\CthreeMismatchTenIdBatteryBestDataDriven}{}\renewcommand{\CthreeMismatchTenIdBatteryBestDataDriven}{0.007}
\providecommand{\CthreeMismatchTenIdBatteryBestDataDrivenMethod}{}\renewcommand{\CthreeMismatchTenIdBatteryBestDataDrivenMethod}{greybox\_ekf}
\providecommand{\CthreeMismatchTenIdBatteryBestLearned}{}\renewcommand{\CthreeMismatchTenIdBatteryBestLearned}{0.0448}
\providecommand{\CthreeMismatchTenIdBatteryBestWorking}{}\renewcommand{\CthreeMismatchTenIdBatteryBestWorking}{0.0574}
\providecommand{\CthreeMismatchTenIdBatteryClimatology}{}\renewcommand{\CthreeMismatchTenIdBatteryClimatology}{0.743}
\providecommand{\CthreeMismatchTenIdBatteryFailedFilters}{}\renewcommand{\CthreeMismatchTenIdBatteryFailedFilters}{1/8}
\providecommand{\CthreeMismatchTenIdBatteryVsDataDriven}{}\renewcommand{\CthreeMismatchTenIdBatteryVsDataDriven}{+540\%}
\providecommand{\CthreeMismatchTenIdBatteryVsWorking}{}\renewcommand{\CthreeMismatchTenIdBatteryVsWorking}{-22\%}
\providecommand{\CthreeMismatchTenIdBioreactorBestDataDriven}{}\renewcommand{\CthreeMismatchTenIdBioreactorBestDataDriven}{0.0159}
\providecommand{\CthreeMismatchTenIdBioreactorBestDataDrivenMethod}{}\renewcommand{\CthreeMismatchTenIdBioreactorBestDataDrivenMethod}{nlox}
\providecommand{\CthreeMismatchTenIdBioreactorBestLearned}{}\renewcommand{\CthreeMismatchTenIdBioreactorBestLearned}{0.02}
\providecommand{\CthreeMismatchTenIdBioreactorBestWorking}{}\renewcommand{\CthreeMismatchTenIdBioreactorBestWorking}{0.372}
\providecommand{\CthreeMismatchTenIdBioreactorClimatology}{}\renewcommand{\CthreeMismatchTenIdBioreactorClimatology}{0.794}
\providecommand{\CthreeMismatchTenIdBioreactorFailedFilters}{}\renewcommand{\CthreeMismatchTenIdBioreactorFailedFilters}{5/8}
\providecommand{\CthreeMismatchTenIdBioreactorVsDataDriven}{}\renewcommand{\CthreeMismatchTenIdBioreactorVsDataDriven}{+28\%}
\providecommand{\CthreeMismatchTenIdBioreactorVsWorking}{}\renewcommand{\CthreeMismatchTenIdBioreactorVsWorking}{-95\%}
\providecommand{\CthreeMismatchTenIdLorenzBestDataDriven}{}\renewcommand{\CthreeMismatchTenIdLorenzBestDataDriven}{0.0202}
\providecommand{\CthreeMismatchTenIdLorenzBestDataDrivenMethod}{}\renewcommand{\CthreeMismatchTenIdLorenzBestDataDrivenMethod}{greybox\_ekf\_oracle}
\providecommand{\CthreeMismatchTenIdLorenzBestLearned}{}\renewcommand{\CthreeMismatchTenIdLorenzBestLearned}{0.0215}
\providecommand{\CthreeMismatchTenIdLorenzBestWorking}{}\renewcommand{\CthreeMismatchTenIdLorenzBestWorking}{0.158}
\providecommand{\CthreeMismatchTenIdLorenzClimatology}{}\renewcommand{\CthreeMismatchTenIdLorenzClimatology}{0.912}
\providecommand{\CthreeMismatchTenIdLorenzFailedFilters}{}\renewcommand{\CthreeMismatchTenIdLorenzFailedFilters}{0/8}
\providecommand{\CthreeMismatchTenIdLorenzVsDataDriven}{}\renewcommand{\CthreeMismatchTenIdLorenzVsDataDriven}{+7\%}
\providecommand{\CthreeMismatchTenIdLorenzVsWorking}{}\renewcommand{\CthreeMismatchTenIdLorenzVsWorking}{-86\%}
\providecommand{\CthreeMismatchTenIdQuadTankBestDataDriven}{}\renewcommand{\CthreeMismatchTenIdQuadTankBestDataDriven}{0.014}
\providecommand{\CthreeMismatchTenIdQuadTankBestDataDrivenMethod}{}\renewcommand{\CthreeMismatchTenIdQuadTankBestDataDrivenMethod}{greybox\_ekf}
\providecommand{\CthreeMismatchTenIdQuadTankBestLearned}{}\renewcommand{\CthreeMismatchTenIdQuadTankBestLearned}{0.0367}
\providecommand{\CthreeMismatchTenIdQuadTankBestWorking}{}\renewcommand{\CthreeMismatchTenIdQuadTankBestWorking}{0.429}
\providecommand{\CthreeMismatchTenIdQuadTankClimatology}{}\renewcommand{\CthreeMismatchTenIdQuadTankClimatology}{0.554}
\providecommand{\CthreeMismatchTenIdQuadTankFailedFilters}{}\renewcommand{\CthreeMismatchTenIdQuadTankFailedFilters}{7/8}
\providecommand{\CthreeMismatchTenIdQuadTankVsDataDriven}{}\renewcommand{\CthreeMismatchTenIdQuadTankVsDataDriven}{+162\%}
\providecommand{\CthreeMismatchTenIdQuadTankVsWorking}{}\renewcommand{\CthreeMismatchTenIdQuadTankVsWorking}{-91\%}
\providecommand{\CthreeMismatchTenIdVdpBestDataDriven}{}\renewcommand{\CthreeMismatchTenIdVdpBestDataDriven}{0.00601}
\providecommand{\CthreeMismatchTenIdVdpBestDataDrivenMethod}{}\renewcommand{\CthreeMismatchTenIdVdpBestDataDrivenMethod}{greybox\_ukf}
\providecommand{\CthreeMismatchTenIdVdpBestLearned}{}\renewcommand{\CthreeMismatchTenIdVdpBestLearned}{0.0172}
\providecommand{\CthreeMismatchTenIdVdpBestWorking}{}\renewcommand{\CthreeMismatchTenIdVdpBestWorking}{0.0244}
\providecommand{\CthreeMismatchTenIdVdpClimatology}{}\renewcommand{\CthreeMismatchTenIdVdpClimatology}{1.06}
\providecommand{\CthreeMismatchTenIdVdpFailedFilters}{}\renewcommand{\CthreeMismatchTenIdVdpFailedFilters}{0/8}
\providecommand{\CthreeMismatchTenIdVdpVsDataDriven}{}\renewcommand{\CthreeMismatchTenIdVdpVsDataDriven}{+187\%}
\providecommand{\CthreeMismatchTenIdVdpVsWorking}{}\renewcommand{\CthreeMismatchTenIdVdpVsWorking}{-29\%}
\providecommand{\CthreeMismatchTenIdVsAny}{}\renewcommand{\CthreeMismatchTenIdVsAny}{5/5}
\providecommand{\CthreeMismatchTenIdVsDataDriven}{}\renewcommand{\CthreeMismatchTenIdVsDataDriven}{0/5}
\providecommand{\CthreeMismatchTenIdVsNominal}{}\renewcommand{\CthreeMismatchTenIdVsNominal}{5/5}
\providecommand{\CthreeMismatchTenIdVsRobust}{}\renewcommand{\CthreeMismatchTenIdVsRobust}{5/5}
\providecommand{\CthreeMismatchTenIdVsWorking}{}\renewcommand{\CthreeMismatchTenIdVsWorking}{5/5}
\providecommand{\CthreeMismatchTenSeqBatteryBestDataDriven}{}\renewcommand{\CthreeMismatchTenSeqBatteryBestDataDriven}{0.00715}
\providecommand{\CthreeMismatchTenSeqBatteryBestDataDrivenMethod}{}\renewcommand{\CthreeMismatchTenSeqBatteryBestDataDrivenMethod}{greybox\_ekf}
\providecommand{\CthreeMismatchTenSeqBatteryBestLearned}{}\renewcommand{\CthreeMismatchTenSeqBatteryBestLearned}{0.0646}
\providecommand{\CthreeMismatchTenSeqBatteryBestWorking}{}\renewcommand{\CthreeMismatchTenSeqBatteryBestWorking}{0.0605}
\providecommand{\CthreeMismatchTenSeqBatteryClimatology}{}\renewcommand{\CthreeMismatchTenSeqBatteryClimatology}{0.776}
\providecommand{\CthreeMismatchTenSeqBatteryFailedFilters}{}\renewcommand{\CthreeMismatchTenSeqBatteryFailedFilters}{1/8}
\providecommand{\CthreeMismatchTenSeqBatteryVsDataDriven}{}\renewcommand{\CthreeMismatchTenSeqBatteryVsDataDriven}{+800\%}
\providecommand{\CthreeMismatchTenSeqBatteryVsWorking}{}\renewcommand{\CthreeMismatchTenSeqBatteryVsWorking}{+7\%}
\providecommand{\CthreeMismatchTenSeqBioreactorBestDataDriven}{}\renewcommand{\CthreeMismatchTenSeqBioreactorBestDataDriven}{0.0214}
\providecommand{\CthreeMismatchTenSeqBioreactorBestDataDrivenMethod}{}\renewcommand{\CthreeMismatchTenSeqBioreactorBestDataDrivenMethod}{greybox\_ekf\_oracle}
\providecommand{\CthreeMismatchTenSeqBioreactorBestLearned}{}\renewcommand{\CthreeMismatchTenSeqBioreactorBestLearned}{0.0613}
\providecommand{\CthreeMismatchTenSeqBioreactorBestWorking}{}\renewcommand{\CthreeMismatchTenSeqBioreactorBestWorking}{0.356}
\providecommand{\CthreeMismatchTenSeqBioreactorClimatology}{}\renewcommand{\CthreeMismatchTenSeqBioreactorClimatology}{0.841}
\providecommand{\CthreeMismatchTenSeqBioreactorFailedFilters}{}\renewcommand{\CthreeMismatchTenSeqBioreactorFailedFilters}{5/8}
\providecommand{\CthreeMismatchTenSeqBioreactorVsDataDriven}{}\renewcommand{\CthreeMismatchTenSeqBioreactorVsDataDriven}{+187\%}
\providecommand{\CthreeMismatchTenSeqBioreactorVsWorking}{}\renewcommand{\CthreeMismatchTenSeqBioreactorVsWorking}{-83\%}
\providecommand{\CthreeMismatchTenSeqLorenzBestDataDriven}{}\renewcommand{\CthreeMismatchTenSeqLorenzBestDataDriven}{0.0201}
\providecommand{\CthreeMismatchTenSeqLorenzBestDataDrivenMethod}{}\renewcommand{\CthreeMismatchTenSeqLorenzBestDataDrivenMethod}{greybox\_ekf\_oracle}
\providecommand{\CthreeMismatchTenSeqLorenzBestLearned}{}\renewcommand{\CthreeMismatchTenSeqLorenzBestLearned}{0.0219}
\providecommand{\CthreeMismatchTenSeqLorenzBestWorking}{}\renewcommand{\CthreeMismatchTenSeqLorenzBestWorking}{0.157}
\providecommand{\CthreeMismatchTenSeqLorenzClimatology}{}\renewcommand{\CthreeMismatchTenSeqLorenzClimatology}{0.911}
\providecommand{\CthreeMismatchTenSeqLorenzFailedFilters}{}\renewcommand{\CthreeMismatchTenSeqLorenzFailedFilters}{0/8}
\providecommand{\CthreeMismatchTenSeqLorenzVsDataDriven}{}\renewcommand{\CthreeMismatchTenSeqLorenzVsDataDriven}{+9\%}
\providecommand{\CthreeMismatchTenSeqLorenzVsWorking}{}\renewcommand{\CthreeMismatchTenSeqLorenzVsWorking}{-86\%}
\providecommand{\CthreeMismatchTenSeqQuadTankBestDataDriven}{}\renewcommand{\CthreeMismatchTenSeqQuadTankBestDataDriven}{0.0135}
\providecommand{\CthreeMismatchTenSeqQuadTankBestDataDrivenMethod}{}\renewcommand{\CthreeMismatchTenSeqQuadTankBestDataDrivenMethod}{greybox\_ekf}
\providecommand{\CthreeMismatchTenSeqQuadTankBestLearned}{}\renewcommand{\CthreeMismatchTenSeqQuadTankBestLearned}{0.0513}
\providecommand{\CthreeMismatchTenSeqQuadTankBestWorking}{}\renewcommand{\CthreeMismatchTenSeqQuadTankBestWorking}{0.434}
\providecommand{\CthreeMismatchTenSeqQuadTankClimatology}{}\renewcommand{\CthreeMismatchTenSeqQuadTankClimatology}{0.607}
\providecommand{\CthreeMismatchTenSeqQuadTankFailedFilters}{}\renewcommand{\CthreeMismatchTenSeqQuadTankFailedFilters}{7/8}
\providecommand{\CthreeMismatchTenSeqQuadTankVsDataDriven}{}\renewcommand{\CthreeMismatchTenSeqQuadTankVsDataDriven}{+279\%}
\providecommand{\CthreeMismatchTenSeqQuadTankVsWorking}{}\renewcommand{\CthreeMismatchTenSeqQuadTankVsWorking}{-88\%}
\providecommand{\CthreeMismatchTenSeqVdpBestDataDriven}{}\renewcommand{\CthreeMismatchTenSeqVdpBestDataDriven}{0.00634}
\providecommand{\CthreeMismatchTenSeqVdpBestDataDrivenMethod}{}\renewcommand{\CthreeMismatchTenSeqVdpBestDataDrivenMethod}{greybox\_ekf}
\providecommand{\CthreeMismatchTenSeqVdpBestLearned}{}\renewcommand{\CthreeMismatchTenSeqVdpBestLearned}{0.02}
\providecommand{\CthreeMismatchTenSeqVdpBestWorking}{}\renewcommand{\CthreeMismatchTenSeqVdpBestWorking}{0.0268}
\providecommand{\CthreeMismatchTenSeqVdpClimatology}{}\renewcommand{\CthreeMismatchTenSeqVdpClimatology}{1.11}
\providecommand{\CthreeMismatchTenSeqVdpFailedFilters}{}\renewcommand{\CthreeMismatchTenSeqVdpFailedFilters}{0/8}
\providecommand{\CthreeMismatchTenSeqVdpVsDataDriven}{}\renewcommand{\CthreeMismatchTenSeqVdpVsDataDriven}{+216\%}
\providecommand{\CthreeMismatchTenSeqVdpVsWorking}{}\renewcommand{\CthreeMismatchTenSeqVdpVsWorking}{-25\%}
\providecommand{\CthreeMismatchTenSeqVsAny}{}\renewcommand{\CthreeMismatchTenSeqVsAny}{4/5}
\providecommand{\CthreeMismatchTenSeqVsDataDriven}{}\renewcommand{\CthreeMismatchTenSeqVsDataDriven}{0/5}
\providecommand{\CthreeMismatchTenSeqVsNominal}{}\renewcommand{\CthreeMismatchTenSeqVsNominal}{4/5}
\providecommand{\CthreeMismatchTenSeqVsRobust}{}\renewcommand{\CthreeMismatchTenSeqVsRobust}{4/5}
\providecommand{\CthreeMismatchTenSeqVsWorking}{}\renewcommand{\CthreeMismatchTenSeqVsWorking}{4/5}
\providecommand{\CthreeMismatchThirtyIdBatteryBestDataDriven}{}\renewcommand{\CthreeMismatchThirtyIdBatteryBestDataDriven}{0.0127}
\providecommand{\CthreeMismatchThirtyIdBatteryBestDataDrivenMethod}{}\renewcommand{\CthreeMismatchThirtyIdBatteryBestDataDrivenMethod}{greybox\_ekf}
\providecommand{\CthreeMismatchThirtyIdBatteryBestLearned}{}\renewcommand{\CthreeMismatchThirtyIdBatteryBestLearned}{0.0382}
\providecommand{\CthreeMismatchThirtyIdBatteryBestWorking}{}\renewcommand{\CthreeMismatchThirtyIdBatteryBestWorking}{0.18}
\providecommand{\CthreeMismatchThirtyIdBatteryClimatology}{}\renewcommand{\CthreeMismatchThirtyIdBatteryClimatology}{0.735}
\providecommand{\CthreeMismatchThirtyIdBatteryFailedFilters}{}\renewcommand{\CthreeMismatchThirtyIdBatteryFailedFilters}{1/8}
\providecommand{\CthreeMismatchThirtyIdBatteryVsDataDriven}{}\renewcommand{\CthreeMismatchThirtyIdBatteryVsDataDriven}{+200\%}
\providecommand{\CthreeMismatchThirtyIdBatteryVsWorking}{}\renewcommand{\CthreeMismatchThirtyIdBatteryVsWorking}{-79\%}
\providecommand{\CthreeMismatchThirtyIdBioreactorBestDataDriven}{}\renewcommand{\CthreeMismatchThirtyIdBioreactorBestDataDriven}{0.0143}
\providecommand{\CthreeMismatchThirtyIdBioreactorBestDataDrivenMethod}{}\renewcommand{\CthreeMismatchThirtyIdBioreactorBestDataDrivenMethod}{nlox}
\providecommand{\CthreeMismatchThirtyIdBioreactorBestLearned}{}\renewcommand{\CthreeMismatchThirtyIdBioreactorBestLearned}{0.0329}
\providecommand{\CthreeMismatchThirtyIdBioreactorBestWorking}{}\renewcommand{\CthreeMismatchThirtyIdBioreactorBestWorking}{0.505}
\providecommand{\CthreeMismatchThirtyIdBioreactorClimatology}{}\renewcommand{\CthreeMismatchThirtyIdBioreactorClimatology}{0.655}
\providecommand{\CthreeMismatchThirtyIdBioreactorFailedFilters}{}\renewcommand{\CthreeMismatchThirtyIdBioreactorFailedFilters}{6/8}
\providecommand{\CthreeMismatchThirtyIdBioreactorVsDataDriven}{}\renewcommand{\CthreeMismatchThirtyIdBioreactorVsDataDriven}{+140\%}
\providecommand{\CthreeMismatchThirtyIdBioreactorVsWorking}{}\renewcommand{\CthreeMismatchThirtyIdBioreactorVsWorking}{-93\%}
\providecommand{\CthreeMismatchThirtyIdLorenzBestDataDriven}{}\renewcommand{\CthreeMismatchThirtyIdLorenzBestDataDriven}{0.00669}
\providecommand{\CthreeMismatchThirtyIdLorenzBestDataDrivenMethod}{}\renewcommand{\CthreeMismatchThirtyIdLorenzBestDataDrivenMethod}{nlox}
\providecommand{\CthreeMismatchThirtyIdLorenzBestLearned}{}\renewcommand{\CthreeMismatchThirtyIdLorenzBestLearned}{0.00664}
\providecommand{\CthreeMismatchThirtyIdLorenzBestWorking}{}\renewcommand{\CthreeMismatchThirtyIdLorenzBestWorking}{0.0129}
\providecommand{\CthreeMismatchThirtyIdLorenzClimatology}{}\renewcommand{\CthreeMismatchThirtyIdLorenzClimatology}{0.614}
\providecommand{\CthreeMismatchThirtyIdLorenzFailedFilters}{}\renewcommand{\CthreeMismatchThirtyIdLorenzFailedFilters}{0/8}
\providecommand{\CthreeMismatchThirtyIdLorenzVsDataDriven}{}\renewcommand{\CthreeMismatchThirtyIdLorenzVsDataDriven}{-0\%}
\providecommand{\CthreeMismatchThirtyIdLorenzVsWorking}{}\renewcommand{\CthreeMismatchThirtyIdLorenzVsWorking}{-49\%}
\providecommand{\CthreeMismatchThirtyIdQuadTankBestDataDriven}{}\renewcommand{\CthreeMismatchThirtyIdQuadTankBestDataDriven}{0.0134}
\providecommand{\CthreeMismatchThirtyIdQuadTankBestDataDrivenMethod}{}\renewcommand{\CthreeMismatchThirtyIdQuadTankBestDataDrivenMethod}{greybox\_ekf}
\providecommand{\CthreeMismatchThirtyIdQuadTankBestLearned}{}\renewcommand{\CthreeMismatchThirtyIdQuadTankBestLearned}{0.0221}
\providecommand{\CthreeMismatchThirtyIdQuadTankClimatology}{}\renewcommand{\CthreeMismatchThirtyIdQuadTankClimatology}{0.53}
\providecommand{\CthreeMismatchThirtyIdQuadTankFailedFilters}{}\renewcommand{\CthreeMismatchThirtyIdQuadTankFailedFilters}{8/8}
\providecommand{\CthreeMismatchThirtyIdQuadTankVsDataDriven}{}\renewcommand{\CthreeMismatchThirtyIdQuadTankVsDataDriven}{+65\%}
\providecommand{\CthreeMismatchThirtyIdVdpBestDataDriven}{}\renewcommand{\CthreeMismatchThirtyIdVdpBestDataDriven}{0.00893}
\providecommand{\CthreeMismatchThirtyIdVdpBestDataDrivenMethod}{}\renewcommand{\CthreeMismatchThirtyIdVdpBestDataDrivenMethod}{greybox\_ekf}
\providecommand{\CthreeMismatchThirtyIdVdpBestLearned}{}\renewcommand{\CthreeMismatchThirtyIdVdpBestLearned}{0.0179}
\providecommand{\CthreeMismatchThirtyIdVdpBestWorking}{}\renewcommand{\CthreeMismatchThirtyIdVdpBestWorking}{0.0523}
\providecommand{\CthreeMismatchThirtyIdVdpClimatology}{}\renewcommand{\CthreeMismatchThirtyIdVdpClimatology}{1.13}
\providecommand{\CthreeMismatchThirtyIdVdpFailedFilters}{}\renewcommand{\CthreeMismatchThirtyIdVdpFailedFilters}{0/8}
\providecommand{\CthreeMismatchThirtyIdVdpVsDataDriven}{}\renewcommand{\CthreeMismatchThirtyIdVdpVsDataDriven}{+101\%}
\providecommand{\CthreeMismatchThirtyIdVdpVsWorking}{}\renewcommand{\CthreeMismatchThirtyIdVdpVsWorking}{-66\%}
\providecommand{\CthreeMismatchThirtyIdVsAny}{}\renewcommand{\CthreeMismatchThirtyIdVsAny}{5/5}
\providecommand{\CthreeMismatchThirtyIdVsDataDriven}{}\renewcommand{\CthreeMismatchThirtyIdVsDataDriven}{0/5}
\providecommand{\CthreeMismatchThirtyIdVsNominal}{}\renewcommand{\CthreeMismatchThirtyIdVsNominal}{5/5}
\providecommand{\CthreeMismatchThirtyIdVsRobust}{}\renewcommand{\CthreeMismatchThirtyIdVsRobust}{5/5}
\providecommand{\CthreeMismatchThirtyIdVsWorking}{}\renewcommand{\CthreeMismatchThirtyIdVsWorking}{5/5}
\providecommand{\CthreeMismatchThirtySeqBatteryBestDataDriven}{}\renewcommand{\CthreeMismatchThirtySeqBatteryBestDataDriven}{0.0151}
\providecommand{\CthreeMismatchThirtySeqBatteryBestDataDrivenMethod}{}\renewcommand{\CthreeMismatchThirtySeqBatteryBestDataDrivenMethod}{greybox\_ekf}
\providecommand{\CthreeMismatchThirtySeqBatteryBestLearned}{}\renewcommand{\CthreeMismatchThirtySeqBatteryBestLearned}{0.0515}
\providecommand{\CthreeMismatchThirtySeqBatteryBestWorking}{}\renewcommand{\CthreeMismatchThirtySeqBatteryBestWorking}{0.177}
\providecommand{\CthreeMismatchThirtySeqBatteryClimatology}{}\renewcommand{\CthreeMismatchThirtySeqBatteryClimatology}{0.765}
\providecommand{\CthreeMismatchThirtySeqBatteryFailedFilters}{}\renewcommand{\CthreeMismatchThirtySeqBatteryFailedFilters}{1/8}
\providecommand{\CthreeMismatchThirtySeqBatteryVsDataDriven}{}\renewcommand{\CthreeMismatchThirtySeqBatteryVsDataDriven}{+241\%}
\providecommand{\CthreeMismatchThirtySeqBatteryVsWorking}{}\renewcommand{\CthreeMismatchThirtySeqBatteryVsWorking}{-71\%}
\providecommand{\CthreeMismatchThirtySeqBioreactorBestDataDriven}{}\renewcommand{\CthreeMismatchThirtySeqBioreactorBestDataDriven}{0.0353}
\providecommand{\CthreeMismatchThirtySeqBioreactorBestDataDrivenMethod}{}\renewcommand{\CthreeMismatchThirtySeqBioreactorBestDataDrivenMethod}{nlox}
\providecommand{\CthreeMismatchThirtySeqBioreactorBestLearned}{}\renewcommand{\CthreeMismatchThirtySeqBioreactorBestLearned}{0.0977}
\providecommand{\CthreeMismatchThirtySeqBioreactorBestWorking}{}\renewcommand{\CthreeMismatchThirtySeqBioreactorBestWorking}{0.247}
\providecommand{\CthreeMismatchThirtySeqBioreactorClimatology}{}\renewcommand{\CthreeMismatchThirtySeqBioreactorClimatology}{0.708}
\providecommand{\CthreeMismatchThirtySeqBioreactorFailedFilters}{}\renewcommand{\CthreeMismatchThirtySeqBioreactorFailedFilters}{6/8}
\providecommand{\CthreeMismatchThirtySeqBioreactorVsDataDriven}{}\renewcommand{\CthreeMismatchThirtySeqBioreactorVsDataDriven}{+194\%}
\providecommand{\CthreeMismatchThirtySeqBioreactorVsWorking}{}\renewcommand{\CthreeMismatchThirtySeqBioreactorVsWorking}{-60\%}
\providecommand{\CthreeMismatchThirtySeqLorenzBestDataDriven}{}\renewcommand{\CthreeMismatchThirtySeqLorenzBestDataDriven}{0.00682}
\providecommand{\CthreeMismatchThirtySeqLorenzBestDataDrivenMethod}{}\renewcommand{\CthreeMismatchThirtySeqLorenzBestDataDrivenMethod}{nlox}
\providecommand{\CthreeMismatchThirtySeqLorenzBestLearned}{}\renewcommand{\CthreeMismatchThirtySeqLorenzBestLearned}{0.007}
\providecommand{\CthreeMismatchThirtySeqLorenzBestWorking}{}\renewcommand{\CthreeMismatchThirtySeqLorenzBestWorking}{0.0142}
\providecommand{\CthreeMismatchThirtySeqLorenzClimatology}{}\renewcommand{\CthreeMismatchThirtySeqLorenzClimatology}{0.614}
\providecommand{\CthreeMismatchThirtySeqLorenzFailedFilters}{}\renewcommand{\CthreeMismatchThirtySeqLorenzFailedFilters}{1/8}
\providecommand{\CthreeMismatchThirtySeqLorenzVsDataDriven}{}\renewcommand{\CthreeMismatchThirtySeqLorenzVsDataDriven}{+3\%}
\providecommand{\CthreeMismatchThirtySeqLorenzVsWorking}{}\renewcommand{\CthreeMismatchThirtySeqLorenzVsWorking}{-51\%}
\providecommand{\CthreeMismatchThirtySeqQuadTankBestDataDriven}{}\renewcommand{\CthreeMismatchThirtySeqQuadTankBestDataDriven}{0.0129}
\providecommand{\CthreeMismatchThirtySeqQuadTankBestDataDrivenMethod}{}\renewcommand{\CthreeMismatchThirtySeqQuadTankBestDataDrivenMethod}{greybox\_ekf}
\providecommand{\CthreeMismatchThirtySeqQuadTankBestLearned}{}\renewcommand{\CthreeMismatchThirtySeqQuadTankBestLearned}{0.0422}
\providecommand{\CthreeMismatchThirtySeqQuadTankClimatology}{}\renewcommand{\CthreeMismatchThirtySeqQuadTankClimatology}{0.581}
\providecommand{\CthreeMismatchThirtySeqQuadTankFailedFilters}{}\renewcommand{\CthreeMismatchThirtySeqQuadTankFailedFilters}{8/8}
\providecommand{\CthreeMismatchThirtySeqQuadTankVsDataDriven}{}\renewcommand{\CthreeMismatchThirtySeqQuadTankVsDataDriven}{+226\%}
\providecommand{\CthreeMismatchThirtySeqVdpBestDataDriven}{}\renewcommand{\CthreeMismatchThirtySeqVdpBestDataDriven}{0.00976}
\providecommand{\CthreeMismatchThirtySeqVdpBestDataDrivenMethod}{}\renewcommand{\CthreeMismatchThirtySeqVdpBestDataDrivenMethod}{greybox\_ekf}
\providecommand{\CthreeMismatchThirtySeqVdpBestLearned}{}\renewcommand{\CthreeMismatchThirtySeqVdpBestLearned}{0.0205}
\providecommand{\CthreeMismatchThirtySeqVdpBestWorking}{}\renewcommand{\CthreeMismatchThirtySeqVdpBestWorking}{0.0578}
\providecommand{\CthreeMismatchThirtySeqVdpClimatology}{}\renewcommand{\CthreeMismatchThirtySeqVdpClimatology}{1.16}
\providecommand{\CthreeMismatchThirtySeqVdpFailedFilters}{}\renewcommand{\CthreeMismatchThirtySeqVdpFailedFilters}{0/8}
\providecommand{\CthreeMismatchThirtySeqVdpVsDataDriven}{}\renewcommand{\CthreeMismatchThirtySeqVdpVsDataDriven}{+110\%}
\providecommand{\CthreeMismatchThirtySeqVdpVsWorking}{}\renewcommand{\CthreeMismatchThirtySeqVdpVsWorking}{-65\%}
\providecommand{\CthreeMismatchThirtySeqVsAny}{}\renewcommand{\CthreeMismatchThirtySeqVsAny}{5/5}
\providecommand{\CthreeMismatchThirtySeqVsDataDriven}{}\renewcommand{\CthreeMismatchThirtySeqVsDataDriven}{0/5}
\providecommand{\CthreeMismatchThirtySeqVsNominal}{}\renewcommand{\CthreeMismatchThirtySeqVsNominal}{5/5}
\providecommand{\CthreeMismatchThirtySeqVsRobust}{}\renewcommand{\CthreeMismatchThirtySeqVsRobust}{5/5}
\providecommand{\CthreeMismatchThirtySeqVsWorking}{}\renewcommand{\CthreeMismatchThirtySeqVsWorking}{5/5}
\providecommand{\CthreeMismatchTwentyIdBatteryBestDataDriven}{}\renewcommand{\CthreeMismatchTwentyIdBatteryBestDataDriven}{0.00977}
\providecommand{\CthreeMismatchTwentyIdBatteryBestDataDrivenMethod}{}\renewcommand{\CthreeMismatchTwentyIdBatteryBestDataDrivenMethod}{greybox\_ekf}
\providecommand{\CthreeMismatchTwentyIdBatteryBestLearned}{}\renewcommand{\CthreeMismatchTwentyIdBatteryBestLearned}{0.043}
\providecommand{\CthreeMismatchTwentyIdBatteryBestWorking}{}\renewcommand{\CthreeMismatchTwentyIdBatteryBestWorking}{0.117}
\providecommand{\CthreeMismatchTwentyIdBatteryClimatology}{}\renewcommand{\CthreeMismatchTwentyIdBatteryClimatology}{0.739}
\providecommand{\CthreeMismatchTwentyIdBatteryFailedFilters}{}\renewcommand{\CthreeMismatchTwentyIdBatteryFailedFilters}{1/8}
\providecommand{\CthreeMismatchTwentyIdBatteryVsDataDriven}{}\renewcommand{\CthreeMismatchTwentyIdBatteryVsDataDriven}{+340\%}
\providecommand{\CthreeMismatchTwentyIdBatteryVsWorking}{}\renewcommand{\CthreeMismatchTwentyIdBatteryVsWorking}{-63\%}
\providecommand{\CthreeMismatchTwentyIdBioreactorBestDataDriven}{}\renewcommand{\CthreeMismatchTwentyIdBioreactorBestDataDriven}{0.0149}
\providecommand{\CthreeMismatchTwentyIdBioreactorBestDataDrivenMethod}{}\renewcommand{\CthreeMismatchTwentyIdBioreactorBestDataDrivenMethod}{nlox}
\providecommand{\CthreeMismatchTwentyIdBioreactorBestLearned}{}\renewcommand{\CthreeMismatchTwentyIdBioreactorBestLearned}{0.024}
\providecommand{\CthreeMismatchTwentyIdBioreactorBestWorking}{}\renewcommand{\CthreeMismatchTwentyIdBioreactorBestWorking}{0.466}
\providecommand{\CthreeMismatchTwentyIdBioreactorClimatology}{}\renewcommand{\CthreeMismatchTwentyIdBioreactorClimatology}{0.718}
\providecommand{\CthreeMismatchTwentyIdBioreactorFailedFilters}{}\renewcommand{\CthreeMismatchTwentyIdBioreactorFailedFilters}{6/8}
\providecommand{\CthreeMismatchTwentyIdBioreactorVsDataDriven}{}\renewcommand{\CthreeMismatchTwentyIdBioreactorVsDataDriven}{+64\%}
\providecommand{\CthreeMismatchTwentyIdBioreactorVsWorking}{}\renewcommand{\CthreeMismatchTwentyIdBioreactorVsWorking}{-95\%}
\providecommand{\CthreeMismatchTwentyIdLorenzBestDataDriven}{}\renewcommand{\CthreeMismatchTwentyIdLorenzBestDataDriven}{0.00891}
\providecommand{\CthreeMismatchTwentyIdLorenzBestDataDrivenMethod}{}\renewcommand{\CthreeMismatchTwentyIdLorenzBestDataDrivenMethod}{nlox}
\providecommand{\CthreeMismatchTwentyIdLorenzBestLearned}{}\renewcommand{\CthreeMismatchTwentyIdLorenzBestLearned}{0.0069}
\providecommand{\CthreeMismatchTwentyIdLorenzBestWorking}{}\renewcommand{\CthreeMismatchTwentyIdLorenzBestWorking}{0.0146}
\providecommand{\CthreeMismatchTwentyIdLorenzClimatology}{}\renewcommand{\CthreeMismatchTwentyIdLorenzClimatology}{0.616}
\providecommand{\CthreeMismatchTwentyIdLorenzFailedFilters}{}\renewcommand{\CthreeMismatchTwentyIdLorenzFailedFilters}{0/8}
\providecommand{\CthreeMismatchTwentyIdLorenzVsDataDriven}{}\renewcommand{\CthreeMismatchTwentyIdLorenzVsDataDriven}{-22\%}
\providecommand{\CthreeMismatchTwentyIdLorenzVsWorking}{}\renewcommand{\CthreeMismatchTwentyIdLorenzVsWorking}{-53\%}
\providecommand{\CthreeMismatchTwentyIdQuadTankBestDataDriven}{}\renewcommand{\CthreeMismatchTwentyIdQuadTankBestDataDriven}{0.0135}
\providecommand{\CthreeMismatchTwentyIdQuadTankBestDataDrivenMethod}{}\renewcommand{\CthreeMismatchTwentyIdQuadTankBestDataDrivenMethod}{greybox\_ekf}
\providecommand{\CthreeMismatchTwentyIdQuadTankBestLearned}{}\renewcommand{\CthreeMismatchTwentyIdQuadTankBestLearned}{0.0255}
\providecommand{\CthreeMismatchTwentyIdQuadTankClimatology}{}\renewcommand{\CthreeMismatchTwentyIdQuadTankClimatology}{0.534}
\providecommand{\CthreeMismatchTwentyIdQuadTankFailedFilters}{}\renewcommand{\CthreeMismatchTwentyIdQuadTankFailedFilters}{8/8}
\providecommand{\CthreeMismatchTwentyIdQuadTankVsDataDriven}{}\renewcommand{\CthreeMismatchTwentyIdQuadTankVsDataDriven}{+89\%}
\providecommand{\CthreeMismatchTwentyIdVdpBestDataDriven}{}\renewcommand{\CthreeMismatchTwentyIdVdpBestDataDriven}{0.00706}
\providecommand{\CthreeMismatchTwentyIdVdpBestDataDrivenMethod}{}\renewcommand{\CthreeMismatchTwentyIdVdpBestDataDrivenMethod}{greybox\_ekf}
\providecommand{\CthreeMismatchTwentyIdVdpBestLearned}{}\renewcommand{\CthreeMismatchTwentyIdVdpBestLearned}{0.0172}
\providecommand{\CthreeMismatchTwentyIdVdpBestWorking}{}\renewcommand{\CthreeMismatchTwentyIdVdpBestWorking}{0.0388}
\providecommand{\CthreeMismatchTwentyIdVdpClimatology}{}\renewcommand{\CthreeMismatchTwentyIdVdpClimatology}{1.07}
\providecommand{\CthreeMismatchTwentyIdVdpFailedFilters}{}\renewcommand{\CthreeMismatchTwentyIdVdpFailedFilters}{0/8}
\providecommand{\CthreeMismatchTwentyIdVdpVsDataDriven}{}\renewcommand{\CthreeMismatchTwentyIdVdpVsDataDriven}{+144\%}
\providecommand{\CthreeMismatchTwentyIdVdpVsWorking}{}\renewcommand{\CthreeMismatchTwentyIdVdpVsWorking}{-56\%}
\providecommand{\CthreeMismatchTwentyIdVsAny}{}\renewcommand{\CthreeMismatchTwentyIdVsAny}{5/5}
\providecommand{\CthreeMismatchTwentyIdVsDataDriven}{}\renewcommand{\CthreeMismatchTwentyIdVsDataDriven}{1/5}
\providecommand{\CthreeMismatchTwentyIdVsNominal}{}\renewcommand{\CthreeMismatchTwentyIdVsNominal}{5/5}
\providecommand{\CthreeMismatchTwentyIdVsRobust}{}\renewcommand{\CthreeMismatchTwentyIdVsRobust}{5/5}
\providecommand{\CthreeMismatchTwentyIdVsWorking}{}\renewcommand{\CthreeMismatchTwentyIdVsWorking}{5/5}
\providecommand{\CthreeMismatchTwentySeqBatteryBestDataDriven}{}\renewcommand{\CthreeMismatchTwentySeqBatteryBestDataDriven}{0.0111}
\providecommand{\CthreeMismatchTwentySeqBatteryBestDataDrivenMethod}{}\renewcommand{\CthreeMismatchTwentySeqBatteryBestDataDrivenMethod}{greybox\_ekf}
\providecommand{\CthreeMismatchTwentySeqBatteryBestLearned}{}\renewcommand{\CthreeMismatchTwentySeqBatteryBestLearned}{0.0569}
\providecommand{\CthreeMismatchTwentySeqBatteryBestWorking}{}\renewcommand{\CthreeMismatchTwentySeqBatteryBestWorking}{0.12}
\providecommand{\CthreeMismatchTwentySeqBatteryClimatology}{}\renewcommand{\CthreeMismatchTwentySeqBatteryClimatology}{0.77}
\providecommand{\CthreeMismatchTwentySeqBatteryFailedFilters}{}\renewcommand{\CthreeMismatchTwentySeqBatteryFailedFilters}{1/8}
\providecommand{\CthreeMismatchTwentySeqBatteryVsDataDriven}{}\renewcommand{\CthreeMismatchTwentySeqBatteryVsDataDriven}{+413\%}
\providecommand{\CthreeMismatchTwentySeqBatteryVsWorking}{}\renewcommand{\CthreeMismatchTwentySeqBatteryVsWorking}{-53\%}
\providecommand{\CthreeMismatchTwentySeqBioreactorBestDataDriven}{}\renewcommand{\CthreeMismatchTwentySeqBioreactorBestDataDriven}{0.0291}
\providecommand{\CthreeMismatchTwentySeqBioreactorBestDataDrivenMethod}{}\renewcommand{\CthreeMismatchTwentySeqBioreactorBestDataDrivenMethod}{greybox\_ekf\_oracle}
\providecommand{\CthreeMismatchTwentySeqBioreactorBestLearned}{}\renewcommand{\CthreeMismatchTwentySeqBioreactorBestLearned}{0.0783}
\providecommand{\CthreeMismatchTwentySeqBioreactorBestWorking}{}\renewcommand{\CthreeMismatchTwentySeqBioreactorBestWorking}{0.401}
\providecommand{\CthreeMismatchTwentySeqBioreactorClimatology}{}\renewcommand{\CthreeMismatchTwentySeqBioreactorClimatology}{0.769}
\providecommand{\CthreeMismatchTwentySeqBioreactorFailedFilters}{}\renewcommand{\CthreeMismatchTwentySeqBioreactorFailedFilters}{6/8}
\providecommand{\CthreeMismatchTwentySeqBioreactorVsDataDriven}{}\renewcommand{\CthreeMismatchTwentySeqBioreactorVsDataDriven}{+169\%}
\providecommand{\CthreeMismatchTwentySeqBioreactorVsWorking}{}\renewcommand{\CthreeMismatchTwentySeqBioreactorVsWorking}{-80\%}
\providecommand{\CthreeMismatchTwentySeqLorenzBestDataDriven}{}\renewcommand{\CthreeMismatchTwentySeqLorenzBestDataDriven}{0.0108}
\providecommand{\CthreeMismatchTwentySeqLorenzBestDataDrivenMethod}{}\renewcommand{\CthreeMismatchTwentySeqLorenzBestDataDrivenMethod}{nlox}
\providecommand{\CthreeMismatchTwentySeqLorenzBestLearned}{}\renewcommand{\CthreeMismatchTwentySeqLorenzBestLearned}{0.00772}
\providecommand{\CthreeMismatchTwentySeqLorenzBestWorking}{}\renewcommand{\CthreeMismatchTwentySeqLorenzBestWorking}{0.0248}
\providecommand{\CthreeMismatchTwentySeqLorenzClimatology}{}\renewcommand{\CthreeMismatchTwentySeqLorenzClimatology}{0.616}
\providecommand{\CthreeMismatchTwentySeqLorenzFailedFilters}{}\renewcommand{\CthreeMismatchTwentySeqLorenzFailedFilters}{0/8}
\providecommand{\CthreeMismatchTwentySeqLorenzVsDataDriven}{}\renewcommand{\CthreeMismatchTwentySeqLorenzVsDataDriven}{-28\%}
\providecommand{\CthreeMismatchTwentySeqLorenzVsWorking}{}\renewcommand{\CthreeMismatchTwentySeqLorenzVsWorking}{-69\%}
\providecommand{\CthreeMismatchTwentySeqQuadTankBestDataDriven}{}\renewcommand{\CthreeMismatchTwentySeqQuadTankBestDataDriven}{0.0131}
\providecommand{\CthreeMismatchTwentySeqQuadTankBestDataDrivenMethod}{}\renewcommand{\CthreeMismatchTwentySeqQuadTankBestDataDrivenMethod}{greybox\_ukf}
\providecommand{\CthreeMismatchTwentySeqQuadTankBestLearned}{}\renewcommand{\CthreeMismatchTwentySeqQuadTankBestLearned}{0.0421}
\providecommand{\CthreeMismatchTwentySeqQuadTankClimatology}{}\renewcommand{\CthreeMismatchTwentySeqQuadTankClimatology}{0.587}
\providecommand{\CthreeMismatchTwentySeqQuadTankFailedFilters}{}\renewcommand{\CthreeMismatchTwentySeqQuadTankFailedFilters}{8/8}
\providecommand{\CthreeMismatchTwentySeqQuadTankVsDataDriven}{}\renewcommand{\CthreeMismatchTwentySeqQuadTankVsDataDriven}{+222\%}
\providecommand{\CthreeMismatchTwentySeqVdpBestDataDriven}{}\renewcommand{\CthreeMismatchTwentySeqVdpBestDataDriven}{0.00755}
\providecommand{\CthreeMismatchTwentySeqVdpBestDataDrivenMethod}{}\renewcommand{\CthreeMismatchTwentySeqVdpBestDataDrivenMethod}{greybox\_ekf}
\providecommand{\CthreeMismatchTwentySeqVdpBestLearned}{}\renewcommand{\CthreeMismatchTwentySeqVdpBestLearned}{0.0199}
\providecommand{\CthreeMismatchTwentySeqVdpBestWorking}{}\renewcommand{\CthreeMismatchTwentySeqVdpBestWorking}{0.0429}
\providecommand{\CthreeMismatchTwentySeqVdpClimatology}{}\renewcommand{\CthreeMismatchTwentySeqVdpClimatology}{1.09}
\providecommand{\CthreeMismatchTwentySeqVdpFailedFilters}{}\renewcommand{\CthreeMismatchTwentySeqVdpFailedFilters}{0/8}
\providecommand{\CthreeMismatchTwentySeqVdpVsDataDriven}{}\renewcommand{\CthreeMismatchTwentySeqVdpVsDataDriven}{+165\%}
\providecommand{\CthreeMismatchTwentySeqVdpVsWorking}{}\renewcommand{\CthreeMismatchTwentySeqVdpVsWorking}{-54\%}
\providecommand{\CthreeMismatchTwentySeqVsAny}{}\renewcommand{\CthreeMismatchTwentySeqVsAny}{5/5}
\providecommand{\CthreeMismatchTwentySeqVsDataDriven}{}\renewcommand{\CthreeMismatchTwentySeqVsDataDriven}{1/5}
\providecommand{\CthreeMismatchTwentySeqVsNominal}{}\renewcommand{\CthreeMismatchTwentySeqVsNominal}{5/5}
\providecommand{\CthreeMismatchTwentySeqVsRobust}{}\renewcommand{\CthreeMismatchTwentySeqVsRobust}{5/5}
\providecommand{\CthreeMismatchTwentySeqVsWorking}{}\renewcommand{\CthreeMismatchTwentySeqVsWorking}{5/5}
\providecommand{\CthreeModelWrongWins}{}\renewcommand{\CthreeModelWrongWins}{28/30}  % model-wrong cells where best KKL beats best working filter
\providecommand{\CthreeNoiseCtIdBatteryBestDataDriven}{}\renewcommand{\CthreeNoiseCtIdBatteryBestDataDriven}{0.0109}
\providecommand{\CthreeNoiseCtIdBatteryBestDataDrivenMethod}{}\renewcommand{\CthreeNoiseCtIdBatteryBestDataDrivenMethod}{greybox\_ekf\_oracle}
\providecommand{\CthreeNoiseCtIdBatteryBestLearned}{}\renewcommand{\CthreeNoiseCtIdBatteryBestLearned}{0.0476}
\providecommand{\CthreeNoiseCtIdBatteryBestWorking}{}\renewcommand{\CthreeNoiseCtIdBatteryBestWorking}{0.0109}
\providecommand{\CthreeNoiseCtIdBatteryClimatology}{}\renewcommand{\CthreeNoiseCtIdBatteryClimatology}{1}
\providecommand{\CthreeNoiseCtIdBatteryFailedFilters}{}\renewcommand{\CthreeNoiseCtIdBatteryFailedFilters}{1/8}
\providecommand{\CthreeNoiseCtIdBatteryVsDataDriven}{}\renewcommand{\CthreeNoiseCtIdBatteryVsDataDriven}{+337\%}
\providecommand{\CthreeNoiseCtIdBatteryVsWorking}{}\renewcommand{\CthreeNoiseCtIdBatteryVsWorking}{+338\%}
\providecommand{\CthreeNoiseCtIdBioreactorBestDataDriven}{}\renewcommand{\CthreeNoiseCtIdBioreactorBestDataDriven}{0.00626}
\providecommand{\CthreeNoiseCtIdBioreactorBestDataDrivenMethod}{}\renewcommand{\CthreeNoiseCtIdBioreactorBestDataDrivenMethod}{greybox\_ekf\_oracle}
\providecommand{\CthreeNoiseCtIdBioreactorBestLearned}{}\renewcommand{\CthreeNoiseCtIdBioreactorBestLearned}{0.0227}
\providecommand{\CthreeNoiseCtIdBioreactorBestWorking}{}\renewcommand{\CthreeNoiseCtIdBioreactorBestWorking}{0.00625}
\providecommand{\CthreeNoiseCtIdBioreactorClimatology}{}\renewcommand{\CthreeNoiseCtIdBioreactorClimatology}{1}
\providecommand{\CthreeNoiseCtIdBioreactorFailedFilters}{}\renewcommand{\CthreeNoiseCtIdBioreactorFailedFilters}{0/8}
\providecommand{\CthreeNoiseCtIdBioreactorVsDataDriven}{}\renewcommand{\CthreeNoiseCtIdBioreactorVsDataDriven}{+262\%}
\providecommand{\CthreeNoiseCtIdBioreactorVsWorking}{}\renewcommand{\CthreeNoiseCtIdBioreactorVsWorking}{+263\%}
\providecommand{\CthreeNoiseCtIdLorenzBestDataDriven}{}\renewcommand{\CthreeNoiseCtIdLorenzBestDataDriven}{0.0188}
\providecommand{\CthreeNoiseCtIdLorenzBestDataDrivenMethod}{}\renewcommand{\CthreeNoiseCtIdLorenzBestDataDrivenMethod}{greybox\_ekf\_oracle}
\providecommand{\CthreeNoiseCtIdLorenzBestLearned}{}\renewcommand{\CthreeNoiseCtIdLorenzBestLearned}{0.0441}
\providecommand{\CthreeNoiseCtIdLorenzBestWorking}{}\renewcommand{\CthreeNoiseCtIdLorenzBestWorking}{0.0181}
\providecommand{\CthreeNoiseCtIdLorenzClimatology}{}\renewcommand{\CthreeNoiseCtIdLorenzClimatology}{1}
\providecommand{\CthreeNoiseCtIdLorenzFailedFilters}{}\renewcommand{\CthreeNoiseCtIdLorenzFailedFilters}{0/8}
\providecommand{\CthreeNoiseCtIdLorenzVsDataDriven}{}\renewcommand{\CthreeNoiseCtIdLorenzVsDataDriven}{+134\%}
\providecommand{\CthreeNoiseCtIdLorenzVsWorking}{}\renewcommand{\CthreeNoiseCtIdLorenzVsWorking}{+143\%}
\providecommand{\CthreeNoiseCtIdQuadTankBestDataDriven}{}\renewcommand{\CthreeNoiseCtIdQuadTankBestDataDriven}{0.0173}
\providecommand{\CthreeNoiseCtIdQuadTankBestDataDrivenMethod}{}\renewcommand{\CthreeNoiseCtIdQuadTankBestDataDrivenMethod}{greybox\_ekf\_oracle}
\providecommand{\CthreeNoiseCtIdQuadTankBestLearned}{}\renewcommand{\CthreeNoiseCtIdQuadTankBestLearned}{0.0547}
\providecommand{\CthreeNoiseCtIdQuadTankBestWorking}{}\renewcommand{\CthreeNoiseCtIdQuadTankBestWorking}{0.0172}
\providecommand{\CthreeNoiseCtIdQuadTankClimatology}{}\renewcommand{\CthreeNoiseCtIdQuadTankClimatology}{1}
\providecommand{\CthreeNoiseCtIdQuadTankFailedFilters}{}\renewcommand{\CthreeNoiseCtIdQuadTankFailedFilters}{0/8}
\providecommand{\CthreeNoiseCtIdQuadTankVsDataDriven}{}\renewcommand{\CthreeNoiseCtIdQuadTankVsDataDriven}{+215\%}
\providecommand{\CthreeNoiseCtIdQuadTankVsWorking}{}\renewcommand{\CthreeNoiseCtIdQuadTankVsWorking}{+218\%}
\providecommand{\CthreeNoiseCtIdVdpBestDataDriven}{}\renewcommand{\CthreeNoiseCtIdVdpBestDataDriven}{0.0132}
\providecommand{\CthreeNoiseCtIdVdpBestDataDrivenMethod}{}\renewcommand{\CthreeNoiseCtIdVdpBestDataDrivenMethod}{greybox\_ekf\_oracle}
\providecommand{\CthreeNoiseCtIdVdpBestLearned}{}\renewcommand{\CthreeNoiseCtIdVdpBestLearned}{0.0314}
\providecommand{\CthreeNoiseCtIdVdpBestWorking}{}\renewcommand{\CthreeNoiseCtIdVdpBestWorking}{0.0132}
\providecommand{\CthreeNoiseCtIdVdpClimatology}{}\renewcommand{\CthreeNoiseCtIdVdpClimatology}{1}
\providecommand{\CthreeNoiseCtIdVdpFailedFilters}{}\renewcommand{\CthreeNoiseCtIdVdpFailedFilters}{0/8}
\providecommand{\CthreeNoiseCtIdVdpVsDataDriven}{}\renewcommand{\CthreeNoiseCtIdVdpVsDataDriven}{+137\%}
\providecommand{\CthreeNoiseCtIdVdpVsWorking}{}\renewcommand{\CthreeNoiseCtIdVdpVsWorking}{+137\%}
\providecommand{\CthreeNoiseCtIdVsAny}{}\renewcommand{\CthreeNoiseCtIdVsAny}{0/5}
\providecommand{\CthreeNoiseCtIdVsDataDriven}{}\renewcommand{\CthreeNoiseCtIdVsDataDriven}{0/5}
\providecommand{\CthreeNoiseCtIdVsNominal}{}\renewcommand{\CthreeNoiseCtIdVsNominal}{0/5}
\providecommand{\CthreeNoiseCtIdVsRobust}{}\renewcommand{\CthreeNoiseCtIdVsRobust}{0/5}
\providecommand{\CthreeNoiseCtIdVsWorking}{}\renewcommand{\CthreeNoiseCtIdVsWorking}{0/5}
\providecommand{\CthreeNoiseCtSeqBatteryBestDataDriven}{}\renewcommand{\CthreeNoiseCtSeqBatteryBestDataDriven}{0.00989}
\providecommand{\CthreeNoiseCtSeqBatteryBestDataDrivenMethod}{}\renewcommand{\CthreeNoiseCtSeqBatteryBestDataDrivenMethod}{greybox\_ekf\_oracle}
\providecommand{\CthreeNoiseCtSeqBatteryBestLearned}{}\renewcommand{\CthreeNoiseCtSeqBatteryBestLearned}{0.0595}
\providecommand{\CthreeNoiseCtSeqBatteryBestWorking}{}\renewcommand{\CthreeNoiseCtSeqBatteryBestWorking}{0.00989}
\providecommand{\CthreeNoiseCtSeqBatteryClimatology}{}\renewcommand{\CthreeNoiseCtSeqBatteryClimatology}{0.78}
\providecommand{\CthreeNoiseCtSeqBatteryFailedFilters}{}\renewcommand{\CthreeNoiseCtSeqBatteryFailedFilters}{1/8}
\providecommand{\CthreeNoiseCtSeqBatteryVsDataDriven}{}\renewcommand{\CthreeNoiseCtSeqBatteryVsDataDriven}{+502\%}
\providecommand{\CthreeNoiseCtSeqBatteryVsWorking}{}\renewcommand{\CthreeNoiseCtSeqBatteryVsWorking}{+502\%}
\providecommand{\CthreeNoiseCtSeqBioreactorBestDataDriven}{}\renewcommand{\CthreeNoiseCtSeqBioreactorBestDataDriven}{0.00625}
\providecommand{\CthreeNoiseCtSeqBioreactorBestDataDrivenMethod}{}\renewcommand{\CthreeNoiseCtSeqBioreactorBestDataDrivenMethod}{greybox\_ekf\_oracle}
\providecommand{\CthreeNoiseCtSeqBioreactorBestLearned}{}\renewcommand{\CthreeNoiseCtSeqBioreactorBestLearned}{0.0645}
\providecommand{\CthreeNoiseCtSeqBioreactorBestWorking}{}\renewcommand{\CthreeNoiseCtSeqBioreactorBestWorking}{0.00624}
\providecommand{\CthreeNoiseCtSeqBioreactorClimatology}{}\renewcommand{\CthreeNoiseCtSeqBioreactorClimatology}{0.928}
\providecommand{\CthreeNoiseCtSeqBioreactorFailedFilters}{}\renewcommand{\CthreeNoiseCtSeqBioreactorFailedFilters}{0/8}
\providecommand{\CthreeNoiseCtSeqBioreactorVsDataDriven}{}\renewcommand{\CthreeNoiseCtSeqBioreactorVsDataDriven}{+932\%}
\providecommand{\CthreeNoiseCtSeqBioreactorVsWorking}{}\renewcommand{\CthreeNoiseCtSeqBioreactorVsWorking}{+933\%}
\providecommand{\CthreeNoiseCtSeqLorenzBestDataDriven}{}\renewcommand{\CthreeNoiseCtSeqLorenzBestDataDriven}{0.0183}
\providecommand{\CthreeNoiseCtSeqLorenzBestDataDrivenMethod}{}\renewcommand{\CthreeNoiseCtSeqLorenzBestDataDrivenMethod}{greybox\_ekf\_oracle}
\providecommand{\CthreeNoiseCtSeqLorenzBestLearned}{}\renewcommand{\CthreeNoiseCtSeqLorenzBestLearned}{0.0406}
\providecommand{\CthreeNoiseCtSeqLorenzBestWorking}{}\renewcommand{\CthreeNoiseCtSeqLorenzBestWorking}{0.0177}
\providecommand{\CthreeNoiseCtSeqLorenzClimatology}{}\renewcommand{\CthreeNoiseCtSeqLorenzClimatology}{1}
\providecommand{\CthreeNoiseCtSeqLorenzFailedFilters}{}\renewcommand{\CthreeNoiseCtSeqLorenzFailedFilters}{0/8}
\providecommand{\CthreeNoiseCtSeqLorenzVsDataDriven}{}\renewcommand{\CthreeNoiseCtSeqLorenzVsDataDriven}{+122\%}
\providecommand{\CthreeNoiseCtSeqLorenzVsWorking}{}\renewcommand{\CthreeNoiseCtSeqLorenzVsWorking}{+129\%}
\providecommand{\CthreeNoiseCtSeqQuadTankBestDataDriven}{}\renewcommand{\CthreeNoiseCtSeqQuadTankBestDataDriven}{0.017}
\providecommand{\CthreeNoiseCtSeqQuadTankBestDataDrivenMethod}{}\renewcommand{\CthreeNoiseCtSeqQuadTankBestDataDrivenMethod}{greybox\_ekf\_oracle}
\providecommand{\CthreeNoiseCtSeqQuadTankBestLearned}{}\renewcommand{\CthreeNoiseCtSeqQuadTankBestLearned}{0.0696}
\providecommand{\CthreeNoiseCtSeqQuadTankBestWorking}{}\renewcommand{\CthreeNoiseCtSeqQuadTankBestWorking}{0.0169}
\providecommand{\CthreeNoiseCtSeqQuadTankClimatology}{}\renewcommand{\CthreeNoiseCtSeqQuadTankClimatology}{0.651}
\providecommand{\CthreeNoiseCtSeqQuadTankFailedFilters}{}\renewcommand{\CthreeNoiseCtSeqQuadTankFailedFilters}{0/8}
\providecommand{\CthreeNoiseCtSeqQuadTankVsDataDriven}{}\renewcommand{\CthreeNoiseCtSeqQuadTankVsDataDriven}{+310\%}
\providecommand{\CthreeNoiseCtSeqQuadTankVsWorking}{}\renewcommand{\CthreeNoiseCtSeqQuadTankVsWorking}{+313\%}
\providecommand{\CthreeNoiseCtSeqVdpBestDataDriven}{}\renewcommand{\CthreeNoiseCtSeqVdpBestDataDriven}{0.0137}
\providecommand{\CthreeNoiseCtSeqVdpBestDataDrivenMethod}{}\renewcommand{\CthreeNoiseCtSeqVdpBestDataDrivenMethod}{greybox\_ekf\_oracle}
\providecommand{\CthreeNoiseCtSeqVdpBestLearned}{}\renewcommand{\CthreeNoiseCtSeqVdpBestLearned}{0.035}
\providecommand{\CthreeNoiseCtSeqVdpBestWorking}{}\renewcommand{\CthreeNoiseCtSeqVdpBestWorking}{0.0134}
\providecommand{\CthreeNoiseCtSeqVdpClimatology}{}\renewcommand{\CthreeNoiseCtSeqVdpClimatology}{1.12}
\providecommand{\CthreeNoiseCtSeqVdpFailedFilters}{}\renewcommand{\CthreeNoiseCtSeqVdpFailedFilters}{0/8}
\providecommand{\CthreeNoiseCtSeqVdpVsDataDriven}{}\renewcommand{\CthreeNoiseCtSeqVdpVsDataDriven}{+156\%}
\providecommand{\CthreeNoiseCtSeqVdpVsWorking}{}\renewcommand{\CthreeNoiseCtSeqVdpVsWorking}{+161\%}
\providecommand{\CthreeNoiseCtSeqVsAny}{}\renewcommand{\CthreeNoiseCtSeqVsAny}{0/5}
\providecommand{\CthreeNoiseCtSeqVsDataDriven}{}\renewcommand{\CthreeNoiseCtSeqVsDataDriven}{0/5}
\providecommand{\CthreeNoiseCtSeqVsNominal}{}\renewcommand{\CthreeNoiseCtSeqVsNominal}{0/5}
\providecommand{\CthreeNoiseCtSeqVsRobust}{}\renewcommand{\CthreeNoiseCtSeqVsRobust}{0/5}
\providecommand{\CthreeNoiseCtSeqVsWorking}{}\renewcommand{\CthreeNoiseCtSeqVsWorking}{0/5}
\providecommand{\CthreePartialIdBatteryBestDataDriven}{}\renewcommand{\CthreePartialIdBatteryBestDataDriven}{0.027}
\providecommand{\CthreePartialIdBatteryBestDataDrivenMethod}{}\renewcommand{\CthreePartialIdBatteryBestDataDrivenMethod}{greybox\_ukf}
\providecommand{\CthreePartialIdBatteryBestLearned}{}\renewcommand{\CthreePartialIdBatteryBestLearned}{0.0476}
\providecommand{\CthreePartialIdBatteryBestWorking}{}\renewcommand{\CthreePartialIdBatteryBestWorking}{0.397}
\providecommand{\CthreePartialIdBatteryClimatology}{}\renewcommand{\CthreePartialIdBatteryClimatology}{1}
\providecommand{\CthreePartialIdBatteryFailedFilters}{}\renewcommand{\CthreePartialIdBatteryFailedFilters}{1/8}
\providecommand{\CthreePartialIdBatteryVsDataDriven}{}\renewcommand{\CthreePartialIdBatteryVsDataDriven}{+77\%}
\providecommand{\CthreePartialIdBatteryVsWorking}{}\renewcommand{\CthreePartialIdBatteryVsWorking}{-88\%}
\providecommand{\CthreePartialIdBioreactorBestDataDriven}{}\renewcommand{\CthreePartialIdBioreactorBestDataDriven}{0.017}
\providecommand{\CthreePartialIdBioreactorBestDataDrivenMethod}{}\renewcommand{\CthreePartialIdBioreactorBestDataDrivenMethod}{nlox}
\providecommand{\CthreePartialIdBioreactorBestLearned}{}\renewcommand{\CthreePartialIdBioreactorBestLearned}{0.0239}
\providecommand{\CthreePartialIdBioreactorClimatology}{}\renewcommand{\CthreePartialIdBioreactorClimatology}{1}
\providecommand{\CthreePartialIdBioreactorFailedFilters}{}\renewcommand{\CthreePartialIdBioreactorFailedFilters}{8/8}
\providecommand{\CthreePartialIdBioreactorVsDataDriven}{}\renewcommand{\CthreePartialIdBioreactorVsDataDriven}{+42\%}
\providecommand{\CthreePartialIdLorenzBestDataDriven}{}\renewcommand{\CthreePartialIdLorenzBestDataDriven}{0.0285}
\providecommand{\CthreePartialIdLorenzBestDataDrivenMethod}{}\renewcommand{\CthreePartialIdLorenzBestDataDrivenMethod}{nlox}
\providecommand{\CthreePartialIdLorenzBestLearned}{}\renewcommand{\CthreePartialIdLorenzBestLearned}{0.0298}
\providecommand{\CthreePartialIdLorenzBestWorking}{}\renewcommand{\CthreePartialIdLorenzBestWorking}{0.533}
\providecommand{\CthreePartialIdLorenzClimatology}{}\renewcommand{\CthreePartialIdLorenzClimatology}{1}
\providecommand{\CthreePartialIdLorenzFailedFilters}{}\renewcommand{\CthreePartialIdLorenzFailedFilters}{5/8}
\providecommand{\CthreePartialIdLorenzVsDataDriven}{}\renewcommand{\CthreePartialIdLorenzVsDataDriven}{+4\%}
\providecommand{\CthreePartialIdLorenzVsWorking}{}\renewcommand{\CthreePartialIdLorenzVsWorking}{-94\%}
\providecommand{\CthreePartialIdQuadTankBestDataDriven}{}\renewcommand{\CthreePartialIdQuadTankBestDataDriven}{0.015}
\providecommand{\CthreePartialIdQuadTankBestDataDrivenMethod}{}\renewcommand{\CthreePartialIdQuadTankBestDataDrivenMethod}{greybox\_ekf}
\providecommand{\CthreePartialIdQuadTankBestLearned}{}\renewcommand{\CthreePartialIdQuadTankBestLearned}{0.0517}
\providecommand{\CthreePartialIdQuadTankBestWorking}{}\renewcommand{\CthreePartialIdQuadTankBestWorking}{0.0285}
\providecommand{\CthreePartialIdQuadTankClimatology}{}\renewcommand{\CthreePartialIdQuadTankClimatology}{1}
\providecommand{\CthreePartialIdQuadTankFailedFilters}{}\renewcommand{\CthreePartialIdQuadTankFailedFilters}{0/8}
\providecommand{\CthreePartialIdQuadTankVsDataDriven}{}\renewcommand{\CthreePartialIdQuadTankVsDataDriven}{+244\%}
\providecommand{\CthreePartialIdQuadTankVsWorking}{}\renewcommand{\CthreePartialIdQuadTankVsWorking}{+82\%}
\providecommand{\CthreePartialIdVdpBestDataDriven}{}\renewcommand{\CthreePartialIdVdpBestDataDriven}{0.0196}
\providecommand{\CthreePartialIdVdpBestDataDrivenMethod}{}\renewcommand{\CthreePartialIdVdpBestDataDrivenMethod}{greybox\_ekf}
\providecommand{\CthreePartialIdVdpBestLearned}{}\renewcommand{\CthreePartialIdVdpBestLearned}{0.0157}
\providecommand{\CthreePartialIdVdpBestWorking}{}\renewcommand{\CthreePartialIdVdpBestWorking}{0.064}
\providecommand{\CthreePartialIdVdpClimatology}{}\renewcommand{\CthreePartialIdVdpClimatology}{1}
\providecommand{\CthreePartialIdVdpFailedFilters}{}\renewcommand{\CthreePartialIdVdpFailedFilters}{0/8}
\providecommand{\CthreePartialIdVdpVsDataDriven}{}\renewcommand{\CthreePartialIdVdpVsDataDriven}{-18\%}
\providecommand{\CthreePartialIdVdpVsWorking}{}\renewcommand{\CthreePartialIdVdpVsWorking}{-75\%}
\providecommand{\CthreePartialIdVsAny}{}\renewcommand{\CthreePartialIdVsAny}{4/5}
\providecommand{\CthreePartialIdVsDataDriven}{}\renewcommand{\CthreePartialIdVsDataDriven}{0/5}
\providecommand{\CthreePartialIdVsNominal}{}\renewcommand{\CthreePartialIdVsNominal}{4/5}
\providecommand{\CthreePartialIdVsRobust}{}\renewcommand{\CthreePartialIdVsRobust}{4/5}
\providecommand{\CthreePartialIdVsWorking}{}\renewcommand{\CthreePartialIdVsWorking}{4/5}
\providecommand{\CthreePartialPatternHolds}{}\renewcommand{\CthreePartialPatternHolds}{4/4}  % hypothesis pattern holds on k of the named plants in the current dev set (5 pre-registered; duffing excluded -> 4)
\providecommand{\CthreePartialSeqBatteryBestDataDriven}{}\renewcommand{\CthreePartialSeqBatteryBestDataDriven}{0.0191}
\providecommand{\CthreePartialSeqBatteryBestDataDrivenMethod}{}\renewcommand{\CthreePartialSeqBatteryBestDataDrivenMethod}{greybox\_ekf}
\providecommand{\CthreePartialSeqBatteryBestLearned}{}\renewcommand{\CthreePartialSeqBatteryBestLearned}{0.0744}
\providecommand{\CthreePartialSeqBatteryBestWorking}{}\renewcommand{\CthreePartialSeqBatteryBestWorking}{0.237}
\providecommand{\CthreePartialSeqBatteryClimatology}{}\renewcommand{\CthreePartialSeqBatteryClimatology}{0.78}
\providecommand{\CthreePartialSeqBatteryFailedFilters}{}\renewcommand{\CthreePartialSeqBatteryFailedFilters}{1/8}
\providecommand{\CthreePartialSeqBatteryVsDataDriven}{}\renewcommand{\CthreePartialSeqBatteryVsDataDriven}{+295\%}
\providecommand{\CthreePartialSeqBatteryVsWorking}{}\renewcommand{\CthreePartialSeqBatteryVsWorking}{-69\%}
\providecommand{\CthreePartialSeqBioreactorBestDataDriven}{}\renewcommand{\CthreePartialSeqBioreactorBestDataDriven}{0.0271}
\providecommand{\CthreePartialSeqBioreactorBestDataDrivenMethod}{}\renewcommand{\CthreePartialSeqBioreactorBestDataDrivenMethod}{koopman\_kf}
\providecommand{\CthreePartialSeqBioreactorBestLearned}{}\renewcommand{\CthreePartialSeqBioreactorBestLearned}{0.0753}
\providecommand{\CthreePartialSeqBioreactorClimatology}{}\renewcommand{\CthreePartialSeqBioreactorClimatology}{0.928}
\providecommand{\CthreePartialSeqBioreactorFailedFilters}{}\renewcommand{\CthreePartialSeqBioreactorFailedFilters}{8/8}
\providecommand{\CthreePartialSeqBioreactorVsDataDriven}{}\renewcommand{\CthreePartialSeqBioreactorVsDataDriven}{+179\%}
\providecommand{\CthreePartialSeqLorenzBestDataDriven}{}\renewcommand{\CthreePartialSeqLorenzBestDataDriven}{0.0287}
\providecommand{\CthreePartialSeqLorenzBestDataDrivenMethod}{}\renewcommand{\CthreePartialSeqLorenzBestDataDrivenMethod}{nlox}
\providecommand{\CthreePartialSeqLorenzBestLearned}{}\renewcommand{\CthreePartialSeqLorenzBestLearned}{0.0308}
\providecommand{\CthreePartialSeqLorenzBestWorking}{}\renewcommand{\CthreePartialSeqLorenzBestWorking}{0.529}
\providecommand{\CthreePartialSeqLorenzClimatology}{}\renewcommand{\CthreePartialSeqLorenzClimatology}{1}
\providecommand{\CthreePartialSeqLorenzFailedFilters}{}\renewcommand{\CthreePartialSeqLorenzFailedFilters}{5/8}
\providecommand{\CthreePartialSeqLorenzVsDataDriven}{}\renewcommand{\CthreePartialSeqLorenzVsDataDriven}{+7\%}
\providecommand{\CthreePartialSeqLorenzVsWorking}{}\renewcommand{\CthreePartialSeqLorenzVsWorking}{-94\%}
\providecommand{\CthreePartialSeqQuadTankBestDataDriven}{}\renewcommand{\CthreePartialSeqQuadTankBestDataDriven}{0.0142}
\providecommand{\CthreePartialSeqQuadTankBestDataDrivenMethod}{}\renewcommand{\CthreePartialSeqQuadTankBestDataDrivenMethod}{greybox\_ekf}
\providecommand{\CthreePartialSeqQuadTankBestLearned}{}\renewcommand{\CthreePartialSeqQuadTankBestLearned}{0.0717}
\providecommand{\CthreePartialSeqQuadTankBestWorking}{}\renewcommand{\CthreePartialSeqQuadTankBestWorking}{0.0202}
\providecommand{\CthreePartialSeqQuadTankClimatology}{}\renewcommand{\CthreePartialSeqQuadTankClimatology}{0.651}
\providecommand{\CthreePartialSeqQuadTankFailedFilters}{}\renewcommand{\CthreePartialSeqQuadTankFailedFilters}{0/8}
\providecommand{\CthreePartialSeqQuadTankVsDataDriven}{}\renewcommand{\CthreePartialSeqQuadTankVsDataDriven}{+406\%}
\providecommand{\CthreePartialSeqQuadTankVsWorking}{}\renewcommand{\CthreePartialSeqQuadTankVsWorking}{+255\%}
\providecommand{\CthreePartialSeqVdpBestDataDriven}{}\renewcommand{\CthreePartialSeqVdpBestDataDriven}{0.0257}
\providecommand{\CthreePartialSeqVdpBestDataDrivenMethod}{}\renewcommand{\CthreePartialSeqVdpBestDataDrivenMethod}{greybox\_ekf}
\providecommand{\CthreePartialSeqVdpBestLearned}{}\renewcommand{\CthreePartialSeqVdpBestLearned}{0.0179}
\providecommand{\CthreePartialSeqVdpBestWorking}{}\renewcommand{\CthreePartialSeqVdpBestWorking}{0.0819}
\providecommand{\CthreePartialSeqVdpClimatology}{}\renewcommand{\CthreePartialSeqVdpClimatology}{1.12}
\providecommand{\CthreePartialSeqVdpFailedFilters}{}\renewcommand{\CthreePartialSeqVdpFailedFilters}{0/8}
\providecommand{\CthreePartialSeqVdpVsDataDriven}{}\renewcommand{\CthreePartialSeqVdpVsDataDriven}{-30\%}
\providecommand{\CthreePartialSeqVdpVsWorking}{}\renewcommand{\CthreePartialSeqVdpVsWorking}{-78\%}
\providecommand{\CthreePartialSeqVsAny}{}\renewcommand{\CthreePartialSeqVsAny}{4/5}
\providecommand{\CthreePartialSeqVsDataDriven}{}\renewcommand{\CthreePartialSeqVsDataDriven}{1/5}
\providecommand{\CthreePartialSeqVsNominal}{}\renewcommand{\CthreePartialSeqVsNominal}{4/5}
\providecommand{\CthreePartialSeqVsRobust}{}\renewcommand{\CthreePartialSeqVsRobust}{4/5}
\providecommand{\CthreePartialSeqVsWorking}{}\renewcommand{\CthreePartialSeqVsWorking}{4/5}
\providecommand{\CthreeRealisticIdBatteryBestDataDriven}{}\renewcommand{\CthreeRealisticIdBatteryBestDataDriven}{0.0185}
\providecommand{\CthreeRealisticIdBatteryBestDataDrivenMethod}{}\renewcommand{\CthreeRealisticIdBatteryBestDataDrivenMethod}{greybox\_ekf\_oracle}
\providecommand{\CthreeRealisticIdBatteryBestLearned}{}\renewcommand{\CthreeRealisticIdBatteryBestLearned}{0.0506}
\providecommand{\CthreeRealisticIdBatteryBestWorking}{}\renewcommand{\CthreeRealisticIdBatteryBestWorking}{0.37}
\providecommand{\CthreeRealisticIdBatteryClimatology}{}\renewcommand{\CthreeRealisticIdBatteryClimatology}{0.735}
\providecommand{\CthreeRealisticIdBatteryFailedFilters}{}\renewcommand{\CthreeRealisticIdBatteryFailedFilters}{1/8}
\providecommand{\CthreeRealisticIdBatteryVsDataDriven}{}\renewcommand{\CthreeRealisticIdBatteryVsDataDriven}{+174\%}
\providecommand{\CthreeRealisticIdBatteryVsWorking}{}\renewcommand{\CthreeRealisticIdBatteryVsWorking}{-86\%}
\providecommand{\CthreeRealisticIdBioreactorBestDataDriven}{}\renewcommand{\CthreeRealisticIdBioreactorBestDataDriven}{0.0171}
\providecommand{\CthreeRealisticIdBioreactorBestDataDrivenMethod}{}\renewcommand{\CthreeRealisticIdBioreactorBestDataDrivenMethod}{nlox}
\providecommand{\CthreeRealisticIdBioreactorBestLearned}{}\renewcommand{\CthreeRealisticIdBioreactorBestLearned}{0.0313}
\providecommand{\CthreeRealisticIdBioreactorClimatology}{}\renewcommand{\CthreeRealisticIdBioreactorClimatology}{0.655}
\providecommand{\CthreeRealisticIdBioreactorFailedFilters}{}\renewcommand{\CthreeRealisticIdBioreactorFailedFilters}{8/8}
\providecommand{\CthreeRealisticIdBioreactorVsDataDriven}{}\renewcommand{\CthreeRealisticIdBioreactorVsDataDriven}{+85\%}
\providecommand{\CthreeRealisticIdLorenzBestDataDriven}{}\renewcommand{\CthreeRealisticIdLorenzBestDataDriven}{0.0134}
\providecommand{\CthreeRealisticIdLorenzBestDataDrivenMethod}{}\renewcommand{\CthreeRealisticIdLorenzBestDataDrivenMethod}{nlox}
\providecommand{\CthreeRealisticIdLorenzBestLearned}{}\renewcommand{\CthreeRealisticIdLorenzBestLearned}{0.011}
\providecommand{\CthreeRealisticIdLorenzBestWorking}{}\renewcommand{\CthreeRealisticIdLorenzBestWorking}{0.0273}
\providecommand{\CthreeRealisticIdLorenzClimatology}{}\renewcommand{\CthreeRealisticIdLorenzClimatology}{0.614}
\providecommand{\CthreeRealisticIdLorenzFailedFilters}{}\renewcommand{\CthreeRealisticIdLorenzFailedFilters}{7/8}
\providecommand{\CthreeRealisticIdLorenzVsDataDriven}{}\renewcommand{\CthreeRealisticIdLorenzVsDataDriven}{-18\%}
\providecommand{\CthreeRealisticIdLorenzVsWorking}{}\renewcommand{\CthreeRealisticIdLorenzVsWorking}{-58\%}
\providecommand{\CthreeRealisticIdQuadTankBestDataDriven}{}\renewcommand{\CthreeRealisticIdQuadTankBestDataDriven}{0.0168}
\providecommand{\CthreeRealisticIdQuadTankBestDataDrivenMethod}{}\renewcommand{\CthreeRealisticIdQuadTankBestDataDrivenMethod}{greybox\_ekf}
\providecommand{\CthreeRealisticIdQuadTankBestLearned}{}\renewcommand{\CthreeRealisticIdQuadTankBestLearned}{0.0245}
\providecommand{\CthreeRealisticIdQuadTankClimatology}{}\renewcommand{\CthreeRealisticIdQuadTankClimatology}{0.53}
\providecommand{\CthreeRealisticIdQuadTankFailedFilters}{}\renewcommand{\CthreeRealisticIdQuadTankFailedFilters}{8/8}
\providecommand{\CthreeRealisticIdQuadTankVsDataDriven}{}\renewcommand{\CthreeRealisticIdQuadTankVsDataDriven}{+46\%}
\providecommand{\CthreeRealisticIdVdpBestDataDriven}{}\renewcommand{\CthreeRealisticIdVdpBestDataDriven}{0.031}
\providecommand{\CthreeRealisticIdVdpBestDataDrivenMethod}{}\renewcommand{\CthreeRealisticIdVdpBestDataDrivenMethod}{greybox\_ekf}
\providecommand{\CthreeRealisticIdVdpBestLearned}{}\renewcommand{\CthreeRealisticIdVdpBestLearned}{0.0294}
\providecommand{\CthreeRealisticIdVdpBestWorking}{}\renewcommand{\CthreeRealisticIdVdpBestWorking}{0.108}
\providecommand{\CthreeRealisticIdVdpClimatology}{}\renewcommand{\CthreeRealisticIdVdpClimatology}{1.13}
\providecommand{\CthreeRealisticIdVdpFailedFilters}{}\renewcommand{\CthreeRealisticIdVdpFailedFilters}{0/8}
\providecommand{\CthreeRealisticIdVdpVsDataDriven}{}\renewcommand{\CthreeRealisticIdVdpVsDataDriven}{-5\%}
\providecommand{\CthreeRealisticIdVdpVsWorking}{}\renewcommand{\CthreeRealisticIdVdpVsWorking}{-73\%}
\providecommand{\CthreeRealisticIdVsAny}{}\renewcommand{\CthreeRealisticIdVsAny}{5/5}
\providecommand{\CthreeRealisticIdVsDataDriven}{}\renewcommand{\CthreeRealisticIdVsDataDriven}{1/5}
\providecommand{\CthreeRealisticIdVsNominal}{}\renewcommand{\CthreeRealisticIdVsNominal}{5/5}
\providecommand{\CthreeRealisticIdVsRobust}{}\renewcommand{\CthreeRealisticIdVsRobust}{5/5}
\providecommand{\CthreeRealisticIdVsWorking}{}\renewcommand{\CthreeRealisticIdVsWorking}{5/5}
\providecommand{\CthreeRealisticSeqBatteryBestDataDriven}{}\renewcommand{\CthreeRealisticSeqBatteryBestDataDriven}{0.0214}
\providecommand{\CthreeRealisticSeqBatteryBestDataDrivenMethod}{}\renewcommand{\CthreeRealisticSeqBatteryBestDataDrivenMethod}{greybox\_ekf\_oracle}
\providecommand{\CthreeRealisticSeqBatteryBestLearned}{}\renewcommand{\CthreeRealisticSeqBatteryBestLearned}{0.0642}
\providecommand{\CthreeRealisticSeqBatteryBestWorking}{}\renewcommand{\CthreeRealisticSeqBatteryBestWorking}{0.29}
\providecommand{\CthreeRealisticSeqBatteryClimatology}{}\renewcommand{\CthreeRealisticSeqBatteryClimatology}{0.765}
\providecommand{\CthreeRealisticSeqBatteryFailedFilters}{}\renewcommand{\CthreeRealisticSeqBatteryFailedFilters}{1/8}
\providecommand{\CthreeRealisticSeqBatteryVsDataDriven}{}\renewcommand{\CthreeRealisticSeqBatteryVsDataDriven}{+202\%}
\providecommand{\CthreeRealisticSeqBatteryVsWorking}{}\renewcommand{\CthreeRealisticSeqBatteryVsWorking}{-78\%}
\providecommand{\CthreeRealisticSeqBioreactorBestDataDriven}{}\renewcommand{\CthreeRealisticSeqBioreactorBestDataDriven}{0.0389}
\providecommand{\CthreeRealisticSeqBioreactorBestDataDrivenMethod}{}\renewcommand{\CthreeRealisticSeqBioreactorBestDataDrivenMethod}{nlox}
\providecommand{\CthreeRealisticSeqBioreactorBestLearned}{}\renewcommand{\CthreeRealisticSeqBioreactorBestLearned}{0.0987}
\providecommand{\CthreeRealisticSeqBioreactorClimatology}{}\renewcommand{\CthreeRealisticSeqBioreactorClimatology}{0.708}
\providecommand{\CthreeRealisticSeqBioreactorFailedFilters}{}\renewcommand{\CthreeRealisticSeqBioreactorFailedFilters}{8/8}
\providecommand{\CthreeRealisticSeqBioreactorVsDataDriven}{}\renewcommand{\CthreeRealisticSeqBioreactorVsDataDriven}{+157\%}
\providecommand{\CthreeRealisticSeqLorenzBestDataDriven}{}\renewcommand{\CthreeRealisticSeqLorenzBestDataDriven}{0.0129}
\providecommand{\CthreeRealisticSeqLorenzBestDataDrivenMethod}{}\renewcommand{\CthreeRealisticSeqLorenzBestDataDrivenMethod}{nlox}
\providecommand{\CthreeRealisticSeqLorenzBestLearned}{}\renewcommand{\CthreeRealisticSeqLorenzBestLearned}{0.0104}
\providecommand{\CthreeRealisticSeqLorenzBestWorking}{}\renewcommand{\CthreeRealisticSeqLorenzBestWorking}{0.0257}
\providecommand{\CthreeRealisticSeqLorenzClimatology}{}\renewcommand{\CthreeRealisticSeqLorenzClimatology}{0.614}
\providecommand{\CthreeRealisticSeqLorenzFailedFilters}{}\renewcommand{\CthreeRealisticSeqLorenzFailedFilters}{7/8}
\providecommand{\CthreeRealisticSeqLorenzVsDataDriven}{}\renewcommand{\CthreeRealisticSeqLorenzVsDataDriven}{-19\%}
\providecommand{\CthreeRealisticSeqLorenzVsWorking}{}\renewcommand{\CthreeRealisticSeqLorenzVsWorking}{-58\%}
\providecommand{\CthreeRealisticSeqQuadTankBestDataDriven}{}\renewcommand{\CthreeRealisticSeqQuadTankBestDataDriven}{0.017}
\providecommand{\CthreeRealisticSeqQuadTankBestDataDrivenMethod}{}\renewcommand{\CthreeRealisticSeqQuadTankBestDataDrivenMethod}{greybox\_ekf}
\providecommand{\CthreeRealisticSeqQuadTankBestLearned}{}\renewcommand{\CthreeRealisticSeqQuadTankBestLearned}{0.0441}
\providecommand{\CthreeRealisticSeqQuadTankClimatology}{}\renewcommand{\CthreeRealisticSeqQuadTankClimatology}{0.581}
\providecommand{\CthreeRealisticSeqQuadTankFailedFilters}{}\renewcommand{\CthreeRealisticSeqQuadTankFailedFilters}{8/8}
\providecommand{\CthreeRealisticSeqQuadTankVsDataDriven}{}\renewcommand{\CthreeRealisticSeqQuadTankVsDataDriven}{+160\%}
\providecommand{\CthreeRealisticSeqVdpBestDataDriven}{}\renewcommand{\CthreeRealisticSeqVdpBestDataDriven}{0.0388}
\providecommand{\CthreeRealisticSeqVdpBestDataDrivenMethod}{}\renewcommand{\CthreeRealisticSeqVdpBestDataDrivenMethod}{greybox\_ekf}
\providecommand{\CthreeRealisticSeqVdpBestLearned}{}\renewcommand{\CthreeRealisticSeqVdpBestLearned}{0.0324}
\providecommand{\CthreeRealisticSeqVdpBestWorking}{}\renewcommand{\CthreeRealisticSeqVdpBestWorking}{0.125}
\providecommand{\CthreeRealisticSeqVdpClimatology}{}\renewcommand{\CthreeRealisticSeqVdpClimatology}{1.16}
\providecommand{\CthreeRealisticSeqVdpFailedFilters}{}\renewcommand{\CthreeRealisticSeqVdpFailedFilters}{0/8}
\providecommand{\CthreeRealisticSeqVdpVsDataDriven}{}\renewcommand{\CthreeRealisticSeqVdpVsDataDriven}{-16\%}
\providecommand{\CthreeRealisticSeqVdpVsWorking}{}\renewcommand{\CthreeRealisticSeqVdpVsWorking}{-74\%}
\providecommand{\CthreeRealisticSeqVsAny}{}\renewcommand{\CthreeRealisticSeqVsAny}{5/5}
\providecommand{\CthreeRealisticSeqVsDataDriven}{}\renewcommand{\CthreeRealisticSeqVsDataDriven}{2/5}
\providecommand{\CthreeRealisticSeqVsNominal}{}\renewcommand{\CthreeRealisticSeqVsNominal}{5/5}
\providecommand{\CthreeRealisticSeqVsRobust}{}\renewcommand{\CthreeRealisticSeqVsRobust}{5/5}
\providecommand{\CthreeRealisticSeqVsWorking}{}\renewcommand{\CthreeRealisticSeqVsWorking}{5/5}
\providecommand{\CthreeRegimesBurstOutageGreyboxWins}{}\renewcommand{\CthreeRegimesBurstOutageGreyboxWins}{3/5}
\providecommand{\CthreeRegimesBurstOutageKklWins}{}\renewcommand{\CthreeRegimesBurstOutageKklWins}{1/5}  % best KKL beats best grey-box, seq OOD
\providecommand{\CthreeRegimesBurstOutageRegime}{}\renewcommand{\CthreeRegimesBurstOutageRegime}{FAIL}
\providecommand{\CthreeRegimesCleanGreyboxWins}{}\renewcommand{\CthreeRegimesCleanGreyboxWins}{2/5}
\providecommand{\CthreeRegimesCleanKklWins}{}\renewcommand{\CthreeRegimesCleanKklWins}{3/5}  % best KKL beats best grey-box, seq OOD
\providecommand{\CthreeRegimesCleanRegime}{}\renewcommand{\CthreeRegimesCleanRegime}{FAIL}
\providecommand{\CthreeRegimesNoiseTenGreyboxWins}{}\renewcommand{\CthreeRegimesNoiseTenGreyboxWins}{3/5}
\providecommand{\CthreeRegimesNoiseTenKklWins}{}\renewcommand{\CthreeRegimesNoiseTenKklWins}{2/5}  % best KKL beats best grey-box, seq OOD
\providecommand{\CthreeRegimesNoiseTenRegime}{}\renewcommand{\CthreeRegimesNoiseTenRegime}{FAIL}
\providecommand{\CthreeRegimesOutliersGreyboxWins}{}\renewcommand{\CthreeRegimesOutliersGreyboxWins}{3/5}
\providecommand{\CthreeRegimesOutliersKklWins}{}\renewcommand{\CthreeRegimesOutliersKklWins}{1/5}  % best KKL beats best grey-box, seq OOD
\providecommand{\CthreeRegimesOutliersRegime}{}\renewcommand{\CthreeRegimesOutliersRegime}{FAIL}
\providecommand{\CthreeRegimesPlantDriftTenGreyboxWins}{}\renewcommand{\CthreeRegimesPlantDriftTenGreyboxWins}{2/5}
\providecommand{\CthreeRegimesPlantDriftTenKklWins}{}\renewcommand{\CthreeRegimesPlantDriftTenKklWins}{3/5}  % best KKL beats best grey-box, seq OOD
\providecommand{\CthreeRegimesPlantDriftTenRegime}{}\renewcommand{\CthreeRegimesPlantDriftTenRegime}{FAIL}
\providecommand{\CthreeRegimesPoorInitGreyboxWins}{}\renewcommand{\CthreeRegimesPoorInitGreyboxWins}{2/5}
\providecommand{\CthreeRegimesPoorInitKklWins}{}\renewcommand{\CthreeRegimesPoorInitKklWins}{3/5}  % best KKL beats best grey-box, seq OOD
\providecommand{\CthreeRegimesPoorInitRegime}{}\renewcommand{\CthreeRegimesPoorInitRegime}{FAIL}
\providecommand{\CthreeTESTCThreeIdCartpoleBestDataDriven}{}\renewcommand{\CthreeTESTCThreeIdCartpoleBestDataDriven}{0.0406}
\providecommand{\CthreeTESTCThreeIdCartpoleBestDataDrivenMethod}{}\renewcommand{\CthreeTESTCThreeIdCartpoleBestDataDrivenMethod}{greybox\_ukf}
\providecommand{\CthreeTESTCThreeIdCartpoleBestLearned}{}\renewcommand{\CthreeTESTCThreeIdCartpoleBestLearned}{0.326}
\providecommand{\CthreeTESTCThreeIdCartpoleBestWorking}{}\renewcommand{\CthreeTESTCThreeIdCartpoleBestWorking}{0.154}
\providecommand{\CthreeTESTCThreeIdCartpoleClimatology}{}\renewcommand{\CthreeTESTCThreeIdCartpoleClimatology}{0.966}
\providecommand{\CthreeTESTCThreeIdCartpoleFailedFilters}{}\renewcommand{\CthreeTESTCThreeIdCartpoleFailedFilters}{1/8}
\providecommand{\CthreeTESTCThreeIdCartpoleVsDataDriven}{}\renewcommand{\CthreeTESTCThreeIdCartpoleVsDataDriven}{+697\%}
\providecommand{\CthreeTESTCThreeIdCartpoleVsWorking}{}\renewcommand{\CthreeTESTCThreeIdCartpoleVsWorking}{+111\%}
\providecommand{\CthreeTESTCThreeIdFhnBestDataDriven}{}\renewcommand{\CthreeTESTCThreeIdFhnBestDataDriven}{0.015}
\providecommand{\CthreeTESTCThreeIdFhnBestDataDrivenMethod}{}\renewcommand{\CthreeTESTCThreeIdFhnBestDataDrivenMethod}{greybox\_ekf}
\providecommand{\CthreeTESTCThreeIdFhnBestLearned}{}\renewcommand{\CthreeTESTCThreeIdFhnBestLearned}{0.0421}
\providecommand{\CthreeTESTCThreeIdFhnBestWorking}{}\renewcommand{\CthreeTESTCThreeIdFhnBestWorking}{0.132}
\providecommand{\CthreeTESTCThreeIdFhnClimatology}{}\renewcommand{\CthreeTESTCThreeIdFhnClimatology}{0.889}
\providecommand{\CthreeTESTCThreeIdFhnFailedFilters}{}\renewcommand{\CthreeTESTCThreeIdFhnFailedFilters}{1/8}
\providecommand{\CthreeTESTCThreeIdFhnVsDataDriven}{}\renewcommand{\CthreeTESTCThreeIdFhnVsDataDriven}{+180\%}
\providecommand{\CthreeTESTCThreeIdFhnVsWorking}{}\renewcommand{\CthreeTESTCThreeIdFhnVsWorking}{-68\%}
\providecommand{\CthreeTESTCThreeIdHighwayTrafficBestDataDriven}{}\renewcommand{\CthreeTESTCThreeIdHighwayTrafficBestDataDriven}{0.036}
\providecommand{\CthreeTESTCThreeIdHighwayTrafficBestDataDrivenMethod}{}\renewcommand{\CthreeTESTCThreeIdHighwayTrafficBestDataDrivenMethod}{greybox\_ekf}
\providecommand{\CthreeTESTCThreeIdHighwayTrafficBestLearned}{}\renewcommand{\CthreeTESTCThreeIdHighwayTrafficBestLearned}{0.127}
\providecommand{\CthreeTESTCThreeIdHighwayTrafficBestWorking}{}\renewcommand{\CthreeTESTCThreeIdHighwayTrafficBestWorking}{0.144}
\providecommand{\CthreeTESTCThreeIdHighwayTrafficClimatology}{}\renewcommand{\CthreeTESTCThreeIdHighwayTrafficClimatology}{0.725}
\providecommand{\CthreeTESTCThreeIdHighwayTrafficFailedFilters}{}\renewcommand{\CthreeTESTCThreeIdHighwayTrafficFailedFilters}{1/8}
\providecommand{\CthreeTESTCThreeIdHighwayTrafficVsDataDriven}{}\renewcommand{\CthreeTESTCThreeIdHighwayTrafficVsDataDriven}{+261\%}
\providecommand{\CthreeTESTCThreeIdHighwayTrafficVsWorking}{}\renewcommand{\CthreeTESTCThreeIdHighwayTrafficVsWorking}{-12\%}
\providecommand{\CthreeTESTCThreeIdNetworkBestLearned}{}\renewcommand{\CthreeTESTCThreeIdNetworkBestLearned}{0.0507}
\providecommand{\CthreeTESTCThreeIdNetworkBestWorking}{}\renewcommand{\CthreeTESTCThreeIdNetworkBestWorking}{0.226}
\providecommand{\CthreeTESTCThreeIdNetworkClimatology}{}\renewcommand{\CthreeTESTCThreeIdNetworkClimatology}{0.993}
\providecommand{\CthreeTESTCThreeIdNetworkFailedFilters}{}\renewcommand{\CthreeTESTCThreeIdNetworkFailedFilters}{0/8}
\providecommand{\CthreeTESTCThreeIdNetworkVsWorking}{}\renewcommand{\CthreeTESTCThreeIdNetworkVsWorking}{-78\%}
\providecommand{\CthreeTESTCThreeIdRosslerBestDataDriven}{}\renewcommand{\CthreeTESTCThreeIdRosslerBestDataDriven}{0.37}
\providecommand{\CthreeTESTCThreeIdRosslerBestDataDrivenMethod}{}\renewcommand{\CthreeTESTCThreeIdRosslerBestDataDrivenMethod}{nlox}
\providecommand{\CthreeTESTCThreeIdRosslerBestLearned}{}\renewcommand{\CthreeTESTCThreeIdRosslerBestLearned}{0.246}
\providecommand{\CthreeTESTCThreeIdRosslerBestWorking}{}\renewcommand{\CthreeTESTCThreeIdRosslerBestWorking}{1.11}
\providecommand{\CthreeTESTCThreeIdRosslerClimatology}{}\renewcommand{\CthreeTESTCThreeIdRosslerClimatology}{1.65}
\providecommand{\CthreeTESTCThreeIdRosslerFailedFilters}{}\renewcommand{\CthreeTESTCThreeIdRosslerFailedFilters}{0/8}
\providecommand{\CthreeTESTCThreeIdRosslerVsDataDriven}{}\renewcommand{\CthreeTESTCThreeIdRosslerVsDataDriven}{-33\%}
\providecommand{\CthreeTESTCThreeIdRosslerVsWorking}{}\renewcommand{\CthreeTESTCThreeIdRosslerVsWorking}{-78\%}
\providecommand{\CthreeTESTCThreeIdTwoLinkArmBestLearned}{}\renewcommand{\CthreeTESTCThreeIdTwoLinkArmBestLearned}{0.279}
\providecommand{\CthreeTESTCThreeIdTwoLinkArmBestWorking}{}\renewcommand{\CthreeTESTCThreeIdTwoLinkArmBestWorking}{0.28}
\providecommand{\CthreeTESTCThreeIdTwoLinkArmClimatology}{}\renewcommand{\CthreeTESTCThreeIdTwoLinkArmClimatology}{0.676}
\providecommand{\CthreeTESTCThreeIdTwoLinkArmFailedFilters}{}\renewcommand{\CthreeTESTCThreeIdTwoLinkArmFailedFilters}{2/8}
\providecommand{\CthreeTESTCThreeIdTwoLinkArmVsWorking}{}\renewcommand{\CthreeTESTCThreeIdTwoLinkArmVsWorking}{+0\%}
\providecommand{\CthreeTESTCThreeIdVsAny}{}\renewcommand{\CthreeTESTCThreeIdVsAny}{3/6}
\providecommand{\CthreeTESTCThreeIdVsDataDriven}{}\renewcommand{\CthreeTESTCThreeIdVsDataDriven}{1/4}
\providecommand{\CthreeTESTCThreeIdVsNominal}{}\renewcommand{\CthreeTESTCThreeIdVsNominal}{4/6}
\providecommand{\CthreeTESTCThreeIdVsRobust}{}\renewcommand{\CthreeTESTCThreeIdVsRobust}{3/6}
\providecommand{\CthreeTESTCThreeIdVsWorking}{}\renewcommand{\CthreeTESTCThreeIdVsWorking}{3/6}
\providecommand{\CthreeTESTCThreeSeqCartpoleBestDataDriven}{}\renewcommand{\CthreeTESTCThreeSeqCartpoleBestDataDriven}{0.064}
\providecommand{\CthreeTESTCThreeSeqCartpoleBestDataDrivenMethod}{}\renewcommand{\CthreeTESTCThreeSeqCartpoleBestDataDrivenMethod}{greybox\_ukf}
\providecommand{\CthreeTESTCThreeSeqCartpoleBestLearned}{}\renewcommand{\CthreeTESTCThreeSeqCartpoleBestLearned}{0.597}
\providecommand{\CthreeTESTCThreeSeqCartpoleBestWorking}{}\renewcommand{\CthreeTESTCThreeSeqCartpoleBestWorking}{0.205}
\providecommand{\CthreeTESTCThreeSeqCartpoleClimatology}{}\renewcommand{\CthreeTESTCThreeSeqCartpoleClimatology}{1.3}
\providecommand{\CthreeTESTCThreeSeqCartpoleFailedFilters}{}\renewcommand{\CthreeTESTCThreeSeqCartpoleFailedFilters}{1/8}
\providecommand{\CthreeTESTCThreeSeqCartpoleVsDataDriven}{}\renewcommand{\CthreeTESTCThreeSeqCartpoleVsDataDriven}{+826\%}
\providecommand{\CthreeTESTCThreeSeqCartpoleVsWorking}{}\renewcommand{\CthreeTESTCThreeSeqCartpoleVsWorking}{+192\%}
\providecommand{\CthreeTESTCThreeSeqFhnBestDataDriven}{}\renewcommand{\CthreeTESTCThreeSeqFhnBestDataDriven}{0.0152}
\providecommand{\CthreeTESTCThreeSeqFhnBestDataDrivenMethod}{}\renewcommand{\CthreeTESTCThreeSeqFhnBestDataDrivenMethod}{greybox\_ekf}
\providecommand{\CthreeTESTCThreeSeqFhnBestLearned}{}\renewcommand{\CthreeTESTCThreeSeqFhnBestLearned}{0.0678}
\providecommand{\CthreeTESTCThreeSeqFhnBestWorking}{}\renewcommand{\CthreeTESTCThreeSeqFhnBestWorking}{0.0971}
\providecommand{\CthreeTESTCThreeSeqFhnClimatology}{}\renewcommand{\CthreeTESTCThreeSeqFhnClimatology}{0.863}
\providecommand{\CthreeTESTCThreeSeqFhnFailedFilters}{}\renewcommand{\CthreeTESTCThreeSeqFhnFailedFilters}{1/8}
\providecommand{\CthreeTESTCThreeSeqFhnVsDataDriven}{}\renewcommand{\CthreeTESTCThreeSeqFhnVsDataDriven}{+347\%}
\providecommand{\CthreeTESTCThreeSeqFhnVsWorking}{}\renewcommand{\CthreeTESTCThreeSeqFhnVsWorking}{-30\%}
\providecommand{\CthreeTESTCThreeSeqHighwayTrafficBestDataDriven}{}\renewcommand{\CthreeTESTCThreeSeqHighwayTrafficBestDataDriven}{0.0273}
\providecommand{\CthreeTESTCThreeSeqHighwayTrafficBestDataDrivenMethod}{}\renewcommand{\CthreeTESTCThreeSeqHighwayTrafficBestDataDrivenMethod}{greybox\_ekf}
\providecommand{\CthreeTESTCThreeSeqHighwayTrafficBestLearned}{}\renewcommand{\CthreeTESTCThreeSeqHighwayTrafficBestLearned}{0.13}
\providecommand{\CthreeTESTCThreeSeqHighwayTrafficBestWorking}{}\renewcommand{\CthreeTESTCThreeSeqHighwayTrafficBestWorking}{0.139}
\providecommand{\CthreeTESTCThreeSeqHighwayTrafficClimatology}{}\renewcommand{\CthreeTESTCThreeSeqHighwayTrafficClimatology}{0.744}
\providecommand{\CthreeTESTCThreeSeqHighwayTrafficFailedFilters}{}\renewcommand{\CthreeTESTCThreeSeqHighwayTrafficFailedFilters}{1/8}
\providecommand{\CthreeTESTCThreeSeqHighwayTrafficVsDataDriven}{}\renewcommand{\CthreeTESTCThreeSeqHighwayTrafficVsDataDriven}{+379\%}
\providecommand{\CthreeTESTCThreeSeqHighwayTrafficVsWorking}{}\renewcommand{\CthreeTESTCThreeSeqHighwayTrafficVsWorking}{-6\%}
\providecommand{\CthreeTESTCThreeSeqNetworkBestLearned}{}\renewcommand{\CthreeTESTCThreeSeqNetworkBestLearned}{0.0638}
\providecommand{\CthreeTESTCThreeSeqNetworkBestWorking}{}\renewcommand{\CthreeTESTCThreeSeqNetworkBestWorking}{0.226}
\providecommand{\CthreeTESTCThreeSeqNetworkClimatology}{}\renewcommand{\CthreeTESTCThreeSeqNetworkClimatology}{1}
\providecommand{\CthreeTESTCThreeSeqNetworkFailedFilters}{}\renewcommand{\CthreeTESTCThreeSeqNetworkFailedFilters}{0/8}
\providecommand{\CthreeTESTCThreeSeqNetworkVsWorking}{}\renewcommand{\CthreeTESTCThreeSeqNetworkVsWorking}{-72\%}
\providecommand{\CthreeTESTCThreeSeqRosslerBestDataDriven}{}\renewcommand{\CthreeTESTCThreeSeqRosslerBestDataDriven}{0.376}
\providecommand{\CthreeTESTCThreeSeqRosslerBestDataDrivenMethod}{}\renewcommand{\CthreeTESTCThreeSeqRosslerBestDataDrivenMethod}{nlox}
\providecommand{\CthreeTESTCThreeSeqRosslerBestLearned}{}\renewcommand{\CthreeTESTCThreeSeqRosslerBestLearned}{0.247}
\providecommand{\CthreeTESTCThreeSeqRosslerBestWorking}{}\renewcommand{\CthreeTESTCThreeSeqRosslerBestWorking}{1.14}
\providecommand{\CthreeTESTCThreeSeqRosslerClimatology}{}\renewcommand{\CthreeTESTCThreeSeqRosslerClimatology}{1.62}
\providecommand{\CthreeTESTCThreeSeqRosslerFailedFilters}{}\renewcommand{\CthreeTESTCThreeSeqRosslerFailedFilters}{0/8}
\providecommand{\CthreeTESTCThreeSeqRosslerVsDataDriven}{}\renewcommand{\CthreeTESTCThreeSeqRosslerVsDataDriven}{-34\%}
\providecommand{\CthreeTESTCThreeSeqRosslerVsWorking}{}\renewcommand{\CthreeTESTCThreeSeqRosslerVsWorking}{-78\%}
\providecommand{\CthreeTESTCThreeSeqTwoLinkArmBestLearned}{}\renewcommand{\CthreeTESTCThreeSeqTwoLinkArmBestLearned}{0.271}
\providecommand{\CthreeTESTCThreeSeqTwoLinkArmBestWorking}{}\renewcommand{\CthreeTESTCThreeSeqTwoLinkArmBestWorking}{0.287}
\providecommand{\CthreeTESTCThreeSeqTwoLinkArmClimatology}{}\renewcommand{\CthreeTESTCThreeSeqTwoLinkArmClimatology}{0.597}
\providecommand{\CthreeTESTCThreeSeqTwoLinkArmFailedFilters}{}\renewcommand{\CthreeTESTCThreeSeqTwoLinkArmFailedFilters}{2/8}
\providecommand{\CthreeTESTCThreeSeqTwoLinkArmVsWorking}{}\renewcommand{\CthreeTESTCThreeSeqTwoLinkArmVsWorking}{-5\%}
\providecommand{\CthreeTESTCThreeSeqVsAny}{}\renewcommand{\CthreeTESTCThreeSeqVsAny}{3/6}
\providecommand{\CthreeTESTCThreeSeqVsDataDriven}{}\renewcommand{\CthreeTESTCThreeSeqVsDataDriven}{1/4}
\providecommand{\CthreeTESTCThreeSeqVsNominal}{}\renewcommand{\CthreeTESTCThreeSeqVsNominal}{4/6}
\providecommand{\CthreeTESTCThreeSeqVsRobust}{}\renewcommand{\CthreeTESTCThreeSeqVsRobust}{3/6}
\providecommand{\CthreeTESTCThreeSeqVsWorking}{}\renewcommand{\CthreeTESTCThreeSeqVsWorking}{3/6}
\providecommand{\CtwoFloorFiveBatteryHyperKKLId}{}\renewcommand{\CtwoFloorFiveBatteryHyperKKLId}{+33\%}
\providecommand{\CtwoFloorFiveBatteryHyperKKLSeq}{}\renewcommand{\CtwoFloorFiveBatteryHyperKKLSeq}{+36\%}
\providecommand{\CtwoFloorFiveBioreactorHyperKKLId}{}\renewcommand{\CtwoFloorFiveBioreactorHyperKKLId}{-8\%}
\providecommand{\CtwoFloorFiveBioreactorHyperKKLSeq}{}\renewcommand{\CtwoFloorFiveBioreactorHyperKKLSeq}{+13\%}
\providecommand{\CtwoFloorFiveLorenzHyperKKLId}{}\renewcommand{\CtwoFloorFiveLorenzHyperKKLId}{+5\%}
\providecommand{\CtwoFloorFiveLorenzHyperKKLSeq}{}\renewcommand{\CtwoFloorFiveLorenzHyperKKLSeq}{+7\%}
\providecommand{\CtwoFloorFiveNoLoss}{}\renewcommand{\CtwoFloorFiveNoLoss}{4/5}
\providecommand{\CtwoFloorFiveProbeAboveFloor}{}\renewcommand{\CtwoFloorFiveProbeAboveFloor}{0/5}
\providecommand{\CtwoFloorFiveProbeRatioBattery}{}\renewcommand{\CtwoFloorFiveProbeRatioBattery}{0.77}
\providecommand{\CtwoFloorFiveProbeRatioBioreactor}{}\renewcommand{\CtwoFloorFiveProbeRatioBioreactor}{0.69}
\providecommand{\CtwoFloorFiveProbeRatioLorenz}{}\renewcommand{\CtwoFloorFiveProbeRatioLorenz}{0.01}
\providecommand{\CtwoFloorFiveProbeRatioQuadTank}{}\renewcommand{\CtwoFloorFiveProbeRatioQuadTank}{0.87}
\providecommand{\CtwoFloorFiveProbeRatioVdp}{}\renewcommand{\CtwoFloorFiveProbeRatioVdp}{0.81}
\providecommand{\CtwoFloorFiveQuadTankHyperKKLId}{}\renewcommand{\CtwoFloorFiveQuadTankHyperKKLId}{+2\%}
\providecommand{\CtwoFloorFiveQuadTankHyperKKLSeq}{}\renewcommand{\CtwoFloorFiveQuadTankHyperKKLSeq}{-8\%}
\providecommand{\CtwoFloorFiveTostNonInferior}{}\renewcommand{\CtwoFloorFiveTostNonInferior}{1/5}
\providecommand{\CtwoFloorFiveVdpHyperKKLId}{}\renewcommand{\CtwoFloorFiveVdpHyperKKLId}{-1\%}
\providecommand{\CtwoFloorFiveVdpHyperKKLSeq}{}\renewcommand{\CtwoFloorFiveVdpHyperKKLSeq}{-3\%}
\providecommand{\CtwoFloorOneBatteryHyperKKLId}{}\renewcommand{\CtwoFloorOneBatteryHyperKKLId}{+106\%}
\providecommand{\CtwoFloorOneBatteryHyperKKLSeq}{}\renewcommand{\CtwoFloorOneBatteryHyperKKLSeq}{+104\%}
\providecommand{\CtwoFloorOneBioreactorHyperKKLId}{}\renewcommand{\CtwoFloorOneBioreactorHyperKKLId}{-5\%}
\providecommand{\CtwoFloorOneBioreactorHyperKKLSeq}{}\renewcommand{\CtwoFloorOneBioreactorHyperKKLSeq}{-33\%}
\providecommand{\CtwoFloorOneLorenzHyperKKLId}{}\renewcommand{\CtwoFloorOneLorenzHyperKKLId}{+16\%}
\providecommand{\CtwoFloorOneLorenzHyperKKLSeq}{}\renewcommand{\CtwoFloorOneLorenzHyperKKLSeq}{-7\%}
\providecommand{\CtwoFloorOneNoLoss}{}\renewcommand{\CtwoFloorOneNoLoss}{3/5}
\providecommand{\CtwoFloorOneProbeAboveFloor}{}\renewcommand{\CtwoFloorOneProbeAboveFloor}{0/5}
\providecommand{\CtwoFloorOneProbeRatioBattery}{}\renewcommand{\CtwoFloorOneProbeRatioBattery}{0.92}
\providecommand{\CtwoFloorOneProbeRatioBioreactor}{}\renewcommand{\CtwoFloorOneProbeRatioBioreactor}{0.92}
\providecommand{\CtwoFloorOneProbeRatioLorenz}{}\renewcommand{\CtwoFloorOneProbeRatioLorenz}{0.76}
\providecommand{\CtwoFloorOneProbeRatioQuadTank}{}\renewcommand{\CtwoFloorOneProbeRatioQuadTank}{0.48}
\providecommand{\CtwoFloorOneProbeRatioVdp}{}\renewcommand{\CtwoFloorOneProbeRatioVdp}{0.91}
\providecommand{\CtwoFloorOneQuadTankHyperKKLId}{}\renewcommand{\CtwoFloorOneQuadTankHyperKKLId}{+10\%}
\providecommand{\CtwoFloorOneQuadTankHyperKKLSeq}{}\renewcommand{\CtwoFloorOneQuadTankHyperKKLSeq}{-1\%}
\providecommand{\CtwoFloorOneTostNonInferior}{}\renewcommand{\CtwoFloorOneTostNonInferior}{0/5}
\providecommand{\CtwoFloorOneVdpHyperKKLId}{}\renewcommand{\CtwoFloorOneVdpHyperKKLId}{-1\%}
\providecommand{\CtwoFloorOneVdpHyperKKLSeq}{}\renewcommand{\CtwoFloorOneVdpHyperKKLSeq}{-5\%}
\providecommand{\CtwoFloorTwentyFiveBatteryHyperKKLId}{}\renewcommand{\CtwoFloorTwentyFiveBatteryHyperKKLId}{-1\%}
\providecommand{\CtwoFloorTwentyFiveBatteryHyperKKLSeq}{}\renewcommand{\CtwoFloorTwentyFiveBatteryHyperKKLSeq}{-6\%}
\providecommand{\CtwoFloorTwentyFiveBioreactorHyperKKLId}{}\renewcommand{\CtwoFloorTwentyFiveBioreactorHyperKKLId}{+17\%}
\providecommand{\CtwoFloorTwentyFiveBioreactorHyperKKLSeq}{}\renewcommand{\CtwoFloorTwentyFiveBioreactorHyperKKLSeq}{+24\%}
\providecommand{\CtwoFloorTwentyFiveLorenzHyperKKLId}{}\renewcommand{\CtwoFloorTwentyFiveLorenzHyperKKLId}{+3\%}
\providecommand{\CtwoFloorTwentyFiveLorenzHyperKKLSeq}{}\renewcommand{\CtwoFloorTwentyFiveLorenzHyperKKLSeq}{+6\%}
\providecommand{\CtwoFloorTwentyFiveNoLoss}{}\renewcommand{\CtwoFloorTwentyFiveNoLoss}{5/5}
\providecommand{\CtwoFloorTwentyFiveProbeAboveFloor}{}\renewcommand{\CtwoFloorTwentyFiveProbeAboveFloor}{1/5}
\providecommand{\CtwoFloorTwentyFiveProbeRatioBattery}{}\renewcommand{\CtwoFloorTwentyFiveProbeRatioBattery}{0.85}
\providecommand{\CtwoFloorTwentyFiveProbeRatioBioreactor}{}\renewcommand{\CtwoFloorTwentyFiveProbeRatioBioreactor}{0.38}
\providecommand{\CtwoFloorTwentyFiveProbeRatioLorenz}{}\renewcommand{\CtwoFloorTwentyFiveProbeRatioLorenz}{0.10}
\providecommand{\CtwoFloorTwentyFiveProbeRatioQuadTank}{}\renewcommand{\CtwoFloorTwentyFiveProbeRatioQuadTank}{1.82}
\providecommand{\CtwoFloorTwentyFiveProbeRatioVdp}{}\renewcommand{\CtwoFloorTwentyFiveProbeRatioVdp}{0.79}
\providecommand{\CtwoFloorTwentyFiveQuadTankHyperKKLId}{}\renewcommand{\CtwoFloorTwentyFiveQuadTankHyperKKLId}{+2\%}
\providecommand{\CtwoFloorTwentyFiveQuadTankHyperKKLSeq}{}\renewcommand{\CtwoFloorTwentyFiveQuadTankHyperKKLSeq}{-7\%}
\providecommand{\CtwoFloorTwentyFiveTostNonInferior}{}\renewcommand{\CtwoFloorTwentyFiveTostNonInferior}{3/5}
\providecommand{\CtwoFloorTwentyFiveVdpHyperKKLId}{}\renewcommand{\CtwoFloorTwentyFiveVdpHyperKKLId}{-4\%}
\providecommand{\CtwoFloorTwentyFiveVdpHyperKKLSeq}{}\renewcommand{\CtwoFloorTwentyFiveVdpHyperKKLSeq}{-5\%}
\providecommand{\CtwoKneeBatteryId}{}\renewcommand{\CtwoKneeBatteryId}{None}
\providecommand{\CtwoKneeBatterySeq}{}\renewcommand{\CtwoKneeBatterySeq}{None}
\providecommand{\CtwoKneeBioreactorId}{}\renewcommand{\CtwoKneeBioreactorId}{64}
\providecommand{\CtwoKneeBioreactorSeq}{}\renewcommand{\CtwoKneeBioreactorSeq}{32}
\providecommand{\CtwoKneeLorenzId}{}\renewcommand{\CtwoKneeLorenzId}{None}
\providecommand{\CtwoKneeLorenzSeq}{}\renewcommand{\CtwoKneeLorenzSeq}{None}
\providecommand{\CtwoKneeQuadTankId}{}\renewcommand{\CtwoKneeQuadTankId}{128}
\providecommand{\CtwoKneeQuadTankSeq}{}\renewcommand{\CtwoKneeQuadTankSeq}{None}
\providecommand{\CtwoKneeVdpId}{}\renewcommand{\CtwoKneeVdpId}{32}
\providecommand{\CtwoKneeVdpSeq}{}\renewcommand{\CtwoKneeVdpSeq}{32}
\providecommand{\CtwoLdecQBattery}{}\renewcommand{\CtwoLdecQBattery}{806}
\providecommand{\CtwoLdecQBioreactor}{}\renewcommand{\CtwoLdecQBioreactor}{44}
\providecommand{\CtwoLdecQLorenz}{}\renewcommand{\CtwoLdecQLorenz}{10}
\providecommand{\CtwoLdecQQuadTank}{}\renewcommand{\CtwoLdecQQuadTank}{14}
\providecommand{\CtwoLdecQVdp}{}\renewcommand{\CtwoLdecQVdp}{23}
\providecommand{\CtwoLrefinedBattery}{}\renewcommand{\CtwoLrefinedBattery}{435}
\providecommand{\CtwoLrefinedBioreactor}{}\renewcommand{\CtwoLrefinedBioreactor}{219}
\providecommand{\CtwoLrefinedLorenz}{}\renewcommand{\CtwoLrefinedLorenz}{239}
\providecommand{\CtwoLrefinedQuadTank}{}\renewcommand{\CtwoLrefinedQuadTank}{1121}
\providecommand{\CtwoLrefinedVdp}{}\renewcommand{\CtwoLrefinedVdp}{199}
\providecommand{\CtwoLstarBattery}{}\renewcommand{\CtwoLstarBattery}{135}
\providecommand{\CtwoLstarBioreactor}{}\renewcommand{\CtwoLstarBioreactor}{97}
\providecommand{\CtwoLstarLorenz}{}\renewcommand{\CtwoLstarLorenz}{135}
\providecommand{\CtwoLstarQuadTank}{}\renewcommand{\CtwoLstarQuadTank}{599}
\providecommand{\CtwoLstarVdp}{}\renewcommand{\CtwoLstarVdp}{97}
\providecommand{\CtwoProbeNoFloorRateOverAlphaBatteryHyperKKL}{}\renewcommand{\CtwoProbeNoFloorRateOverAlphaBatteryHyperKKL}{0.18}
\providecommand{\CtwoProbeNoFloorRateOverAlphaBioreactorHyperKKL}{}\renewcommand{\CtwoProbeNoFloorRateOverAlphaBioreactorHyperKKL}{0.07}
\providecommand{\CtwoProbeNoFloorRateOverAlphaLorenzHyperKKL}{}\renewcommand{\CtwoProbeNoFloorRateOverAlphaLorenzHyperKKL}{0.02}
\providecommand{\CtwoProbeNoFloorRateOverAlphaQuadTankHyperKKL}{}\renewcommand{\CtwoProbeNoFloorRateOverAlphaQuadTankHyperKKL}{0.47}
\providecommand{\CtwoProbeNoFloorRateOverAlphaVdpHyperKKL}{}\renewcommand{\CtwoProbeNoFloorRateOverAlphaVdpHyperKKL}{0.15}
\providecommand{\CtwoSfmPoleOverAlphaBatteryHyperKKL}{}\renewcommand{\CtwoSfmPoleOverAlphaBatteryHyperKKL}{0.18}
\providecommand{\CtwoSfmPoleOverAlphaBioreactorHyperKKL}{}\renewcommand{\CtwoSfmPoleOverAlphaBioreactorHyperKKL}{0.14}
\providecommand{\CtwoSfmPoleOverAlphaLorenzHyperKKL}{}\renewcommand{\CtwoSfmPoleOverAlphaLorenzHyperKKL}{0.18}
\providecommand{\CtwoSfmPoleOverAlphaQuadTankHyperKKL}{}\renewcommand{\CtwoSfmPoleOverAlphaQuadTankHyperKKL}{0.17}
\providecommand{\CtwoSfmPoleOverAlphaVdpHyperKKL}{}\renewcommand{\CtwoSfmPoleOverAlphaVdpHyperKKL}{0.15}
\providecommand{\CtwoWindowRefinedId}{}\renewcommand{\CtwoWindowRefinedId}{2/3}
\providecommand{\CtwoWindowRefinedSeq}{}\renewcommand{\CtwoWindowRefinedSeq}{2/3}
\providecommand{\CtwoWindowTheoryId}{}\renewcommand{\CtwoWindowTheoryId}{2/4}
\providecommand{\CtwoWindowTheorySeq}{}\renewcommand{\CtwoWindowTheorySeq}{2/4}
\providecommand{\CtwoWindowTwoFiveSixVsMaxBatteryId}{}\renewcommand{\CtwoWindowTwoFiveSixVsMaxBatteryId}{+63\%}
\providecommand{\CtwoWindowTwoFiveSixVsMaxBatterySeq}{}\renewcommand{\CtwoWindowTwoFiveSixVsMaxBatterySeq}{+73\%}
\providecommand{\CtwoWindowTwoFiveSixVsMaxBioreactorId}{}\renewcommand{\CtwoWindowTwoFiveSixVsMaxBioreactorId}{-18\%}
\providecommand{\CtwoWindowTwoFiveSixVsMaxBioreactorSeq}{}\renewcommand{\CtwoWindowTwoFiveSixVsMaxBioreactorSeq}{-42\%}
\providecommand{\CtwoWindowTwoFiveSixVsMaxLorenzId}{}\renewcommand{\CtwoWindowTwoFiveSixVsMaxLorenzId}{+34\%}
\providecommand{\CtwoWindowTwoFiveSixVsMaxLorenzSeq}{}\renewcommand{\CtwoWindowTwoFiveSixVsMaxLorenzSeq}{+12\%}
\providecommand{\CtwoWindowTwoFiveSixVsMaxQuadTankId}{}\renewcommand{\CtwoWindowTwoFiveSixVsMaxQuadTankId}{-0\%}
\providecommand{\CtwoWindowTwoFiveSixVsMaxQuadTankSeq}{}\renewcommand{\CtwoWindowTwoFiveSixVsMaxQuadTankSeq}{+8\%}
\providecommand{\CtwoWindowTwoFiveSixVsMaxVdpId}{}\renewcommand{\CtwoWindowTwoFiveSixVsMaxVdpId}{-1\%}
\providecommand{\CtwoWindowTwoFiveSixVsMaxVdpSeq}{}\renewcommand{\CtwoWindowTwoFiveSixVsMaxVdpSeq}{-5\%}
\providecommand{\HyperKKLRankOneVsFourBatterySeq}{}\renewcommand{\HyperKKLRankOneVsFourBatterySeq}{+29\%}
\providecommand{\HyperKKLRankOneVsFourBatteryShape}{}\renewcommand{\HyperKKLRankOneVsFourBatteryShape}{+33\%}
\providecommand{\HyperKKLRankOneVsFourBioreactorSeq}{}\renewcommand{\HyperKKLRankOneVsFourBioreactorSeq}{+1\%}
\providecommand{\HyperKKLRankOneVsFourBioreactorShape}{}\renewcommand{\HyperKKLRankOneVsFourBioreactorShape}{+8\%}
\providecommand{\HyperKKLRankOneVsFourLorenzSeq}{}\renewcommand{\HyperKKLRankOneVsFourLorenzSeq}{+58\%}
\providecommand{\HyperKKLRankOneVsFourLorenzShape}{}\renewcommand{\HyperKKLRankOneVsFourLorenzShape}{+58\%}
\providecommand{\HyperKKLRankOneVsFourQuadTankSeq}{}\renewcommand{\HyperKKLRankOneVsFourQuadTankSeq}{+3\%}
\providecommand{\HyperKKLRankOneVsFourQuadTankShape}{}\renewcommand{\HyperKKLRankOneVsFourQuadTankShape}{+11\%}
\providecommand{\HyperKKLRankOneVsFourVdpSeq}{}\renewcommand{\HyperKKLRankOneVsFourVdpSeq}{+23\%}
\providecommand{\HyperKKLRankOneVsFourVdpShape}{}\renewcommand{\HyperKKLRankOneVsFourVdpShape}{+25\%}
\providecommand{\HyperKKLRankSixteenVsFourBatterySeq}{}\renewcommand{\HyperKKLRankSixteenVsFourBatterySeq}{+12\%}
\providecommand{\HyperKKLRankSixteenVsFourBatteryShape}{}\renewcommand{\HyperKKLRankSixteenVsFourBatteryShape}{+3\%}
\providecommand{\HyperKKLRankSixteenVsFourBioreactorSeq}{}\renewcommand{\HyperKKLRankSixteenVsFourBioreactorSeq}{+10\%}
\providecommand{\HyperKKLRankSixteenVsFourBioreactorShape}{}\renewcommand{\HyperKKLRankSixteenVsFourBioreactorShape}{+6\%}
\providecommand{\HyperKKLRankSixteenVsFourLorenzSeq}{}\renewcommand{\HyperKKLRankSixteenVsFourLorenzSeq}{-4\%}
\providecommand{\HyperKKLRankSixteenVsFourLorenzShape}{}\renewcommand{\HyperKKLRankSixteenVsFourLorenzShape}{-10\%}
\providecommand{\HyperKKLRankSixteenVsFourQuadTankSeq}{}\renewcommand{\HyperKKLRankSixteenVsFourQuadTankSeq}{+6\%}
\providecommand{\HyperKKLRankSixteenVsFourQuadTankShape}{}\renewcommand{\HyperKKLRankSixteenVsFourQuadTankShape}{+11\%}
\providecommand{\HyperKKLRankSixteenVsFourVdpSeq}{}\renewcommand{\HyperKKLRankSixteenVsFourVdpSeq}{+2\%}
\providecommand{\HyperKKLRankSixteenVsFourVdpShape}{}\renewcommand{\HyperKKLRankSixteenVsFourVdpShape}{+3\%}
\providecommand{\JobsCOneBasewidthSixtyFour}{}\renewcommand{\JobsCOneBasewidthSixtyFour}{18/18}
\providecommand{\JobsCOneBasewidthTwoFiveSix}{}\renewcommand{\JobsCOneBasewidthTwoFiveSix}{18/18}
\providecommand{\JobsCOneBudgetL}{}\renewcommand{\JobsCOneBudgetL}{18/18}
\providecommand{\JobsCOneBudgetS}{}\renewcommand{\JobsCOneBudgetS}{18/18}
\providecommand{\JobsCOneCtx}{}\renewcommand{\JobsCOneCtx}{18/18}
\providecommand{\JobsCOneModelfree}{}\renewcommand{\JobsCOneModelfree}{18/18}
\providecommand{\JobsCOneRankOne}{}\renewcommand{\JobsCOneRankOne}{18/18}
\providecommand{\JobsCOneRankSixteen}{}\renewcommand{\JobsCOneRankSixteen}{18/18}
\providecommand{\JobsCOneTuneSelect}{}\renewcommand{\JobsCOneTuneSelect}{60/60}
\providecommand{\JobsCOneTuned}{}\renewcommand{\JobsCOneTuned}{36/36}
\providecommand{\JobsCOneTunedExtOne}{}\renewcommand{\JobsCOneTunedExtOne}{36/36}
\providecommand{\JobsCOneTunedExtTwo}{}\renewcommand{\JobsCOneTunedExtTwo}{18/18}
\providecommand{\JobsCOneWhere}{}\renewcommand{\JobsCOneWhere}{30/30}
\providecommand{\JobsCThreeDataFifty}{}\renewcommand{\JobsCThreeDataFifty}{18/18}
\providecommand{\JobsCThreeDataOneFiveZero}{}\renewcommand{\JobsCThreeDataOneFiveZero}{18/18}
\providecommand{\JobsCThreeFaults}{}\renewcommand{\JobsCThreeFaults}{30/30}
\providecommand{\JobsCThreeMismatchFifty}{}\renewcommand{\JobsCThreeMismatchFifty}{18/18}
\providecommand{\JobsCThreeMismatchTen}{}\renewcommand{\JobsCThreeMismatchTen}{18/18}
\providecommand{\JobsCThreeMismatchThirty}{}\renewcommand{\JobsCThreeMismatchThirty}{18/18}
\providecommand{\JobsCThreeMismatchTwenty}{}\renewcommand{\JobsCThreeMismatchTwenty}{18/18}
\providecommand{\JobsCThreeNoiseCt}{}\renewcommand{\JobsCThreeNoiseCt}{18/18}
\providecommand{\JobsCThreePartial}{}\renewcommand{\JobsCThreePartial}{18/18}
\providecommand{\JobsCThreePartialModelfree}{}\renewcommand{\JobsCThreePartialModelfree}{18/18}
\providecommand{\JobsCThreeRealistic}{}\renewcommand{\JobsCThreeRealistic}{30/30}
\providecommand{\JobsCThreeRegimes}{}\renewcommand{\JobsCThreeRegimes}{18/18}
\providecommand{\JobsCTwoFloorFive}{}\renewcommand{\JobsCTwoFloorFive}{18/18}
\providecommand{\JobsCTwoFloorOne}{}\renewcommand{\JobsCTwoFloorOne}{18/18}
\providecommand{\JobsCTwoFloorTwentyFive}{}\renewcommand{\JobsCTwoFloorTwentyFive}{18/18}
\providecommand{\JobsCTwoTruncDev}{}\renewcommand{\JobsCTwoTruncDev}{18/18}
\providecommand{\JobsCTwoTruncTest}{}\renewcommand{\JobsCTwoTruncTest}{18/18}
\providecommand{\JobsCTwoWindowAOne}{}\renewcommand{\JobsCTwoWindowAOne}{18/18}
\providecommand{\JobsCTwoWindowATwo}{}\renewcommand{\JobsCTwoWindowATwo}{9/9}
\providecommand{\JobsRZero}{}\renewcommand{\JobsRZero}{6/6}
\providecommand{\JobsTESTCOne}{}\renewcommand{\JobsTESTCOne}{18/18}
\providecommand{\JobsTESTCOneTuned}{}\renewcommand{\JobsTESTCOneTuned}{18/18}
\providecommand{\JobsTESTCOneTunedExtOne}{}\renewcommand{\JobsTESTCOneTunedExtOne}{36/36}
\providecommand{\JobsTESTCOneTunedExtTwo}{}\renewcommand{\JobsTESTCOneTunedExtTwo}{18/18}
\providecommand{\JobsTESTCThree}{}\renewcommand{\JobsTESTCThree}{18/18}
\providecommand{\MergedCThreeMismatchFifty}{}\renewcommand{\MergedCThreeMismatchFifty}{15/15}  % runs with the new data-driven baselines merged
\providecommand{\MergedCThreeMismatchTen}{}\renewcommand{\MergedCThreeMismatchTen}{15/15}  % runs with the new data-driven baselines merged
\providecommand{\MergedCThreeMismatchThirty}{}\renewcommand{\MergedCThreeMismatchThirty}{15/15}  % runs with the new data-driven baselines merged
\providecommand{\MergedCThreeMismatchTwenty}{}\renewcommand{\MergedCThreeMismatchTwenty}{15/15}  % runs with the new data-driven baselines merged
\providecommand{\MergedCThreeNoiseCt}{}\renewcommand{\MergedCThreeNoiseCt}{15/15}  % runs with the new data-driven baselines merged
\providecommand{\MergedCThreePartial}{}\renewcommand{\MergedCThreePartial}{15/15}  % runs with the new data-driven baselines merged
\providecommand{\MergedCThreeRealistic}{}\renewcommand{\MergedCThreeRealistic}{25/25}  % runs with the new data-driven baselines merged
\providecommand{\MergedTESTCThree}{}\renewcommand{\MergedTESTCThree}{14/18}  % runs with the new data-driven baselines merged
\providecommand{\ModelErrBatteryMismatchThirty}{}\renewcommand{\ModelErrBatteryMismatchThirty}{0.09}
\providecommand{\ModelErrBatteryPartial}{}\renewcommand{\ModelErrBatteryPartial}{0.22}
\providecommand{\ModelErrBioreactorMismatchThirty}{}\renewcommand{\ModelErrBioreactorMismatchThirty}{0.50}
\providecommand{\ModelErrBioreactorPartial}{}\renewcommand{\ModelErrBioreactorPartial}{2.15}
\providecommand{\ModelErrCartpoleBMismatchThirty}{}\renewcommand{\ModelErrCartpoleBMismatchThirty}{0.35}
\providecommand{\ModelErrCartpoleBPartial}{}\renewcommand{\ModelErrCartpoleBPartial}{0.03}
\providecommand{\ModelErrCartpoleMismatchThirty}{}\renewcommand{\ModelErrCartpoleMismatchThirty}{0.35}
\providecommand{\ModelErrCartpolePartial}{}\renewcommand{\ModelErrCartpolePartial}{0.03}
\providecommand{\ModelErrFhnMismatchThirty}{}\renewcommand{\ModelErrFhnMismatchThirty}{0.16}
\providecommand{\ModelErrFhnPartial}{}\renewcommand{\ModelErrFhnPartial}{0.04}
\providecommand{\ModelErrHighwayTrafficMismatchThirty}{}\renewcommand{\ModelErrHighwayTrafficMismatchThirty}{0.40}
\providecommand{\ModelErrHighwayTrafficPartial}{}\renewcommand{\ModelErrHighwayTrafficPartial}{0.82}
\providecommand{\ModelErrLorenzMismatchThirty}{}\renewcommand{\ModelErrLorenzMismatchThirty}{0.29}
\providecommand{\ModelErrLorenzPartial}{}\renewcommand{\ModelErrLorenzPartial}{0.71}
\providecommand{\ModelErrNetworkMismatchThirty}{}\renewcommand{\ModelErrNetworkMismatchThirty}{0.21}
\providecommand{\ModelErrNetworkPartial}{}\renewcommand{\ModelErrNetworkPartial}{0.22}
\providecommand{\ModelErrQuadTankMismatchThirty}{}\renewcommand{\ModelErrQuadTankMismatchThirty}{0.69}
\providecommand{\ModelErrQuadTankPartial}{}\renewcommand{\ModelErrQuadTankPartial}{0.15}
\providecommand{\ModelErrRosslerMismatchThirty}{}\renewcommand{\ModelErrRosslerMismatchThirty}{0.25}
\providecommand{\ModelErrRosslerPartial}{}\renewcommand{\ModelErrRosslerPartial}{0.30}
\providecommand{\ModelErrTwoLinkArmBMismatchThirty}{}\renewcommand{\ModelErrTwoLinkArmBMismatchThirty}{0.50}
\providecommand{\ModelErrTwoLinkArmBPartial}{}\renewcommand{\ModelErrTwoLinkArmBPartial}{0.02}
\providecommand{\ModelErrTwoLinkArmMismatchThirty}{}\renewcommand{\ModelErrTwoLinkArmMismatchThirty}{0.29}
\providecommand{\ModelErrTwoLinkArmPartial}{}\renewcommand{\ModelErrTwoLinkArmPartial}{0.04}
\providecommand{\ModelErrVdpMismatchThirty}{}\renewcommand{\ModelErrVdpMismatchThirty}{0.20}
\providecommand{\ModelErrVdpPartial}{}\renewcommand{\ModelErrVdpPartial}{0.66}
\providecommand{\OutcomeCOneBasewidthSixtyFour}{}\renewcommand{\OutcomeCOneBasewidthSixtyFour}{PASS}  % winner changes on shape_ood 0/5; sequence_ood 0/5 ('base-dependent' if >= 3/5 (p
\providecommand{\OutcomeCOneBasewidthTwoFiveSix}{}\renewcommand{\OutcomeCOneBasewidthTwoFiveSix}{PASS}  % winner changes on shape_ood 0/5; sequence_ood 0/5 ('base-dependent' if >= 3/5 (p
\providecommand{\OutcomeCOneBudget}{}\renewcommand{\OutcomeCOneBudget}{PASS}  % stable at BOTH budgets: shape_ood 5/5; sequence_ood 5/5
\providecommand{\OutcomeCOneBudgetL}{}\renewcommand{\OutcomeCOneBudgetL}{PASS}  % C1-where winner best-or-within-CI: shape_ood 5/5; sequence_ood 5/5 (need 4/5 (pr
\providecommand{\OutcomeCOneBudgetS}{}\renewcommand{\OutcomeCOneBudgetS}{PASS}  % C1-where winner best-or-within-CI: shape_ood 5/5; sequence_ood 5/5 (need 4/5 (pr
\providecommand{\OutcomeCOneCtx}{}\renewcommand{\OutcomeCOneCtx}{PASS}  % C1-where winner best-or-within-CI: shape_ood 5/5; sequence_ood 5/5 (need 4/5 (pr
\providecommand{\OutcomeCOneModelfree}{}\renewcommand{\OutcomeCOneModelfree}{FAIL}  % C1-where beats model-free phase 1 (same method, CI): shape 0/5; seq 0/5 (need 4/
\providecommand{\OutcomeCOneRankOne}{}\renewcommand{\OutcomeCOneRankOne}{FAIL}  % lora at this rank within +-10% of rank 4 or better (TOST non-inferiority, 90% CI
\providecommand{\OutcomeCOneRankSixteen}{}\renewcommand{\OutcomeCOneRankSixteen}{FAIL}  % lora at this rank within +-10% of rank 4 or better (TOST non-inferiority, 90% CI
\providecommand{\OutcomeCOneTuneSelect}{}\renewcommand{\OutcomeCOneTuneSelect}{REPORTED}  % selected: latent latent[lr=0.004,latent_depth=3] (-8.1% val vs default); concat 
\providecommand{\OutcomeCOneTuned}{}\renewcommand{\OutcomeCOneTuned}{FAIL}  % tuned HyperKKL (lora) beats tuned latent and concat: dev shape_ood 3/5 (need 4);
\providecommand{\OutcomeCOneTunedExtOne}{}\renewcommand{\OutcomeCOneTunedExtOne}{FAIL}  % EXTENDED 'ext1': tuned HyperKKL beats tuned latent and concat: dev shape_ood 4/5
\providecommand{\OutcomeCOneTunedExtTwo}{}\renewcommand{\OutcomeCOneTunedExtTwo}{FAIL}  % EXTENDED 'ext2': tuned HyperKKL beats tuned latent and concat: dev shape_ood 3/5
\providecommand{\OutcomeCOneWhere}{}\renewcommand{\OutcomeCOneWhere}{FAIL}  % location rule over latent / concat / film / HyperKKL (dense-subspace ablation ex
\providecommand{\OutcomeCThreeDataFifty}{}\renewcommand{\OutcomeCThreeDataFifty}{REPORTED}  % learning curve reported (Figure F3 companion); no pass/fail rule
\providecommand{\OutcomeCThreeDataOneFiveZero}{}\renewcommand{\OutcomeCThreeDataOneFiveZero}{REPORTED}  % learning curve reported (Figure F3 companion); no pass/fail rule
\providecommand{\OutcomeCThreeFaults}{}\renewcommand{\OutcomeCThreeFaults}{REPORTED}  % learned-wins regime (>= 4/5 (pre-reg. 4/6) vs the best filter of either kind, sh
\providecommand{\OutcomeCThreeMismatchCrossover}{}\renewcommand{\OutcomeCThreeMismatchCrossover}{FAIL}  % level outcomes: 10% PASS, 20% PASS, 30% PASS, 50% PASS; first winning level 10%
\providecommand{\OutcomeCThreeMismatchFifty}{}\renewcommand{\OutcomeCThreeMismatchFifty}{PASS}  % id: vs best filter 5/5, vs best working 5/5; seq: vs best filter 5/5, vs best wo
\providecommand{\OutcomeCThreeMismatchFiftyDataDriven}{}\renewcommand{\OutcomeCThreeMismatchFiftyDataDriven}{FAIL}  % best KKL vs best data-driven: id 0/5; seq 0/5 (merged runs 15/15)
\providecommand{\OutcomeCThreeMismatchTen}{}\renewcommand{\OutcomeCThreeMismatchTen}{PASS}  % id: vs best filter 5/5, vs best working 5/5; seq: vs best filter 4/5, vs best wo
\providecommand{\OutcomeCThreeMismatchTenDataDriven}{}\renewcommand{\OutcomeCThreeMismatchTenDataDriven}{FAIL}  % best KKL vs best data-driven: id 0/5; seq 0/5 (merged runs 15/15)
\providecommand{\OutcomeCThreeMismatchThirty}{}\renewcommand{\OutcomeCThreeMismatchThirty}{PASS}  % id: vs best filter 5/5, vs best working 5/5; seq: vs best filter 5/5, vs best wo
\providecommand{\OutcomeCThreeMismatchThirtyDataDriven}{}\renewcommand{\OutcomeCThreeMismatchThirtyDataDriven}{FAIL}  % best KKL vs best data-driven: id 0/5; seq 0/5 (merged runs 15/15)
\providecommand{\OutcomeCThreeMismatchTwenty}{}\renewcommand{\OutcomeCThreeMismatchTwenty}{PASS}  % id: vs best filter 5/5, vs best working 5/5; seq: vs best filter 5/5, vs best wo
\providecommand{\OutcomeCThreeMismatchTwentyDataDriven}{}\renewcommand{\OutcomeCThreeMismatchTwentyDataDriven}{FAIL}  % best KKL vs best data-driven: id 1/5; seq 1/5 (merged runs 15/15)
\providecommand{\OutcomeCThreeNoiseCt}{}\renewcommand{\OutcomeCThreeNoiseCt}{FAIL}  % id: vs best filter 0/5, vs best working 0/5; seq: vs best filter 0/5, vs best wo
\providecommand{\OutcomeCThreeNoiseCtDataDriven}{}\renewcommand{\OutcomeCThreeNoiseCtDataDriven}{FAIL}  % best KKL vs best data-driven: id 0/5; seq 0/5 (merged runs 15/15)
\providecommand{\OutcomeCThreePartial}{}\renewcommand{\OutcomeCThreePartial}{REPORTED}  % id: vs best filter 4/5, vs best working 4/5; seq: vs best filter 4/5, vs best wo
\providecommand{\OutcomeCThreePartialDataDriven}{}\renewcommand{\OutcomeCThreePartialDataDriven}{FAIL}  % best KKL vs best data-driven: id 0/5; seq 1/5 (merged runs 15/15)
\providecommand{\OutcomeCThreePartialModelfree}{}\renewcommand{\OutcomeCThreePartialModelfree}{FAIL}  % C3-partial (partial model in the PDE term) beats model-free phase 1: id 0/5; seq
\providecommand{\OutcomeCThreeRealistic}{}\renewcommand{\OutcomeCThreeRealistic}{PASS}  % id: vs best filter 5/5, vs best working 5/5; seq: vs best filter 5/5, vs best wo
\providecommand{\OutcomeCThreeRealisticDataDriven}{}\renewcommand{\OutcomeCThreeRealisticDataDriven}{FAIL}  % best KKL vs best data-driven: id 1/5; seq 2/5 (merged runs 25/25)
\providecommand{\OutcomeCThreeRegimes}{}\renewcommand{\OutcomeCThreeRegimes}{REPORTED}  % KKL-wins regimes (>= 4/5 (pre-reg. 4/6), sequence_ood, vs best grey-box): none o
\providecommand{\OutcomeCTwoFloorFive}{}\renewcommand{\OutcomeCTwoFloorFive}{FAIL}  % no loss vs no floor (HyperKKL, id & seq): 4/5 (need 5/5 (pre-reg. 5/6)); TOST(10
\providecommand{\OutcomeCTwoFloorOne}{}\renewcommand{\OutcomeCTwoFloorOne}{FAIL}  % no loss vs no floor (HyperKKL, id & seq): 3/5 (need 5/5 (pre-reg. 5/6)); TOST(10
\providecommand{\OutcomeCTwoFloorTwentyFive}{}\renewcommand{\OutcomeCTwoFloorTwentyFive}{FAIL}  % no loss vs no floor (HyperKKL, id & seq): 5/5 (need 5/5 (pre-reg. 5/6)); TOST(10
\providecommand{\OutcomeCTwoTruncDev}{}\renewcommand{\OutcomeCTwoTruncDev}{FAIL}  % (S) within +5% of inf at every arm in [L*, settle]: 2/5; (T) significantly worse
\providecommand{\OutcomeCTwoTruncTest}{}\renewcommand{\OutcomeCTwoTruncTest}{FAIL}  % (S) within +5% of inf at every arm in [L*, settle]: 5/6; (T) significantly worse
\providecommand{\OutcomeCTwoWindowAOne}{}\renewcommand{\OutcomeCTwoWindowAOne}{FAIL}  % e(L >= L*) within 5% of e(512): id 2/4; seq 2/4 (need 5/5 (pre-reg. 5/6)). Refin
\providecommand{\OutcomeCTwoWindowATwo}{}\renewcommand{\OutcomeCTwoWindowATwo}{FAIL}  % knee in samples changes by <= 2x: 0/1 (need 2/3)
\providecommand{\OutcomeOriginalCOneBasewidthSixtyFour}{}\renewcommand{\OutcomeOriginalCOneBasewidthSixtyFour}{PASS}  % winner changes on shape_ood 1/6; sequence_ood 0/6 ('base-dependent' if >= 3/6; P
\providecommand{\OutcomeOriginalCOneBasewidthTwoFiveSix}{}\renewcommand{\OutcomeOriginalCOneBasewidthTwoFiveSix}{PASS}  % winner changes on shape_ood 2/6; sequence_ood 2/6 ('base-dependent' if >= 3/6; P
\providecommand{\OutcomeOriginalCOneBudget}{}\renewcommand{\OutcomeOriginalCOneBudget}{PASS}  % stable at BOTH budgets: shape_ood 4/6; sequence_ood 4/6
\providecommand{\OutcomeOriginalCOneBudgetL}{}\renewcommand{\OutcomeOriginalCOneBudgetL}{PASS}  % C1-where winner best-or-within-CI: shape_ood 6/6; sequence_ood 6/6 (need 4/6); T
\providecommand{\OutcomeOriginalCOneBudgetS}{}\renewcommand{\OutcomeOriginalCOneBudgetS}{PASS}  % C1-where winner best-or-within-CI: shape_ood 4/6; sequence_ood 4/6 (need 4/6); T
\providecommand{\OutcomeOriginalCOneCtx}{}\renewcommand{\OutcomeOriginalCOneCtx}{PASS}  % C1-where winner best-or-within-CI: shape_ood 6/6; sequence_ood 6/6 (need 4/6); T
\providecommand{\OutcomeOriginalCOneModelfree}{}\renewcommand{\OutcomeOriginalCOneModelfree}{FAIL}  % C1-where beats model-free phase 1 (same method, CI): shape 0/6; seq 0/6 (need 4/
\providecommand{\OutcomeOriginalCOneWhere}{}\renewcommand{\OutcomeOriginalCOneWhere}{FAIL}  % preferred location: NONE. Weaker claim (not the pre-registered rule), weight-spa
\providecommand{\OutcomeOriginalCThreeFaults}{}\renewcommand{\OutcomeOriginalCThreeFaults}{REPORTED}  % learned-wins regime (>= 4/6 vs the best filter of either kind, shape and seq): n
\providecommand{\OutcomeOriginalCThreeMismatchCrossover}{}\renewcommand{\OutcomeOriginalCThreeMismatchCrossover}{FAIL}  % level outcomes: 10% PASS, 20% PASS, 30% PASS, 50% PASS; first winning level 10%
\providecommand{\OutcomeOriginalCThreeMismatchFifty}{}\renewcommand{\OutcomeOriginalCThreeMismatchFifty}{PASS}  % id: vs best filter 6/6, vs best working 6/6; seq: vs best filter 5/6, vs best wo
\providecommand{\OutcomeOriginalCThreeMismatchFiftyDataDriven}{}\renewcommand{\OutcomeOriginalCThreeMismatchFiftyDataDriven}{FAIL}  % best KKL vs best data-driven: id 0/6; seq 0/6 (merged runs 18/18)
\providecommand{\OutcomeOriginalCThreeMismatchTen}{}\renewcommand{\OutcomeOriginalCThreeMismatchTen}{PASS}  % id: vs best filter 5/6, vs best working 5/6; seq: vs best filter 4/6, vs best wo
\providecommand{\OutcomeOriginalCThreeMismatchTenDataDriven}{}\renewcommand{\OutcomeOriginalCThreeMismatchTenDataDriven}{FAIL}  % best KKL vs best data-driven: id 0/6; seq 1/6 (merged runs 18/18)
\providecommand{\OutcomeOriginalCThreeMismatchThirty}{}\renewcommand{\OutcomeOriginalCThreeMismatchThirty}{PASS}  % id: vs best filter 5/6, vs best working 5/6; seq: vs best filter 5/6, vs best wo
\providecommand{\OutcomeOriginalCThreeMismatchThirtyDataDriven}{}\renewcommand{\OutcomeOriginalCThreeMismatchThirtyDataDriven}{FAIL}  % best KKL vs best data-driven: id 0/6; seq 0/6 (merged runs 18/18)
\providecommand{\OutcomeOriginalCThreeMismatchTwenty}{}\renewcommand{\OutcomeOriginalCThreeMismatchTwenty}{PASS}  % id: vs best filter 5/6, vs best working 5/6; seq: vs best filter 5/6, vs best wo
\providecommand{\OutcomeOriginalCThreeMismatchTwentyDataDriven}{}\renewcommand{\OutcomeOriginalCThreeMismatchTwentyDataDriven}{FAIL}  % best KKL vs best data-driven: id 1/6; seq 1/6 (merged runs 18/18)
\providecommand{\OutcomeOriginalCThreeNoiseCt}{}\renewcommand{\OutcomeOriginalCThreeNoiseCt}{FAIL}  % id: vs best filter 0/6, vs best working 0/6; seq: vs best filter 0/6, vs best wo
\providecommand{\OutcomeOriginalCThreeNoiseCtDataDriven}{}\renewcommand{\OutcomeOriginalCThreeNoiseCtDataDriven}{FAIL}  % best KKL vs best data-driven: id 0/6; seq 0/6 (merged runs 18/18)
\providecommand{\OutcomeOriginalCThreePartial}{}\renewcommand{\OutcomeOriginalCThreePartial}{REPORTED}  % id: vs best filter 5/6, vs best working 5/6; seq: vs best filter 5/6, vs best wo
\providecommand{\OutcomeOriginalCThreePartialDataDriven}{}\renewcommand{\OutcomeOriginalCThreePartialDataDriven}{FAIL}  % best KKL vs best data-driven: id 1/6; seq 2/6 (merged runs 18/18)
\providecommand{\OutcomeOriginalCThreePartialModelfree}{}\renewcommand{\OutcomeOriginalCThreePartialModelfree}{FAIL}  % C3-partial (partial model in the PDE term) beats model-free phase 1: id 0/6; seq
\providecommand{\OutcomeOriginalCThreeRealistic}{}\renewcommand{\OutcomeOriginalCThreeRealistic}{PASS}  % id: vs best filter 6/6, vs best working 6/6; seq: vs best filter 6/6, vs best wo
\providecommand{\OutcomeOriginalCThreeRealisticDataDriven}{}\renewcommand{\OutcomeOriginalCThreeRealisticDataDriven}{FAIL}  % best KKL vs best data-driven: id 1/6; seq 2/6 (merged runs 30/30)
\providecommand{\OutcomeOriginalCThreeRegimes}{}\renewcommand{\OutcomeOriginalCThreeRegimes}{REPORTED}  % KKL-wins regimes (>= 4/6, sequence_ood, vs best grey-box): none of 6 conditions;
\providecommand{\OutcomeOriginalCTwoFloorFive}{}\renewcommand{\OutcomeOriginalCTwoFloorFive}{FAIL}  % no loss vs no floor (lora & hyper, id & seq): 4/6 (need 5/6); TOST(10%) non-infe
\providecommand{\OutcomeOriginalCTwoFloorOne}{}\renewcommand{\OutcomeOriginalCTwoFloorOne}{FAIL}  % no loss vs no floor (lora & hyper, id & seq): 4/6 (need 5/6); TOST(10%) non-infe
\providecommand{\OutcomeOriginalCTwoFloorTwentyFive}{}\renewcommand{\OutcomeOriginalCTwoFloorTwentyFive}{FAIL}  % no loss vs no floor (lora & hyper, id & seq): 6/6 (need 5/6); TOST(10%) non-infe
\providecommand{\OutcomeOriginalCTwoTruncDev}{}\renewcommand{\OutcomeOriginalCTwoTruncDev}{FAIL}  % (S) within +5% of inf at every arm in [L*, settle]: 3/6; (T) significantly worse
\providecommand{\OutcomeOriginalCTwoTruncTest}{}\renewcommand{\OutcomeOriginalCTwoTruncTest}{FAIL}  % (S) within +5% of inf at every arm in [L*, settle]: 2/6; (T) significantly worse
\providecommand{\OutcomeOriginalTESTCOne}{}\renewcommand{\OutcomeOriginalTESTCOne}{PASS}  % weight-space beats latent and concat: shape_ood 5/6; sequence_ood 5/6 (need 4/6 
\providecommand{\OutcomeOriginalTESTCThree}{}\renewcommand{\OutcomeOriginalTESTCThree}{PASS}  % id: vs best filter 4/6, vs best working 4/6; seq: vs best filter 5/6, vs best wo
\providecommand{\OutcomeOriginalTESTCThreeDataDriven}{}\renewcommand{\OutcomeOriginalTESTCThreeDataDriven}{FAIL*}  % best KKL vs best data-driven: id 1/4; seq 1/4 (merged runs 14/18)
\providecommand{\OutcomeRZero}{}\renewcommand{\OutcomeRZero}{PASS}  % no failed jobs: True; max rejection 0.0% (< 5%: True); every conditioned variant
\providecommand{\OutcomeTESTCOne}{}\renewcommand{\OutcomeTESTCOne}{PASS}  % HyperKKL (r=4) beats latent and concat: shape_ood 5/6; sequence_ood 5/6 (need 4/
\providecommand{\OutcomeTESTCOneTuned}{}\renewcommand{\OutcomeTESTCOneTuned}{FAIL}  % test half of the C1-tuned rule: shape_ood 4/6 (need 4); sequence_ood 3/6 (need 4
\providecommand{\OutcomeTESTCOneTunedExtOne}{}\renewcommand{\OutcomeTESTCOneTunedExtOne}{PASS}  % test half of the EXTENDED 'ext1' rule: shape_ood 4/6 (need 4); sequence_ood 4/6 
\providecommand{\OutcomeTESTCOneTunedExtTwo}{}\renewcommand{\OutcomeTESTCOneTunedExtTwo}{PASS}  % test half of the EXTENDED 'ext2' rule: shape_ood 5/6 (need 4); sequence_ood 5/6 
\providecommand{\OutcomeTESTCThree}{}\renewcommand{\OutcomeTESTCThree}{FAIL}  % id: vs best filter 3/6, vs best working 3/6; seq: vs best filter 3/6, vs best wo
\providecommand{\OutcomeTESTCThreeDataDriven}{}\renewcommand{\OutcomeTESTCThreeDataDriven}{FAIL*}  % best KKL vs best data-driven: id 1/4; seq 1/4 (merged runs 14/18)
}{}
\providecommand{\CthreeFaultsRegimes}{none of 15 conditions}
\providecommand{\CthreeRealisticIdVsWorking}{6/6}\providecommand{\CthreeRealisticSeqVsWorking}{6/6}
\providecommand{\ModelErrDuffingMismatchThirty}{0.02}\providecommand{\ModelErrDuffingPartial}{0.13}
\providecommand{\ModelErrQuadTankPartial}{0.15}\providecommand{\ModelErrBioreactorPartial}{2.15}
\providecommand{\ModelErrLorenzPartial}{0.71}\providecommand{\ModelErrVdpPartial}{0.66}
\providecommand{\ModelErrQuadTankMismatchThirty}{0.69}
\providecommand{\ConeLearnedOverEkfRangeSeq}{2.7--14.1}\providecommand{\ConeLearnedOverEkfMedianSeq}{5.9}
\begin{document}
\maketitle
\paragraph{Reading guide.} Every statement carries a status tag: [proven] under the stated assumptions, [proven, linear
plants], [conjecture], or [empirical]; each is followed by a one-sentence intuition, a proof, what it predicts and which
experiment node tests it (or that none does). In the paper the theory lives in the appendix; three statements go to the
main text: Theorem~\ref{thm:m2-window} (why a short window suffices; two-sided for linear plants, with the practical rule
of Corollary~\ref{cor:m2-rule}), Proposition~\ref{prop:m1-where} (where the input should enter: a stability margin that
depends on the input against extrapolation at unseen contexts), and Proposition~\ref{prop:m3-mismatch} (when a learned
observer beats a filter: a data-independent bias--variance floor under model incompleteness against the learned fit error
on the true plant). Proposition~\ref{prop:m3-iss} is a decomposition with measured constants, not a certificate, and is
presented as such. All numerical constants quoted are produced by \texttt{theory/checks/*.py}.
% =====================================================================
% Block M1 (contribution C1): shared assumptions, decoder target, where the input enters.
% Owns: Assumption ass:m1 (A1)-(A4) used by every block, eq:m1-TL, def:m1-suff, lem:m1-window,
% prop:m1-suff, rem:m1-koopman, rem:m1-capacity, prop:m1-where, cor:m1-amplitude, prop:m1-tight, cor:m1-ood.
% Cites: lem:m2-tail, thm:m2-window, prop:m2-current, lem:m2-sfm (block M2); lem:D, prop:m3-iss, prop:m3-mismatch (block M3).
% Checks: paragraph "Numerical checks" at the end of this section (theory/checks/m1_location.py, check_theory.py).
% Status legend used in all three sections: [proven] / [proven, linear plants] / [conjecture] / [empirical].
% =====================================================================
\section{Where the input enters a learned KKL observer}\label{sec:m1}

\paragraph{Setting and conventions.} Plant $\dot x=f(x,u)$, $y=h(x)$, $x\in\mathbb R^{n_x}$,
$y\in\mathbb R^{n_y}$; no input-affine structure. Inputs are applied through a zero-order hold at the
sample rate with values in a compact $U\ni0$, $\bar u:=\max_{U}|u|$; $u_t$ is the window $u|_{[t-L,t]}$
shifted to $[-L,0]$, a point of $W_L=U^{N_L}$ at sample instants when $L/dt=N_L\in\mathbb N$ (the decoder
runs only there), with the sup norm. $X(x,t;s;u)$ is the value at $s$ of the solution through $x$ at $t$.
The latent filter is $\dot z=Az+By$. \emph{Latent initial condition (one convention for Sections
\ref{sec:m1}--\ref{sec:m3}).} The \emph{label} and training latents start from $z(0)=0$, so
$z(t)=T^0(x(t),t;u):=\int_0^te^{A(t-s)}Bh(X(x(t),t;s;u))\,ds$ exactly; this finite-start map is the only
time-varying map used. The \emph{deployed} observer (decoder-side and latent-injection alike, \texttt{evaluate.py}
\texttt{z0\_map}) starts from $\hat z(0)=\hat T_0(\hat x_0)$, which adds the transient $e^{At}\hat T_0(\hat x_0)$ to
the latent; it is carried as the $\Delta z_0$ term of Proposition~\ref{prop:m3-iss} and never absorbed into a map.
The decoder target is the windowed map
\begin{equation}\label{eq:m1-TL}
T_L(x,w):=\int_{-L}^{0}e^{-A\tau}B\,h\big(X(x,0;\tau;w)\big)\,d\tau,
\end{equation}
which by Lemma~\ref{lem:m2-tail}(a) is within $\tfrac12\varepsilon(L)$ of $T^0(x(t),t;u)$ at $w=u_t$.
$\mathcal P_L$ is the set of pairs $(x(t),u_t)$, $t\ge L$, generated by admissible trajectories.

\begin{assumption}[standing; shared by all sections]\label{ass:m1}
Regularity: $f(\cdot,u)\in C^1$ uniformly in $u\in U$, $h\in C^2$.
\begin{enumerate}
\item[(A1)] \emph{Forward invariance.} A compact $\mathcal X$ contains every solution from the initial set
  $\mathcal X_0\subset\mathcal X$ under every admissible input, for all $t\ge0$.
\item[(A2)] \emph{Bounded output.} $\|h(x)\|\le\bar y$ on $\mathcal X$.
\item[(A3)] \emph{Hurwitz filter.} $-\alpha<0$ is the spectral abscissa of $A$ and
  $\kappa:=\sup_{t\ge0}e^{\alpha t}\|e^{At}\|<\infty$ ($\kappa=1$ for the normal Butterworth $A$ of the code).
\item[(A4)] \emph{Uniform Lipschitz injectivity of the windowed map.} There are $L_0$ and $\ell<\infty$ with
  $|x_1-x_2|\le\ell\,|T_L(x_1,w)-T_L(x_2,w)|$ for all $(x_1,w),(x_2,w)\in\overline{\mathcal P_L}$ and all
  $L\ge L_0$ (i.e.\ $\ell:=\sup_{L\ge L_0}\ell(L)<\infty$; Theorem~\ref{thm:m2-window}(a) varies $L$ at fixed $\ell$).
\end{enumerate}
\end{assumption}
\noindent\emph{Chaotic plants.} (A1)--(A3) hold for the dissipative chaotic benchmarks (Lorenz, R\"ossler), (A4) is
assumed for them as for every other plant, but its constant is not under control when $\alpha$ is below the
\emph{backward} expansion rate $\lambda_b$ of the flow (for Lorenz $\lambda_b\approx14.6$ against $\alpha=2$):
$\partial_xT^0$ contains $\int_0^te^{A(t-s)}B\,\partial h\,\partial_xX(x,t;s;u)\,ds$ and $\|\partial_xX\|$ grows like
$e^{\lambda_b(t-s)}$ along the backward flow, so the forward Lipschitz constant of $T^0(\cdot,t)$ grows with $t$ when
$\alpha<\lambda_b$; this is the eigenvalue condition of \citet{andrieu2006existence} (the spectrum of $A$ must
dominate the Lipschitz constant of $f$ for a $C^1$ stationary map; a H\"older map needs less), and it is a plausible
mechanism (not tested: no node sweeps $\alpha$ or measures $\lambda_b$) behind the phase-1 PDE residual
$\varepsilon_{\rm pde,rms}=74$ on Lorenz against $\le0.3$ elsewhere. Which results survive on such plants is stated at
each result: everything that uses only the true past and (A1)--(A2) does (Lemma~\ref{lem:m2-tail}(a),
Theorem~\ref{thm:m2-window}(a), Proposition~\ref{prop:m3-iss}); the counterfactual-past statements
(Lemma~\ref{lem:m2-tail}(b)) and the linear lower bounds (Theorem~\ref{thm:m2-window}(b),
Proposition~\ref{prop:m2-current}(b)) do not apply to them.

\begin{remark}[existence and injectivity: cited, and what is assumed]\label{rem:m1-exist}
$T^0$ is a $C^1$ solution of $\partial_xT\,f(x,u(t))+\partial_tT=AT+Bh(x)$ whenever backward solutions
exist on $[0,t]$ \citep[Lemma~1]{bernard2019luenberger}, and an observer
$\hat x=T^*(\hat z,t)$ converges if there are $\bar t$ and a concave class-$\mathcal K_\infty$ $\rho_T$ with
$|x_1-x_2|\le\rho_T(|T(x_1,t)-T(x_2,t)|)$ for all $t\ge\bar t$ \citep[Theorem~1]{bernard2019luenberger}.
For a fixed input that is backward distinguishable in finite time, $T^0(\cdot,t)$ is injective for almost
every complex diagonal $A$ with $n_y(n_x+1)$ eigenvalues \citep[Theorem~3]{bernard2019luenberger}, but the
exceptional set depends on the input \citep[Remark~5]{bernard2019luenberger}; the Lipschitz, uniform-in-time
version \citep[Theorem~2]{bernard2019luenberger} needs strong differential observability, inputs of class
$C^{m+1}$ (their time derivatives enter) and a gain $k$ above a threshold, which excludes held, switching
and PRBS inputs and a fixed $(A,B)$. For autonomous systems, $2n_x+1$ generic real dimensions suffice
\citep[Theorem~3.4]{brivadis2023further}, $2n_x+2$ for time-varying ones \citep[Remark~3.6]{brivadis2023further}.
(A4) is therefore \emph{assumed}: it is the quantitative, uniform-in-input version of what these results give
for each regular input. The code's $n_z=5$ for 2-D plants is the autonomous count; the time-varying count
$2n_x+2$ is not met, so (A4) at $n_z=5$ is an assumption whose constant must be reported, not a consequence of the cited
results; the phase-1 proxy for $\ell$ (\texttt{l\_dec\_q95}, Remark~\ref{rem:m3-cert}) is $24$ (vdp), $44$
(bioreactor), $299$ (duffing), $10$ (lorenz), $14$ (quad\_tank), $816$ (battery) for the C1-where phase~1 (used in
\texttt{check\_apriori\_vacuity}); the C2-window-a1 phase~1 of the results layer gives $23,44,295,10,14,806$
(Table~\ref{tab:m2-window-rule}).
\end{remark}

\begin{remark}[Koopman reading; which model classes the variants are]\label{rem:m1-koopman}
For the autonomous plant and a diagonalised filter $A=\mathrm{diag}(\lambda_i)$, $B=(b_i)$, the KKL PDE
$\partial_xT_i\,f=\lambda_iT_i+b_ih$ reads $(L_f-\lambda_i)T_i=b_ih$ with $L_f=f\cdot\nabla$ the Koopman generator on
observables, so
\[
T_i=b_i\,(L_f-\lambda_i)^{-1}h=b_i\int_0^\infty e^{\lambda_i\tau}\,h\big(X(x,0;-\tau;0)\big)\,d\tau ,
\]
the resolvent of the generator at the filter pole, i.e.\ the Laplace transform of the backward Koopman semigroup
applied to the output observable; if $h=\sum_jc_j\varphi_j$ with $L_f\varphi_j=\mu_j\varphi_j$ then
$T_i=\sum_jb_ic_j\varphi_j/(\mu_j-\lambda_i)$. For a linear plant this is the row identity
$p_i=b_iC(F-\lambda_iI)^{-1}$ of the Sylvester solution (check [7]); for a nonlinear plant it gives the spectral
reading of the ``$A$ fast enough'' condition above ($\operatorname{Re}\lambda_i$ must clear the backward spectrum) and of
the zero-input floor $\varepsilon_0$ of Remark~\ref{rem:m1-capacity}(i) as a spectral truncation. Under an input the
same object is the resolvent for the skew-product (state $\times$ input-shift) system, and the windowed map $T_L$ is its
finite-horizon version with the window as a truncated coordinate on the shift space; for a linear plant the context
$\dot q=Aq-PGu$ of Proposition~\ref{prop:m1-suff}(a) has the structure of a linear predictor with inputs
\citep{proctor2016dmdc,korda2018linear}. In these terms, \emph{latent injection} with a stationary $T_0$ is a lifted
Koopman model with inputs (its input-affine case $\Omega(z)u$, and the NLOX baseline $\omega(z)u$, are bilinear lifted
observers with a fixed output-injection gain $B$), whose stability requires the input-dependent generator $A+\partial_z\Phi$
to stay contracting (Proposition~\ref{prop:m1-where}(a)); \emph{decoder-side conditioning} keeps a fixed generator and
parametrises the left inverse of the lifting by the input context. Proposition~\ref{prop:m1-tight} is then the
statement that the Koopman spectrum moves with $u$ (eigenvalue $u-a$; rotation rate $\omega_0+u$) while a fixed $A$ does
not. We state this as positioning only; no result below depends on it.
\end{remark}

\subsection{The decoder target and sufficient contexts}

\begin{definition}[sufficient context]\label{def:m1-suff}
A context map $E:W_L\to\mathbb R^{d_c}$ is \emph{sufficient} if the latent and the context determine the state:
for $(x_i,w_i)\in\overline{\mathcal P_L}$, $E(w_1)=E(w_2)\Rightarrow|x_1-x_2|\le\ell\,|T_L(x_1,w_1)-T_L(x_2,w_2)|$.
It is $\epsilon_E$-sufficient if the same holds with an additive $\epsilon_E$ on the right.
\end{definition}
\noindent The window itself ($E=\mathrm{id}$) is sufficient by (A4). A lower-dimensional $E$ cannot be
injective on $W_L$ (invariance of domain), but it need not be: it only has to resolve what the latent leaves open.

\begin{lemma}[decoder target]\label{lem:m1-window}
Let $E$ be continuous and $\epsilon_E$-sufficient, and $S_L:=\{(T_L(x,w),E(w)):(x,w)\in\overline{\mathcal P_L}\}$.
If $\epsilon_E=0$ there is a continuous $G^*:S_L\to\mathcal X$ with $G^*(T_L(x,w),E(w))=x$, $\ell$-Lipschitz in $z$
for each $c$; extend $G^*(\cdot,c)$ to $\mathbb R^{n_z}$ with the same constant (Kirszbraun). Then along every
admissible trajectory, for $t\ge L$,
$\ |G^*(z(t),E(u_t))-x(t)|\le\tfrac{\ell}{2}\varepsilon(L).$
For $\epsilon_E>0$, any decoder $D$ satisfies $\sup|D(T_L(x,w),E(w))-x|\ge\tfrac12\delta_E$, where $\delta_E$ is the
largest $|x_1-x_2|$ over pairs with equal latent and equal context.
\end{lemma}
\noindent\emph{Intuition.} Over a fixed window the plant has no clock, so the state is a fixed function of
(latent, context) as soon as the context resolves the latent's ambiguity.
\begin{proof}
$\Psi(x,w)=(T_L(x,w),E(w))$ is continuous on the compact $\overline{\mathcal P_L}$ (under the hold, the backward
solution is a composition of $N_L$ $C^1$ flows with constant parameters). By sufficiency, $\Psi(x_1,w_1)=\Psi(x_2,w_2)$
implies $x_1=x_2$, so $\pi_x$ factors as $G^*\circ\Psi$; $G^*$ is continuous because $\Psi$ is a closed map
(continuous on a compact set into a Hausdorff space), and $\ell$-Lipschitz in $z$ by sufficiency. Then
$|G^*(z(t),c)-x(t)|=|G^*(z(t),c)-G^*(T_L(x(t),u_t),c)|\le\ell\,|z(t)-T_L(x(t),u_t)|$ and Lemma~\ref{lem:m2-tail}(a).
The lower bound: $D$ returns one value for two states $\delta_E$ apart.
\end{proof}

\begin{proposition}[sufficient contexts exist, and must carry history]\label{prop:m1-suff}
(a) \emph{Linear plants.} For $\dot x=Fx+Gu$, $y=Cx$, the windowed map is affine in $x$ with a
$w$-independent slope, $T_L(x,w)=P_Lx+Q_Lw$, $P_L=P-e^{AL}Pe^{-FL}$, $PF=AP+BC$. Hence the $n_z$-dimensional linear
context $c=Q_Lw$ is exactly sufficient whenever $P_L$ is injective ($\ell=\|P_L^+\|$), and as $L\to\infty$ (under the
washout $e^{AL}Pe^{-FL}\to0$, (H1) of Theorem~\ref{thm:m2-window}) it is the
output $q$ of the stable filter $\dot q=Aq-PGu$: a spectral filter bank with identity readout.
(b) \emph{Nonlinear plants.} The window ($d_c=n_uN_L$) is sufficient, and by (a) of Lemma~\ref{lem:m2-tail}
any context that determines $u_t$ for some $L$ is $\ell\varepsilon(L)$-sufficient for the infinite past:
$|x_1-x_2|\le\ell\,|T^0_1-T^0_2|+\ell\,\varepsilon(L)$ (each $T^0_i$ is within $\tfrac12\varepsilon(L)$ of $T_L$).
(c) \emph{History is necessary.} The current input alone is not sufficient in general
(Proposition~\ref{prop:m2-current}).
\end{proposition}
\begin{proof}
(a) $X(x,0;\tau;w)=e^{F\tau}x+(\text{linear in }w)$; substitute in \eqref{eq:m1-TL} and use
$\int_{-L}^0e^{-A\tau}BCe^{F\tau}d\tau=P-e^{AL}Pe^{-FL}$ (integrate $\tfrac{d}{d\tau}(e^{-A\tau}Pe^{F\tau})=e^{-A\tau}BCe^{F\tau}$).
If $Q_Lw_1=Q_Lw_2$ then $T_L(x_1,w_1)-T_L(x_2,w_2)=P_L(x_1-x_2)$. The limit: $z-Px$ obeys $\dot{(z-Px)}=A(z-Px)-PGu$.
(b) (A4) and Lemma~\ref{lem:m2-tail}(a).
\end{proof}
\noindent For nonlinear plants and contexts of fixed dimension below $n_uN_L$, finite stable filter banks with a
static readout uniformly approximate every time-invariant fading-memory operator on slew-limited bounded inputs
\citep[Theorem~2]{boyd1985fading}; applying it to the observer operator $(y,u)\mapsto x$ (a fading-memory operator
under (A1)--(A4), on slew-limited signals, and for the filters the theorem chooses rather than ours) motivates the
spectral context but is not a guarantee for the trained encoder.

\begin{remark}[expressivity at a fixed base; what C1 compares]\label{rem:m1-capacity}
(i) \emph{Shared floor.} Every conditioner with exact zero-input collapse, $D(\cdot,0)=T_0^+$, errs at $c=0$ by
$\varepsilon_0:=\sup_z|G^*(z,0)-T_0^+(z)|$: the zero-input target is the inverse of the zero-input windowed map,
not the learned stationary base. (ii) \emph{Additive (concat).} $T_0^+(z)+R(z,c)-R(z,0)$ approximates on $S_L$, as
the width of $R$ grows, every continuous target that equals $T_0^+$ on the zero slice: SiLU satisfies
$\mathrm{SiLU}(s)-\mathrm{SiLU}(-s)=s$, so a two-hidden-layer SiLU network contains every one-hidden-layer one, and
the universal approximation theorem for non-polynomial activations \citep{leshno1993multilayer} applies (with a
Tietze extension of the target). (iii) \emph{Base-capped (FiLM, LoRA, hypernetwork).} For each fixed $c$ the decoder
is a network $N_{\theta(c)}$ of the frozen base's architecture, so for every conditioner budget
\[
\inf_{D}\ \sup_{S_L}|D-G^*|\ \ge\ \sup_{c}\ \inf_{\theta\in\Theta}\ \sup_z|N_\theta(z)-G^*(z,c)|,
\]
with $\Theta$ the reachable parameter set. At the matched budget $\Theta$ is much smaller than the base's weight
space: the hypernetwork's linear bias-free heads place $\theta(c)$ in an affine subspace of dimension equal to its
trunk width; LoRA updates have the configured rank and leave biases fixed; FiLM never modulates the readout layer, so
$D(\cdot,c)=W_L\phi_c(z)+b_L$ with $(W_L,b_L)$ frozen. In the ReLU setting, a LoRA rank of order width times a depth
ratio suffices to represent a smaller target network exactly \citep{zeng2024expressive}; that result does not cover
tanh bases or small fixed ranks. (iv) \emph{Latent injection.} $\Phi$ is a universal class, but as an observer it
represents only targets of the stationary route: a stationary $T_0$ with an exact correction.
The theory therefore does not rank the variants by expressivity; at a matched budget on one base, C1 measures
inductive bias and optimisation within these classes.
\end{remark}

\subsection{Where the input enters: a stability--extrapolation trade-off}

Both observers run the filter on $y+\nu$, $\|\nu\|_\infty\le\delta_y$, with a context $\hat c=E(\hat u)$ of the input
$\hat u$ they are given; $c=E(u)$ and $\delta_c:=\sup_t|\hat c-c|\le\ell_E\|\hat u-u\|_\infty$ (Lemma~\ref{lem:m2-sfm}).
\begin{align*}
\text{(latent)}\quad &\dot{\hat z}=A\hat z+B(y+\nu)+\Phi(\hat z,\hat c), && \hat x=T_0^+(\hat z),\ \ \hat z(0)=T_0(\hat x_0),\\
\text{(decoder side)}\quad &\dot{\hat z}=A\hat z+B(y+\nu), && \hat x=D(\hat z,\hat c),\ \ \hat z(0)=\hat T_0(\hat x_0).
\end{align*}
For the latent variant, $T_0\in C^1$ is a stationary map with left inverse $T_0^+$ ($\ell_{\rm dec}$-Lipschitz,
$\varepsilon_{\rm dec}=\sup_{\mathcal X}|T_0^+T_0(x)-x|$), the residual along trajectories is
$r=\partial_xT_0(x)f(x,u)-AT_0(x)-Bh(x)-\Phi(T_0(x),c)$ with $|r|\le\varepsilon_{\rm res}$, and $\ell_{\Phi,c}$ is the
Lipschitz constant of $\Phi(z,\cdot)$ \emph{uniformly over} $z\in T_0(\mathcal X)+B(r)$ (the set used in (a)).
$\mu_2(M)=\lambda_{\max}(\tfrac12(M+M^\top))$ is the logarithmic norm.

\begin{proposition}[where the input enters]\label{prop:m1-where}
Under Assumption~\ref{ass:m1}:
\begin{enumerate}
\item[(a)] \emph{Latent injection: stability depends on the input} [proven]. Let
$\bar d:=\varepsilon_{\rm res}+\|B\|\delta_y+\ell_{\Phi,c}\,\delta_c$ and let $r\ge|e(0)|+\bar d/\mu$. If
$\mu_2\big(A+\partial_z\Phi(z,\hat c)\big)\le-\mu<0$ for all $z$ in the set $\mathcal Z_r:=T_0(\mathcal X)+B(r)$ and all
contexts $\hat c$ the observer meets, then $e=\hat z-T_0(x)$ stays in $B(r)$ for all $t$ (so the hypothesis is
checked on a set fixed in advance, not on the solution) and
\[
|e(t)|\le e^{-\mu t}|e(0)|+\tfrac1\mu\big(\varepsilon_{\rm res}+\|B\|\delta_y+\ell_{\Phi,c}\,\delta_c\big),\qquad
|\hat x-x|\le\ell_{\rm dec}|e|+\varepsilon_{\rm dec}.
\]
\item[(b)] \emph{Decoder side: stability is free, extrapolation is not} [proven]. The latent is an input-free LTI filter,
so $\hat z-z$ is driven only by $\nu$ and the initial transient, and $|\hat x-x|$ obeys the bound of
Proposition~\ref{prop:m3-iss} for every input, with the decoder's Jacobian norms $\ell_z(\hat c),\ell_c(\hat c)$ taken at
the contexts actually queried. Both conditioner families have $c$-uniform bounds on a bounded context set, but with very
different constants: for an additive conditioner $\ell_z(c)\le\ell_z(T_0^+)+2\sup\|\partial_zR\|$; for multiplicative
ones (FiLM, LoRA, hypernetwork) the available bound is the product form
$\ell_z(c)\le\prod_l\|W_l+\Delta W_l(c)\|\sup|\sigma'|^{d}$, finite because $\theta(c)$ is bounded (tanh trunk, bounded
$c$), but $10^1$--$10^5$ times larger at random initialisation (check [5]) before the tanh base saturates.
\end{enumerate}
\end{proposition}
\noindent\emph{Intuition.} Injecting the input makes the observer a recurrent network whose stability depends on how
strongly the injection reacts to the latent, but its decoder never sees an unfamiliar context; conditioning the decoder
keeps the recurrent part a fixed stable filter, but the network is queried at every unseen context directly.
\begin{proof}
(a) $\dot e=Ae+[\Phi(\hat z,\hat c)-\Phi(T_0x,\hat c)]+d$, $d=B\nu+[\Phi(T_0x,\hat c)-\Phi(T_0x,c)]-r$, $|d|\le\bar d$. While
$e\in B(r)$ the segment $[T_0x,\hat z]$ lies in the ball $T_0x+B(r)\subseteq\mathcal Z_r$ (convexity of the ball,
not of $\mathcal Z_r$, is what is used); write the bracket as $\bar Je$,
$\bar J=\int_0^1\partial_z\Phi(T_0x+se,\hat c)ds$, so $\mu_2(A+\bar J)\le-\mu$ by subadditivity and convexity of $\mu_2$,
and $D^+|e|\le-\mu|e|+\bar d$. The comparison function $\phi(t)=e^{-\mu t}|e(0)|+\bar d(1-e^{-\mu t})/\mu\le r$ bounds
$|e|$ as long as $e\in B(r)$, and $\phi<r$ whenever $|e|=r$ would be reached, so $B(r)$ is forward invariant and the
bound holds for all $t$ (contraction, \citealp{lohmiller1998contraction}; equivalently the perturbation lemmas of
\citealp[Ch.~9]{khalil2002nonlinear}). The state bound is the triangle inequality through $T_0^+(T_0(x))$.
(b) The latent statement is linearity; the bound is Proposition~\ref{prop:m3-iss}; the Jacobian bounds are the chain rule
($R$'s $z$-Jacobian does not involve $c$ multiplicatively, whereas modulated layers multiply $z$-features by $c$-dependent
weights), and $\theta(c)$ is bounded because the trunk is tanh and $c$ is bounded (Lemma~\ref{lem:m2-sfm}).
\end{proof}

\begin{corollary}[the certified margin shrinks with the context]\label{cor:m1-amplitude}
$\mu_2(A+\partial_z\Phi)\le\mu_2(A)+\ell_\Phi$, and $\mu_2(A)=-\alpha$ for normal $A$, so $\mu\ge\alpha-\ell_\Phi$
(for general $A$, the Gronwall version gives $\alpha-\kappa\ell_\Phi$). For an injection with exact zero-input
collapse, $\Phi(z,c)=M(z,c)-M(z,0)$, $M\in C^2$: $\ell_\Phi(c)\le\min\big(|c|\,L_{zc},\,2\sup\|\partial_zM\|\big)$ with
$L_{zc}=\sup\|\partial_c\partial_zM\|$, and $|c|\le\ell_E\|u\|_\infty$ (sup norm) since $E(0)=0$; for an input-affine
injection $\Omega(z)u$, $\ell_\Phi\le\ell_\Omega\|u\|_\infty$. The margin is evaluated at the \emph{observer's} context
$\hat c$, so a corrupted input shrinks it. The margin saturates: if $2\kappa\sup\|\partial_zM\|<\alpha$ the latent observer
is certified at every amplitude, and when the plant input saturates (as through a saturating actuator) the exact
correction's $z$-Lipschitz constant stays bounded too. The condition is \emph{sufficient only}: for the code's trained
injection the certified margin is already negative at moderate gain ($-5.47$, log-norm $+8.1$ in check [3]) while two
observer copies still contract, so a negative margin predicts nothing, and only a measured incremental rate does.
\end{corollary}

\begin{proposition}[the margin cannot be removed at a fixed spectrum]\label{prop:m1-tight}
[proven] There are bounded, observable plants with an \emph{exact} injection ($\varepsilon_{\rm res}=0$) whose latent
observer diverges while a decoder-side observer with a bounded sufficient context is within $\tfrac{\ell}{2}\varepsilon(L)$
of exact. For $\dot x=(u-a)x+b$, $y=x$, $a>\alpha>0$, scalar latent $\dot z=-\alpha z+y$ and constant $u\in(\alpha,a)$:
the plant converges to $b/(a-u)$, the stationary map is $T_0(x)=Px+q$, $P=1/(\alpha-a)$, $q=-Pb/\alpha$, the exact
injection is $\Phi^*(z,u)=u(z-q)$ ($\ell_\Phi=|u|$), and $\dot e=(u-\alpha)e$. The windowed map is
$T_L(x,w)=P_L(u)\,x+Q_L(u)$ with $P_L=\int_{-L}^0e^{(\alpha+u-a)\tau}d\tau>0$ and
$Q_L=\int_{-L}^0e^{\alpha\tau}\,b\,\tfrac{e^{(u-a)\tau}-1}{u-a}\,d\tau$ (both bounded functions of the window input), so the
decoder $\hat x=(\hat z-Q_L)/P_L$ errs by at most $\tfrac{\ell}{2}\varepsilon(L)$ with $\ell=1/P_L$
(Lemma~\ref{lem:m1-window}), for every $u\in(\alpha,a)$.
\end{proposition}
\begin{proof}
$T_0'(x)f(x,0)=P(b-ax)$ and $-\alpha T_0+x=(1-\alpha P)x-\alpha q=-aPx+Pb$; so $T_0$ solves the zero-input PDE and
$\partial_xT_0[f(x,u)-f(x,0)]=Pux=u(T_0(x)-q)$. Then $\dot e=-\alpha\hat z+x+u(\hat z-q)-P((u-a)x+b)=(u-\alpha)e$.
For the window: $X(x,0;\tau;u)=e^{(u-a)\tau}x+b(e^{(u-a)\tau}-1)/(u-a)$ in \eqref{eq:m1-TL}.
\end{proof}
\noindent Here the small-gain threshold $u=\alpha$ is attained exactly. The \emph{exact} time-varying map $z=m(t)x+n(t)$,
$\dot m=1-(\alpha+u-a)m$, $\dot n=-\alpha n-bm$, is also a sufficient context, but for $u<a-\alpha$ its coefficient
$-(\alpha+u-a)$ is positive and $m,n\to\infty$ (check: $m(30)=2\cdot10^{19}$ at $u=0.5$); it is outside the bounded
class of Proposition~\ref{prop:m1-where}(b), which is why the windowed context is the one used. The obstruction is
intrinsic to the stationary route \emph{at a fixed spectrum}: a designer removes it by choosing $\alpha>\bar u$ (here) or
$\alpha>\bar u\ell_\Omega\kappa$ (Corollary~\ref{cor:m1-amplitude}), i.e.\ by a gain that grows with the input bound, which
is also what \citet[Theorem~4]{bernard2019luenberger} require. A second example, the rotating oscillator
$\dot x=(\omega_0+u)Jx$, $y=x_1$ ($|x|$ is conserved for every input), with the drift map $P$ and the exact correction
$\Omega(z)=PJP^+z$, has error rate $\max\operatorname{Re}\lambda(A+uPJP^+)$: stable for all tested $u>-\omega_0$, but
unstable ($+2.65$, $+5.16$ at $u=-2,-4$, $\omega_0=1$, $n_z=5$, $\alpha=1.5$) when the rotation reverses, although the
plant is bounded and observable there; this is the input-independence required by the stationary route
\citep[Theorem~4]{bernard2019luenberger}, which also needs an input-affine plant and a gain growing with $\bar u$.
The decoder-side route needs only Assumption~\ref{ass:m1} and a sufficient context.

\paragraph{What the bounds predict, and what the current nodes test.} Under protocol v2 no tier changes the input
amplitude (every family lives in the same box), so the quantity that varies between tiers is the \emph{distance in context
space}, $d(\hat c,\mathcal C_{\rm train})$, not $|c|$. (P1) The latent variant's excess error on \texttt{shape\_ood} and
\texttt{sequence\_ood} is the extrapolation error of $\Phi$ at the new contexts, filtered at the measured incremental rate;
it degrades sharply only if that rate approaches zero. Decoder-side variants pay $\ell_z(\hat c)\,|\Delta z|$ and
$\ell_c(\hat c)\,|\Delta c|$ at the new contexts (Corollary~\ref{cor:m1-ood}), with the additive conditioner's $\ell_z$ small
and $c$-uniform. (P2) Input noise moves the latent variant's operating point to a corrupted context $\hat c$ (its margin is
evaluated there), whereas decoder-side variants stay stable for every input and pay $\ell_c\ell_E\delta_u$.
\emph{Status: untested.} No current node measures $\mu_2$ or the incremental rate on visited segments, $\ell_z(\hat c)$,
$\ell_c(\hat c)$, or the context distance per tier; under C3-faults \texttt{input\_noise} every variant, latent included,
degrades by about the same factor ($\approx10\times$ on vdp, read from the C3-faults node report; not a number of the
generated results layer), which neither supports nor falsifies (P2). The quantities to
record are eval-only on saved checkpoints: $\min_{c'\in\mathcal C_{\rm train}}\|\hat c-c'\|$ per tier (C1-where,
TEST-C1), the latent margin $\alpha-\kappa\ell_\Phi$ already computed by \texttt{metrics.py}, and the realised gains of
Definition~\ref{def:m3-measured}. Measured Lipschitz constants and log-norms are estimates on the visited set (lower
bounds of suprema), not certificates.

\begin{corollary}[out of distribution: Lipschitz extrapolation, and nothing more]\label{cor:m1-ood}
[proven] Let $\varepsilon_{\rm train}:=\sup|D(z',E(w'))-x'|$ over the training support $S_{\rm train}$ of triples
(label latent $z'=T^0$, window $w'$, state $x'$), and \emph{assume} the target $x=G^*(z,w)$ of Lemma~\ref{lem:m1-window}
has, on its domain $S_L$, Lipschitz constants $\ell$ in $z$ (this is (A4)) and $\ell_w$ in $w$ (sup norm; a new assumption,
not implied by sufficiency), and assume its extension beyond $S_L$ keeps these constants (a componentwise McShane
extension achieves them only up to a factor $\sqrt{n_x}$, which then multiplies $\ell$ and $\ell_w$ in (a)).
(a) For a query $(z,w)$ and any $(z',w')\in S_{\rm train}$, with $D$'s Jacobian norms $\ell^D_z,\ell^D_c$ on the segment,
\[
|D(z,E(w))-G^*(z,w)|\le\varepsilon_{\rm train}+\tfrac{\ell}{2}\varepsilon(L)+(\ell^D_z+\ell)|z-z'|+(\ell^D_c\,\ell_E+\ell_w)\|w-w'\|_\infty ,
\]
the $\tfrac{\ell}{2}\varepsilon(L)$ because $G^*(z',w')=x'$ only at $z'=T_L$, whereas the training latent is $T^0$; and
$|G^*(z(t),u_t)-x(t)|\le\tfrac{\ell}{2}\varepsilon(L)$ (Lemma~\ref{lem:m1-window}).
(b) The linear rate cannot be improved without structure. Among targets that are $2\lambda$-Lipschitz and agree with a
$\lambda$-Lipschitz one on $S_{\rm train}$, $G_\pm=G_0\pm\lambda\,d(\cdot,S_{\rm train})\,e$ ($|e|=1$) both fit the data,
so every learner errs by at least $\lambda\,d(q,S_{\rm train})$ at some query $q$. At the plant level, if $\varphi$ is a
diffeomorphism of $\mathcal X$ with $h\circ\varphi=h$, equal to the identity on an open neighbourhood of the states reached
under the training inputs (so that $D\varphi=I$ there), then the plants $f$ and $\varphi_*f$ produce identical training
data, and under an input that drives the state where $\varphi\neq\mathrm{id}$ they produce identical $(y,u)$ from $x_0$ and
$\varphi(x_0)$, so any observer errs by at least $\tfrac12|\varphi(x(t))-x(t)|$ on one of them.
\end{corollary}
\begin{proof}
(a) Triangle inequality through $D(z',E(w'))$, $G^*(z',w')$ and $G^*(T_L(x',w'),w')=x'$, with
$|E(w)-E(w')|\le\ell_E\|w-w'\|_\infty$ and $|z'-T_L(x',w')|\le\tfrac12\varepsilon(L)$.
(b) $d(\cdot,S)$ is 1-Lipschitz and vanishes on $S$; $|G_+-G_-|=2\lambda d$ and one learner output is at distance
$\ge\lambda d$ from one of them. For the plants: $\varphi$ maps solutions of $f$ to solutions of $\varphi_*f$ and preserves
$y$; on the training set it is the identity.
\end{proof}
\noindent\emph{What it predicts under protocol v2.} Distances are measured in input (window) space so that they are
comparable across encoders. In v2 at most $0.1\%$ of evaluation states lie outside the training box on any dev system
and tier (C1-where node report; not in the generated results layer), and the
input amplitude box is shared by all tiers, so $|z-z'|$ and the state-level construction of (b) are inert: the relevant
distance is $\|w-w'\|_\infty$ (equivalently $d(\hat c,\mathcal C_{\rm train})$ in context space), which \texttt{shape\_ood}
and \texttt{sequence\_ood} change through the waveform and the event timing at the same amplitude. The corollary therefore
predicts a shape/sequence-OOD excess that is at most linear in that context distance and, by (b), can be no better
without structure; it does not predict its size, and the context distance per tier is not yet reported (eval-only; C1-where,
TEST-C1). A model-based filter with the correct model extrapolates by structure; a learned observer cannot, beyond this
Lipschitz rate.

\paragraph{Zero input.} With $E(0)=0$, $\Delta(\cdot,0)=0$ in every decoder conditioner and $\Phi(\cdot,0)=0$, every
conditioned observer is the autonomous KKL observer at $u\equiv0$, so the autonomous theory
\citep[Theorems~2.2, 2.4, 2.7]{andrieu2006existence} and the robustness bound of \citet[Theorem~19]{niazi2025kkl}
(their $\mathrm{cond}(V)$ is our $\kappa$) are the special case $c\equiv0$ of Proposition~\ref{prop:m3-iss}.

\paragraph{Numerical checks.} \texttt{theory/checks/m1\_location.py}: [1] zero-input collapse of all encoders and
conditioners (errors $\le5\cdot10^{-16}$); [2] at fixed $c$, FiLM, LoRA and hypernetwork decoders equal base-architecture
networks with explicit weights ($\le3\cdot10^{-13}$); [3] the injection's $z$-Lipschitz constant is $0,0.61,1.22,2.33,4.26,6.82,7.29$ at
$|c|=0,0.25,0.5,1,2,8,32$ (linear, then saturating), and for two observer copies the measured incremental rate is below the
sup log-norm on the visited segments ($-1.52\le-1.23$; $-1.05\le-0.29$, where the small-gain margin is $-5.47$ and certifies
nothing; at the largest gain the log-norm is $+8.1$ and also certifies nothing while the copies still contract, so the
log-norm condition is sufficient only); [4] the rotating oscillator; [5] \emph{at random initialisation} (trained
decoders are not tested here) the $z$-Jacobian of concat stays below its uniform-in-$c$ bound ($\le1.25$ against $12$ for
$|c|\le256$) while FiLM, LoRA and hypernetwork reach $81$, $3.7\cdot10^5$, $2.3\cdot10^6$ and then fall (LoRA and
hyper to $0$ and $2.9$ at $|c|=256$): $c$-uniformly bounded, with a far larger constant, as stated in
Proposition~\ref{prop:m1-where}(b); [6] the linear windowed map is affine with $w$-independent slope to $10^{-15}$;
[7] the Koopman resolvent identity $p_i=b_iC(F-\lambda_iI)^{-1}$ (to $10^{-16}$) and its Laplace-transform form (to
$10^{-8}$); [8] the a-priori form of Proposition~\ref{prop:m1-where}(a): with $\mu=0.46$, $|e(0)|=0.5$, $\sup|d|=0.3$ the
error never leaves $B(r)$, $r=1.15$. \texttt{theory/checks/check\_theory.py}: \texttt{where\_input\_enters}
(Proposition~\ref{prop:m1-tight}: measured rate $u-\alpha$ to $10^{-3}$; windowed-context error equals
$\tfrac{\ell}{2}\varepsilon(L)$ to three digits at $L=6$, i.e.\ the bound is attained for a constant input; the exact
context's $m(t)$ reaches $10^{11}$--$10^{19}$).

% ---- bibliography entries used (for the main .bib) ----
% bernard2019luenberger: P. Bernard, V. Andrieu. Luenberger observers for nonautonomous nonlinear systems.
%   IEEE TAC 64(1):270--281, 2019. doi:10.1109/TAC.2018.2872202  (numbering checked on HAL hal-02050020)
% andrieu2006existence: V. Andrieu, L. Praly. On the existence of a Kazantzis--Kravaris/Luenberger observer.
%   SIAM J. Control Optim. 45(2):432--456, 2006.  (numbering checked on HAL hal-00362726)
% brivadis2023further: L. Brivadis, V. Andrieu, P. Bernard, U. Serres. Further remarks on KKL observers.
%   Systems & Control Letters 172:105429, 2023.  (numbering checked on HAL hal-03695863)
% niazi2025kkl: M.U.B. Niazi, J. Cao, M. Barreau, K.H. Johansson. KKL observer synthesis for nonlinear systems via
%   physics-informed learning. arXiv:2501.11655v2, 2025.
% boyd1985fading: S. Boyd, L.O. Chua. Fading memory and the problem of approximating nonlinear operators with
%   Volterra series. IEEE Trans. Circuits Syst. 32(11):1150--1161, 1985.
% leshno1993multilayer: M. Leshno, V.Ya. Lin, A. Pinkus, S. Schocken. Multilayer feedforward networks with a
%   nonpolynomial activation function can approximate any function. Neural Networks 6(6):861--867, 1993.
% zeng2024expressive: Y. Zeng, K. Lee. The expressive power of low-rank adaptation. ICLR 2024.
% lohmiller1998contraction: W. Lohmiller, J.-J.E. Slotine. On contraction analysis for non-linear systems.
%   Automatica 34(6):683--696, 1998.
% khalil2002nonlinear: H.K. Khalil. Nonlinear Systems, 3rd ed. Prentice Hall, 2002.
% korda2018linear: M. Korda, I. Mezic. Linear predictors for nonlinear dynamical systems: Koopman operator meets
%   model predictive control. Automatica 93:149--160, 2018.
% proctor2016dmdc: J.L. Proctor, S.L. Brunton, J.N. Kutz. Dynamic mode decomposition with control.
%   SIAM J. Applied Dynamical Systems 15(1):142--161, 2016.

% =====================================================================
% Block M2 (contribution C2): why a short input window suffices.
% Uses the shared assumption block (A1)-(A4) and the maps T^0, T_L of Section~\ref{sec:m1};
% sampled-data constants: Lemma~\ref{lem:D}. Checks: Table~\ref{tab:m2-checks} (theory/checks/m2_checks.py,
% check_theory.py). Owns: lem:m2-tail, thm:m2-window, cor:m2-rule, prop:m2-current, lem:m2-sfm, tab:m2-window-rule.
% No \appendix here: the only appendix (Lemma D) is issued once, in m3_error_bounds.tex.
% =====================================================================
\section{Why a short input window suffices}\label{sec:m2}

\paragraph{Setting.} Assumptions (A1)--(A4) of Section~\ref{sec:m1}: forward invariance of a compact
$\mathcal X$, $\|h\|\le\bar y$ on $\mathcal X$, $\|e^{At}\|\le\kappa e^{-\alpha t}$, and Lipschitz injectivity
(modulus $\ell$, uniform in $L\ge L_0$) of the windowed map. Here we also use that $-\alpha$ is \emph{exactly} the
spectral abscissa of $A$ and that the eigenvalues on $\operatorname{Re}\lambda=-\alpha$ are semisimple (otherwise no
finite $\kappa$ exists for that $\alpha$). The Butterworth real-Jordan $A$ is normal: each block is $sI+\omega J$,
so $\|e^{At}\|=e^{-\alpha t}$ and $\kappa=1$; its real parts are
$\alpha_j=\alpha\,\sin\frac{(2j-1)\pi}{2n_z}/\sin\frac{\pi}{2n_z}$, $j=1,\dots,\lceil n_z/2\rceil$, so
$\alpha=\alpha_1\le\alpha_j\le\alpha_{\max}=\alpha/\sin\frac{\pi}{2n_z}$. The \emph{label} latent starts at
$z(0)=0$ (the convention of Section~\ref{sec:m1}), hence $z(t)=T^0(x(t),t;u)$ exactly; the deployed observer's
initial transient is a separate term (Proposition~\ref{prop:m3-iss}) and plays no role in this section.
$T_L(x,w)$ is the windowed map \eqref{eq:m1-TL} (zero history before the window), and
\[
 \varepsilon(L):=\frac{2\kappa\|B\|\bar y}{\alpha}\,e^{-\alpha L}.
\]

% ---------------------------------------------------------------------
\begin{lemma}[fading-memory tail]\label{lem:m2-tail}
(a) [proven] Along every admissible trajectory and $t\ge L$,
$\ \|T^0(x(t),t;u)-T_L(x(t),u_t)\|\le\tfrac12\varepsilon(L)$.
(b) [proven under its hypothesis] If $u=v$ on $[t-L,t]$ and \emph{both} backward solutions from $x$ stay in
$\mathcal X$ on $[0,t]$, then
\[
 T^0(x,t;u)-T^0(x,t;v)=e^{AL}\big[T^0(\xi,t-L;u)-T^0(\xi,t-L;v)\big],\qquad
 \|T^0(x,t;u)-T^0(x,t;v)\|\le\varepsilon(L),
\]
with $\xi$ the common backward state at $t-L$. In particular the map forgets at exactly rate $\alpha$:
for a unit left eigenvector $p$ of an eigenvalue with real part $-\alpha$,
$|p^*[T^0(x,t;u)-T^0(x,t;v)]|=e^{-\alpha L}\,|p^*[T^0(\xi,t-L;u)-T^0(\xi,t-L;v)]|$.
\end{lemma}
\emph{Intuition.} Inside the window the two backward solutions are one curve; everything older reaches
the latent through $e^{AL}$.
\begin{proof}
(a) $T^0-T_L$ is the integral over $[0,t-L]$ of $e^{A(t-s)}Bh(x(s))$ along the true past, which lies in
$\mathcal X$ by (A1); bound the integrand by $\kappa\|B\|\bar y e^{-\alpha(t-s)}$. (b) On $[t-L,t]$ both
backward solutions solve the same ODE from $x$, so they coincide; split \eqref{eq:m1-TL}-type integrals at
$t-L$ and use the flow property: the window parts cancel and the rest is $e^{AL}T^0(\xi,t-L;\cdot)$. Bound
each map's remainder as in (a). Finally $p^*e^{AL}=e^{\lambda L}p^*$, $|e^{\lambda L}|=e^{-\alpha L}$.
\end{proof}
No Lipschitz property of $f$ in $u$ is used; the plant enters only through $\bar y$. \emph{Where (b) does not
apply.} Part (a) uses only the true past and (A1), so it holds for every benchmark. Part (b) needs the
\emph{counterfactual} backward solution (from $x$ under $v$) to stay in $\mathcal X$; for a dissipative chaotic flow
this fails generically, because backward time expands at the rate $\lambda_b$ of the most contracting forward
direction (Lorenz: $\lambda_b\approx14.6$, so a counterfactual past leaves any compact set within a fraction of a
second). For Lorenz and R\"ossler, (b) and everything that rests on it (Proposition~\ref{prop:m1-suff}(b)'s
$\ell\varepsilon(L)$-sufficiency of a context that determines the window; the lower bounds below) are not
available; the sufficient direction (a) is. The zero-order hold changes nothing: if the held samples agree for
$j=k-M,\dots,k-1$, apply the lemma at $t_k$ with $L=M\,dt$ (the sample $u_k$ has not acted yet); the deployed
discrete filter has the same bounds with $e^{-\alpha L}\mapsto e^{-\alpha M dt}$ and constants no larger than the
continuous ones, plus $e_{\rm disc}$ (Lemma~\ref{lem:D}).

% ---------------------------------------------------------------------
\begin{theorem}[window rule for the observer]\label{thm:m2-window}
Let $L^*(\eta)$ be the shortest window for which some decoder $D(z,u_t)$ has worst-case error at most $\eta$.
\begin{enumerate}
\item[(a)] (\emph{sufficient; every plant under (A1)--(A4)}) [proven] $D=T_L^*$ achieves
$\tfrac{\ell}{2}\varepsilon(L)$, so
\[
L^*(\eta)\ \le\ \frac1\alpha\ln\frac{\kappa\ell\|B\|\bar y}{\alpha\eta}\ =\ \frac{\ln(1/\eta)+\ln\ell+\ln(\kappa\|B\|\bar y/\alpha)}{\alpha}.
\]
\item[(b)] (\emph{necessary; proven for linear plants}) For every decoder and every plant,
$\sup\|D-x\|\ge\frac12\delta(L)$, where the fibre diameter $\delta(L)$ is the largest $\|x-x'\|$ over
\emph{admissible} pairs with $T^0(x,t;u)=T^0(x',t;v)$ and $u=v$ on $[t-L,t]$ (each state reached at $t$ under its own
past). For a linear plant $\dot x=Fx+Gu$, $y=Cx$, normal $A$, $(A,PG)$ controllable, symmetric inputs $|u|\le\bar u$, under
\begin{itemize}
\item[(H1)] \emph{washout}: $e^{At}Pe^{-Ft}\to0$, so that after a transient $z=Px+q$ with $PF=AP+BC$, $P$
injective, $q(t)=-\int_{-\infty}^te^{A(t-r)}PGu(r)dr$ (this holds iff $\alpha$ exceeds the fastest
\emph{contraction} rate of $F$; automatically for neutrally stable $F$);
\item[(H2)] \emph{overlapping reachable sets}: for any two admissible input histories and every $d$ with
$\|d\|\le2\kappa R_qe^{-\alpha L}/\sigma_{\min}(P)$ there is a state $x$ reachable at $t$ under the first history with
$x+d$ reachable under the second. (Every $x\in\mathcal X$ being reachable under every history is \emph{not} required and
is false under (A1) whenever the forced response is nonzero: the reachable set under one history is the translate
$e^{Ft}\mathcal X_0+(\text{forced response})$, a strict subset of $\mathcal X$. (H2) holds for $F$ skew-symmetric or
$F=0$ when $\mathcal X_0$ contains a ball whose radius exceeds the forced-response bound plus $\|d\|/2$, as in the
oscillator and integrator examples; it fails for contracting $F$, where the reachable set under a fixed history shrinks
to a point and the fibre, hence the necessary window, is \emph{smaller});
\end{itemize}
the fibre diameter satisfies
\[
 \frac{r_q}{\|P\|}\,e^{-\alpha_k L}\ \le\ \delta(L)\ \le\ \frac{2\kappa R_q}{\sigma_{\min}(P)}\,e^{-\alpha L},
\]
where $R_q=\kappa\|PG\|\bar u/\alpha$, $r_q>0$ is the radius of a ball contained in the set $Q$ of reachable $q$,
and $\alpha_k$ is the real part of the first mode group such that the slowest groups up to it span a dimension
$\ge n_z-n_x+1$. Hence $\alpha\le\alpha_k\le\alpha_{\max}$ and
\[
 \frac1{\alpha_k}\ln\frac{r_q}{2\|P\|\eta}\ \le\ L^*(\eta)\ \le\ \frac1\alpha\ln\frac{\kappa\ell\|B\|\bar y}{\alpha\eta},
 \qquad L^*(\eta)=\Theta\!\Big(\frac1\alpha\ln\frac1\eta\Big)\ \ \text{(linear plants)},
\]
with constants that depend on $n_z$ only through $\alpha_k/\alpha\le1/\sin\frac{\pi}{2n_z}$ ($\alpha_k/\alpha=2.618$ at
$n_z=5$, $n_x=2$; $4.049$ at $n_z=7$, $n_x=3$).
\end{enumerate}
\end{theorem}
\emph{Intuition.} The decoder cannot recover what the window does not contain, and what it does not contain
reaches the latent through $e^{AL}$; a decoder may lean on the faster modes, but never faster than $\alpha_{\max}$.
\begin{proof}
(a) Lemma~\ref{lem:m1-window} with Lemma~\ref{lem:m2-tail}(a). (b) Two points of the fibre receive the same
input $(z,u_t)$, so $D$ errs by $\ge\frac12\|x-x'\|$ on one of them. Linear case: by (H1) the window parts of $q$
coincide, so $Px+q=Px'+q'$ iff $P(x-x')=e^{AL}(m'-m)$ with $m,m'\in Q$ the history terms at $t-L$; every such $d=x'-x$
obeys the upper bound below, so by (H2) the pair can be realised by admissible states, and $\delta(L)=\sup\{\|d\|:e^{-AL}Pd\in Q-Q\}$. Upper:
$\|d\|\le\|Pd\|/\sigma_{\min}(P)\le\|e^{AL}\|\,2R_q/\sigma_{\min}(P)$. Lower: $\dim\operatorname{range}P=n_x$ and
the slow span $S_k$ has dimension $\ge n_z-n_x+1$, so some unit $d$ has $Pd\in S_k$; for normal $A$,
$\|e^{-AL}Pd\|\le e^{\alpha_kL}\|P\|$, and $Q-Q\supseteq Q\supseteq B(r_q)$ ($0\in Q$, controllability), so
$d\,r_qe^{-\alpha_kL}/\|P\|$ is feasible.
\end{proof}
\emph{What is proven and what is not.} The two-sided $\Theta(\ln(1/\eta)/\alpha)$ is a statement about linear plants
(with the hidden factor $\alpha_k/\alpha$ in the lower constant at $n_z=5$). For nonlinear plants only (a) is proven;
``(b) with rate between $\alpha$ and $\alpha_{\max}$'' is a \textbf{conjecture}, and for the chaotic benchmarks even
the counterfactual-past step is unavailable (Lemma~\ref{lem:m2-tail}). \emph{The fibre is a needle.} Writing
$Q=\{\int_0^\infty e^{As}PGu\,ds:|u|\le\bar u\}$, a symmetric convex body with support function
$h_Q(w)=\bar u\int_0^\infty|w^\top e^{As}PG|\,ds$, the fibre diameter is computed exactly as
$\delta(L)=\max_{|d|=1}2/g_Q(e^{-AL}Pd)$ with $g_Q$ the gauge of $Q$ (a linear programme on a Riemann discretisation
of $Q$). For the oscillator with $n_x=2$ the maximiser is the single direction $d$ with $Pd\in S_k$: $\delta(L)$
decays at rate $1.000\alpha$ ($n_z=3$) and $2.617\alpha$ ($n_z=5$; $\alpha_k=2.618\alpha$), every other direction
decays at $\approx\alpha_{\max}$, and the angular half-width of the needle shrinks from $3^\circ$ to $1^\circ$ over
$L=2$--$6$ (Table~\ref{tab:m2-checks}). An RMS error curve averages over directions, so a sweep's knee on a log scale
is expected between the slopes $-\alpha_k$ and $-\alpha_{\max}$, with the slow needle visible only at the largest $L$.
\emph{Tested by} C2-window-a1 (six dev systems, one bandwidth; pre-registered rule ``$e(L\ge\ln(1/\eta)/\alpha)$ within
$5\%$ of $e(512)$ on $\ge5/6$'': \textbf{FAIL}, \texttt{id} $2/5$, \texttt{sequence\_ood} $3/5$) and C2-window-a2 (battery,
duffing, vdp at $2\alpha$, the only test of the $1/\alpha$ scaling; rule ``knee in samples changes by $\le2\times$ on
$\ge2/3$'': \textbf{FAIL}, $1/2$). The $\ell$-free rule that was pre-registered is the sufficient bound with $\ell$ dropped;
the corollary below restores $\ell$ and is \emph{post hoc} consistent with the largest failure (battery), not with lorenz
(Table~\ref{tab:m2-window-rule}). Source:
\texttt{paper/generated/RESULTS.md}.

% ---------------------------------------------------------------------
\begin{corollary}[practical window rule, and why one plant needs more]\label{cor:m2-rule}
[proven: (a) of the theorem, rewritten; its use as an explanation of the sweep is \emph{post hoc}, not pre-registered]
For a target relative error $\eta$ the window that is \emph{sufficient} is
\[
 L^*(\eta)\ \approx\ \frac{\ln(1/\eta)+\ln\ell}{\alpha}\qquad\text{(plus the design constant }\ln(\kappa\|B\|\bar y/\alpha)/\alpha\text{)},
\]
where $\ell$ is the Lipschitz constant of the inverse of the windowed map (A4). The $\ln(1/\eta)/\alpha$ part is
plant-independent; the plant enters through $\ln\ell$ and, in the design constant, through $\bar y$. The lower bound of
the theorem (linear plants) can only \emph{shorten} the necessary window (rate $\alpha_k\ge\alpha$), so a plant whose
sufficient window exceeds $\ln(1/\eta)/\alpha$ does so through $\ell$ (or $\bar y$), not through the fading-memory rate.
\end{corollary}
\emph{Integrator-type states (coulomb counting; a heuristic reading, not a proven statement).} If a state enters the output only weakly, as the battery's state of
charge enters the terminal voltage through the open-circuit-voltage slope $\partial V_{\rm oc}/\partial{\rm SOC}$, then
the latent carries that state at gain $\propto$ slope and inverting it costs $\ell\propto1/{\rm slope}$: large
$\ell$, long $L^*$. The state itself is the integral of the input (coulomb counting), so a decoder must integrate the
window, which is possible only over the window's length; the memory needed is $\ln\ell/\alpha$ longer than for a
well-observed state. This is the sufficient side of the theorem; the digest's earlier attribution of the battery to
``the Lemma~1$'$ worst case'' is withdrawn: Lemma~\ref{lem:m2-tail}(b) is plant-independent (rate exactly $\alpha$ for
every plant) and the lower bound at $n_z-n_x+1=n_z$ (scalar integrator) gives $\alpha_k=\alpha_{\max}$, the
\emph{shortest} necessary window. \emph{Proxy for $\ell$.} No node measures $\ell$ itself; the phase-1 diagnostic
\texttt{l\_dec\_q95} (random-direction quotient of the \emph{stationary} base decoder, $\approx\|J\|_F/\sqrt{n_z}$,
Remark~\ref{rem:m3-cert}) is used as its proxy in Table~\ref{tab:m2-window-rule}. It is an estimate on data of a
related constant, not the constant of the bound.

\begin{table}[ht]\centering\small
\caption{Post-hoc reading of the C2-window-a1 sweep (lora, \texttt{sequence\_ood}, NRMSE at $L\in\{1,\dots,512\}$
samples, $\eta=0.05$). $L_0=\ln(1/\eta)/(\alpha dt)$ is the pre-registered $\ell$-free rule and
$L^*=(\ln(1/\eta)+\ln\ell)/(\alpha dt)$ the rule of Corollary~\ref{cor:m2-rule} with $\ell=\texttt{l\_dec\_q95}$ of the
node's own phase~1 (both in samples; the design constant is omitted; not pre-registered). ``ratio'' is
$\max_{L\ge L_\cdot}e(L)/e(512)$ on the dyadic grid; the node's criterion is $\le1.05$; n.t.\ = not testable (no grid
window in $[L_\cdot,512)$, so the ratio is $e(512)/e(512)$ and says nothing; the results layer leaves these out).
$\ell$ and $L^*$ are the results-layer values. Monotone: no window worse than a shorter one by $>10\%$.}\label{tab:m2-window-rule}
\footnotesize
\begin{tabular}{@{}lrrrrrrccl@{}}\hline
system & $n_z$ & $\alpha$ & $\alpha dt$ & $\ell$ & $L_0$ & $L^*$ & ratio, $L\ge L_0$ & ratio, $L\ge L^*$ & monotone\\\hline
battery    & 7 & 1.5 & 0.0223 & 806 & 135 & 435  & 1.73 (fails) & n.t. & yes\\
bioreactor & 5 & 0.5 & 0.0309 & 44  & 97  & 219  & 1.00 & 1.00 & no (min at 128; $2.4\times$ worse at 512)\\
duffing    & 5 & 3.0 & 0.0309 & 295 & 97  & 281  & 1.00 & n.t. & yes\\
lorenz     & 7 & 2.0 & 0.0223 & 10  & 135 & 239  & 1.07 (fails) & 1.07 (fails) & yes\\
quad\_tank & 6 & 0.1 & 0.0050 & 14  & 599 & 1121 & n.t. & n.t. & yes (falling at 512)\\
vdp        & 5 & 1.5 & 0.0309 & 23  & 97  & 199  & 1.00 & 1.00 & yes\\\hline
\end{tabular}
\end{table}
\noindent\emph{Reading the table (post hoc).} The pre-registered $\ell$-free rule fails where $\ell$ is largest
(battery, $1.73\times$ the $L=512$ error at $L=256$) and marginally on lorenz ($1.07$). With $\ln\ell$ restored the
rule can be evaluated only on bioreactor, lorenz and vdp: on battery, duffing and quad\_tank the refined $L^*$ leaves no
grid window below $512$, so the refined rule is \emph{untested} there (it does not ``hold''); it holds on bioreactor and
vdp and still fails by $7\%$ on lorenz ($2/3$ on \texttt{id} and $2/3$ on \texttt{sequence\_ood} in the results layer).
For battery the refined rule only says that its failure at $L\ge L_0$ is not in contradiction with the theory. This is what a
sufficient-side statement can do: the measured knee sits at or below it (duffing's knee is at $32$ samples against
$L^*=281$; the bound is far from tight there), and it does not predict the knee. Bioreactor is non-monotone (longer
windows are \emph{worse}: an optimisation or over-fitting effect that no truncation bound describes) and quad\_tank's
fixed $dt=0.05$ puts both rules beyond the grid, as predicted below. The refined rule is offered as the theory's
explanation of why the pre-registered test failed, not as a passed test; a pre-registered test of it needs a new node.
The $L=1$ row of the same sweep (current input only, Proposition~\ref{prop:m2-current}(a)) is $+143\%$ to $+1522\%$ worse
than $L=512$ on every dev system (\texttt{sequence\_ood}, lora; paired changes from the C2-window-a1 node report, not a
number of the generated results layer).

% ---------------------------------------------------------------------
\begin{proposition}[the current input is not enough]\label{prop:m2-current}
(a) [proven] Under the zero-order hold, $T^0(\cdot,t_k;u)$ does not depend on $u_k$.
(b) [proven, linear SISO plants] For the linear plant of Theorem~\ref{thm:m2-window} with $n_y=n_u=1$ under (H1)--(H2), but without assuming $(A,PG)$
controllable,
$(z,u_t)$ determines $x$ iff $\operatorname{range}P\cap\langle A\,|\,PG\rangle=\{0\}$, for every $L\ge0$. If $(A,B)$
is controllable, the eigenvalues of $A$ are distinct and none is a transmission zero of $H(s)=C(sI-F)^{-1}G$, then
$\langle A|PG\rangle=\mathbb R^{n_z}$ and the condition fails \emph{for every} $n_z$: no latent dimension makes the
latent store the history. (For $n_y>1$ the condition $p_i^*BH(\lambda_i)\ne0$ is stronger than ``not a transmission
zero''.)
(c) [proven, local] Parametrise the input before $t-L$ by $\theta\in\mathbb R^p$ (e.g.\ held levels in the
interior of $U$), keep the window fixed, and let $\Psi(x,\theta)$ be the resulting latent at $t$. If $\Psi$ is
$C^1$ near $(\bar x,\bar\theta)$ with trajectories in the interior of $\mathcal X$ and
$\operatorname{rank}\partial_\theta\Psi(\bar x,\bar\theta)=n_z$, then every $x$ in a neighbourhood of $\bar x$ has a past
$\theta(x)$ with the same latent and the same window; any decoder $D(z,u_t)$ errs by at least half that
neighbourhood's radius.
\end{proposition}
\emph{Intuition.} With a time-varying map the past input is part of what must be inverted, and the latent
alone cannot hold it.
\begin{proof}
(a) The backward solution on $(-\infty,t_k]$ uses $u$ on $[t_j,t_{j+1})$, $j\le k-1$, only.
(b) By Theorem~\ref{thm:m2-window}(b), ambiguity exists iff $P d=e^{AL}\Delta m$ for some $d\ne0$ and admissible
$\Delta m$; the reachable differences form a neighbourhood of $0$ in $\langle A|PG\rangle$, which $e^{AL}$ maps onto
itself. For distinct eigenvalues, controllability of $(A,PG)$ is $p_i^*PG\ne0$ for every left eigenvector,
and $p_i^*PG=-p_i^*B\,H(\lambda_i)$ (from $p_i^*P(F-\lambda_iI)=p_i^*BC$).
(c) Implicit function theorem: choose $n_z$ components of $\theta$ with invertible $\partial\Psi$-block and solve
$\Psi(x,\theta)=\Psi(\bar x,\bar\theta)$ for them as a $C^1$ function of $x$.
\end{proof}
The rank condition in (c) is checkable on a simulator: for Van der Pol with the saturating actuator,
twelve held levels before the window give $\operatorname{rank}\partial_\theta\Psi=5=n_z$ (smallest singular value
$1.6\cdot10^{-4}$: the check shows existence of the ambiguity, not its size, which is $\sigma_{\min}\times$ the
admissible range of $\theta$). Example: the oscillator $\dot x_1=\omega_0x_2$, $\dot x_2=-\omega_0x_1+u$, $y=x_1$,
$n_z=n_x=2$, with $v\equiv0$, $u=a\mathbf 1_{(-\infty,t-L)}$: $x$ and $x'=x-P^{-1}e^{AL}A^{-1}PGa$ share latent, window
and current input. On the stationary route the current input does suffice (the injection $\Phi(z,u)$ depends on $u(t)$
only; Section~\ref{sec:m1}); the price there is a small-gain margin (Proposition~\ref{prop:m1-where}).
\emph{Tested by} the $L=1$ row of C2-window-a1 (above). No current node compares the context encoders (current, window,
GRU, SFM) at a matched budget; C1-ctx only moves the SFM's parameters into the budget and does not test this
proposition.

% ---------------------------------------------------------------------
\subsection{Predictions for the window sweeps}
With $dt=0.1\sin(\pi/2n_z)/\alpha$, $\alpha\,dt$ does not depend on the bandwidth scale $s$, and a window of $L$
samples covers $(L-1)\,dt$ seconds of relevant history (Proposition~\ref{prop:m2-current}a).
\begin{enumerate}
\item The knee in seconds scales as $1/(s\alpha)$; in \emph{samples} it is the same for all $s$ (vdp, duffing,
battery in C2-window-a2). One e-fold of the rate-$\alpha$ bound is $1/(\alpha dt)=32.4$ samples ($n_z=5$) or
$44.9$ ($n_z=7$); at rate $\alpha_k$ it is $\alpha_k/\alpha$ times fewer ($2.6$ at $n_z=5$, $n_x=2$; $4.0$ for lorenz,
$n_z=7$, $n_x=3$; the \texttt{a2/a} column of \texttt{c2\_predictions} is the second-slowest mode group, which equals
$\alpha_k$ only when $n_z-n_x+1\le4$). The worst-case window of
Theorem~\ref{thm:m2-window}(a) at $\eta=10^{-3}z_{\rm scale}$ ($\ell$ omitted) is $\approx311,313,320,467$ samples;
measured tails are $10^2$--$5\cdot10^2$ times below the bound, i.e.\ the bound overstates the window by
$\ln(485)/(\alpha\,dt)\approx200$ samples. At one bandwidth (C2-window-a1) the knees within $5\%$ of $e(512)$ are $32$ (vdp,
\texttt{id} and \texttt{sequence\_ood}; duffing and bioreactor on \texttt{sequence\_ood}), $64$ (bioreactor \texttt{id}),
$128$ (quad\_tank \texttt{id}) and not reached on the grid for battery and lorenz (both tiers), duffing (\texttt{id}) and
quad\_tank (\texttt{sequence\_ood}); see Table~\ref{tab:m2-window-rule}. At $2\alpha$ (C2-window-a2) the knee in samples
moved by $\le2\times$ on only $1/2$ evaluable systems: the $1/\alpha$ scaling is \emph{not} supported by the current data.
\item quad\_tank has a fixed $dt=0.05$: $\alpha dt=0.0025,0.005,0.01$ for $s=\tfrac12,1,2$; the knee in samples scales as
$1/s$ and the grid ($L\le512$) spans only $1.3,2.6,5.1$ e-folds, so the curve should still fall at $L=512$ for
$s\le1$ (it does at $s=1$): this plant cannot test the scaling.
\item On a dyadic grid ``knee$\times\alpha$ varies by less than $2\times$'' holds only if the knee sits at the same grid
point for every $s$ (neighbours differ by exactly $2$); use ``$\le2\times$'' or a knee interpolated in $\log L$.
\end{enumerate}

% ---------------------------------------------------------------------
\subsection{Context encoders: forgetting by construction}\label{sec:m2-ctx}

\begin{lemma}[SFM context]\label{lem:m2-sfm}
[proven] Let $\xi_{k+1}=A^\xi_d\xi_k+B_du_k/u_{\rm ref}$, $c_k=\psi([\xi_k,u_k/u_{\rm ref}])$, with $A^\xi$ block-normal,
$s_b=-\mathrm{softplus}(\cdot)<0$, $\alpha_\xi=\min_b|s_b|$, and $\psi$ $\ell_\psi$-Lipschitz. Then for $l\ge1$
\[
 \text{(a)}\ \Big\|\frac{\partial c_k}{\partial u_{k-l}}\Big\|\le\frac{\ell_\psi}{u_{\rm ref}}\|B_d\|e^{-\alpha_\xi(l-1)dt},\qquad
 \text{(b)}\ \ell_E\le\frac{\ell_\psi}{u_{\rm ref}}\Big(1+\frac{\|B_d\|}{1-e^{-\alpha_\xi dt}}\Big),
\]
and the per-lag bound is attained ($\psi$ isometric on the slowest block, $B_d$ in that block). (c) (\emph{the
memory probe}) If two input records agree on the last $L$ samples and differ arbitrarily before, with $|u|\le\bar u$,
\[
 \|c_k-c'_k\|\le\frac{\ell_\psi\|B_d\|\,2\bar u}{u_{\rm ref}}\,\frac{e^{-\alpha_\xi L\,dt}}{1-e^{-\alpha_\xi dt}} .
\]
\end{lemma}
\begin{proof}
$\xi_k=\sum_{l\ge1}(A^\xi_d)^{l-1}B_du_{k-l}/u_{\rm ref}$, $\|(A^\xi_d)^{l-1}\|=e^{-\alpha_\xi(l-1)dt}$; chain rule. For $\ell_E$
add the direct path $u_k$ and sum the geometric series. (c) Only lags $l\ge L+1$ differ; sum (a) over them.
\end{proof}
The rate exists for every parameter value but is not bounded below uniformly (softplus$\to0$); a floor
$s_b\le-\alpha_{\min}$ makes it uniform (the C2-floor nodes set $\alpha_\xi\ge c\,\alpha$, $c\in\{0.25,0.5,1\}$).

\emph{What the probe measures.} The C1-where and C2-floor memory probes keep the last $L$ inputs, swap everything
older, and fit a rate to the decoded change $|\Delta\hat x|\le\ell_c\|c-c'\|$. This is the quantity of (c): a test of
the \emph{encoder's} memory, not of Lemma~\ref{lem:m2-tail} (the latent $z$ is unchanged by the probe). Since (c) is an
upper envelope, a fitted rate below $\alpha_\xi$ over a finite $L$ range does not contradict it unless the curve is in
its exponential regime over several decades; the probe tests the rate only there.

\emph{At initialisation} the SFM's fill-in poles are $\alpha/1.5^{b}$, $b=1,\dots,5$, so its slowest mode is
$\alpha/1.5^5=0.132\alpha$ and the next $0.198\alpha$: at initialisation the SFM forgets $\approx7.6\times$ more slowly
than the latent. \emph{Measured (un-floored, C1-where, lora, median over seeds; \texttt{paper/generated/RESULTS.md}):}
the learned SFM's slowest pole is $0.14$--$0.18\,\alpha$ (battery $0.18$, bioreactor $0.14$, duffing $0.14$, lorenz $0.18$,
quad\_tank $0.17$, vdp $0.15$), i.e.\ between the two slowest \emph{initial} poles: training moves the slowest pole
only slightly (by at most $\approx40\%$, never past the second initial pole), so the ``$5$--$7\times$ extra memory'' of the
un-floored SFM is largely inherited from initialisation rather than learned.
\emph{The floor intervention (C2-floor-025/05/1) fails its pre-registered rule.} At $0.5\alpha$: no significant loss on
$4/6$ dev systems (rule: $\ge5/6$; battery $+23$--$36\%$ and bioreactor lose), TOST non-inferiority at $\pm10\%$ on $0/6$,
and the strict probe check ``fitted rate $\ge$ floor'' on $0/6$ (ratios $0.73$--$0.96$: battery $0.78$, bioreactor $0.77$,
duffing $0.73$, quad\_tank $0.96$, vdp $0.82$; lorenz $-0.02$, flat). At $0.25\alpha$: no significant loss on $6/6$ but
TOST $3/6$ and probe $1/6$; at $1\alpha$: $4/6$, $0/6$, $0/6$ (battery $\approx2\times$ worse). So the data are
\emph{consistent with} the memory being unused on four systems but do not establish it (three seeds; see the TOST counts),
and battery, the large-$\ell$ plant of Corollary~\ref{cor:m2-rule}, is the one system where a floor at $0.5\alpha$ clearly
costs accuracy, as the corollary says it should (its sufficient window $\ln(1/\eta)/\alpha+\ln\ell/\alpha$ is attenuated by
$e^{-0.5\alpha L^*}$ under the floor); the bioreactor loss has no such explanation. The probe ratios below one do not
contradict Lemma~\ref{lem:m2-sfm}(c), which is an upper envelope, but they mean the pre-registered check of its
\emph{rate} fails on every system, and the flat lorenz curve ($d(L)$ within $\times2$--$3$ over $L\le5.7$~s) is unexplained:
a probe outside its exponential regime, to be reported as such. For the window prediction the SFM has no truncation loss
on the latent's time scale (it
should match the window encoder's plateau beyond the knee, not its knee), but an input glitch persists $7.6\times$ longer
in $c$ and $\ell_E$ is correspondingly larger ($\|B_d\|/(1-e^{-\alpha_\xi dt})\approx7.6/(\alpha dt)\approx246\|B_d\|$
for $n_z=5$).

\begin{remark}[SFM is an exact context for linear plants]
If $\operatorname{spec}(A^\xi)\supseteq\operatorname{spec}(A)$ with nonzero input weights (true at initialisation), then
$q(t_k)=-\sum_{i<k}A_d^{k-1-i}A^{-1}(A_d-I)PGu_i$ is a fixed linear function of $\xi_k$ (each $2\times2$ block commutes
with $J$): no truncation, whereas a window of $L$ samples truncates at $e^{-\alpha(L-1)dt}$. Once the poles are
trained this holds only approximately.
\end{remark}

\begin{remark}[a GRU has no such guarantee]
(1) An update-gate bias $b$ gives $h_k=\sigma(b)h_{k-1}+\dots$, forgetting at rate $-\ln\sigma(b)\to0$ as $b\to\infty$.
(2) No fading memory in the sense of \citet{boyd1985fading}: with reset gate $\equiv1$ and update gate $\equiv0$,
$h_{k+1}=\tanh(3h_k+2u_k)$ is bistable; one pulse flips the state and the context difference stays $1.99$ for all
later $k$, although the two inputs agree after the pulse. Contractivity of recurrent networks needs weight
constraints \citep{miller2019stable}; the SFM has it by construction. \emph{Tested by} no current node (an encoder
ablation SFM vs GRU at a matched budget is requested; C1-ctx does not contain a GRU variant).
\end{remark}

\begin{table}[ht]\centering\footnotesize
\caption{Numerical checks of this section (CPU).}\label{tab:m2-checks}
\begin{tabular}{@{}lp{4.6cm}p{7.4cm}@{}}\hline
result & check & outcome\\\hline
Lemma~\ref{lem:m2-tail} & integrator; oscillator $n_z=5$; Van der Pol & rate $1.5000=\alpha$; $1.0000=\alpha$ (both); bound loose $135$--$485\times$\\
Thm.~\ref{thm:m2-window}(b) & $\delta(L)$ by gauge/LP, oscillator $n_x=2$ & rate $1.000,2.617$ at $n_z=3,5$ ($\alpha_k=1,2.618$); needle at predicted angle\\
 & bounds $r_qe^{-\alpha_kL}/\|P\|\le\delta\le2R_qe^{-\alpha L}/\sigma_{\min}$ & hold at every $L$ (slack $21\times$, $1.5$--$106\times$)\\
 & $\sigma_{\min}(e^{-AL}P)$ (arithmetic) & same rates $1.000,2.618$\\
Cor.~\ref{cor:m2-rule} & Table~\ref{tab:m2-window-rule} from the C2-window-a1 report & $\ell$-free rule fails $1.73\times$ (battery); with $\ln\ell$: holds on $2/3$ testable systems (lorenz $1.07$), untestable on battery, duffing, quad\_tank\\
Prop.~\ref{prop:m2-current}(c) & Van der Pol, 12 past levels & $\operatorname{rank}\partial_\theta\Psi=5$\\
sampled form & ZOH step identity; discrete constant & exact to $10^{-15}$; $c_d=0.985\le1/\alpha$\\
Lemma~\ref{lem:m2-sfm}(a,b) & tight SFM; random SFM & ratio $1.0000000000$; $\le0.13$\\
Lemma~\ref{lem:m2-sfm}(c) & history swap, floor $0$ and $0.5$; initial poles & observed/bound $\le0.01$; $\alpha_\xi=\alpha/1.5^5$ at init, $\ge0.5\alpha$ with floor\\
GRU remark & latch in the code's GRU encoder & difference $1.990$ at $k=100,1000,4999$\\\hline
\end{tabular}
\end{table}
% scripts: theory/checks/check_theory.py (integrator, GRU gate bias), theory/checks/m2_checks.py (all other rows)

% Bibliography entries used:
% bernard2019luenberger: P. Bernard, V. Andrieu. Luenberger observers for nonautonomous nonlinear systems.
%   IEEE Trans. Automatic Control 64(1):270--281, 2019.
% boyd1985fading: S. Boyd, L. O. Chua. Fading memory and the problem of approximating nonlinear operators
%   with Volterra series. IEEE Trans. Circuits and Systems 32(11):1150--1161, 1985.
% miller2019stable: J. Miller, M. Hardt. Stable recurrent models. ICLR 2019.

% =====================================================================================
% Block M3 (revision 3): error bound for the deployed decoder-side observer, sensor faults, MODEL
% INCOMPLETENESS (prop:m3-mismatch, the C3 anchor), measurement, and Lemma D (the single sampled-data
% lemma; block M2 cites its part (e)). The ONLY \appendix of the document is issued here, before Lemma D.
% No preamble; amsthm environments; natbib keys (entries at the end of this file).
% Standing assumptions: the shared block of Assumption~\ref{ass:m1}, referred to as
% (A1)-(A4) plus its regularity clause; this section adds only (L).
% Checks: theory/checks/m3_checks.py (CPU, ~2 min).
% =====================================================================================

\section{Error bound for the deployed observer}\label{sec:m3}

\paragraph{Setting.} Plant $\dot x=f(x,u)$, $y=h(x)$, input held on $[t_k,t_{k+1})$, $t_k=k\,dt$.
The decoder-side observer, exactly as deployed, is
\begin{equation}\label{eq:m3-obs}
\hat z_{k+1}=A_d\hat z_k+W_0\tilde y_k+W_1\tilde y_{k+1},\qquad
\hat x_k=\hat T^{+}(\hat z_k;\hat c_k),\qquad \hat c_k=E(\tilde u_0,\dots,\tilde u_k),
\end{equation}
with $A_d,W_0,W_1$ the exact step of Lemma~\ref{lem:D}(a), the context held between samples, and
$\hat z_0$ a fixed function of the prior $\hat x_0$. Its response to a sensor sequence $(v_k)$ is
the \emph{latent noise filter}, the sampled form of $(sI-A)^{-1}B$:
\begin{equation}\label{eq:m3-H}
H_0=W_1,\qquad H_m=A_d^{m-1}W_0+A_d^{m}W_1\ (m\ge1),\qquad
\gamma_y:=\sum_{m\ge0}\|H_m\|\le\frac{\kappa\|B\|}{\alpha},\qquad \|H\|_2^2:=\sum_{m\ge0}\|H_m\|_F^2 ,
\end{equation}
where the inequality holds at every $dt$ (Lemma~\ref{lem:D}(c)).

\paragraph{Reference condition.} Every bound is stated relative to the \emph{nominal} deployed
observer $\hat x^{\rm nom}_k$: the same episode, the nominal sensor model (the one used in training,
\texttt{data.sensor}), the true input record and the nominal prior. Deviations are
$v_k:=\tilde y_k-\tilde y^{\rm nom}_k$ (sensor faults, extra noise), $\hat c_k-c^{\rm nom}_k$ (input
corruption) and $\hat z_0-\hat z^{\rm nom}_0$ (initialisation). This is the condition the decoder
was trained in: phase~2 regresses $x_k$ on the deployed latent of the nominal-noise record, so its
validation error measures the nominal error directly, and no noise-free quantity is needed.

\paragraph{Assumptions.} (A1)--(A4) and the regularity clause ($f(\cdot,u)\in C^1$, $h\in C^2$,
hence $M_2:=\sup\|\ddot y\|$, $L_y:=\sup\|\dot y\|$ on each hold interval) of
Assumption~\ref{ass:m1}, and
\begin{enumerate}
\item[(L)] $\hat T^{+}\in C^1$, and the context encoder satisfies
$\|E(\tilde u_{0:k})-E(u_{0:k})\|\le\ell_E\sup_{j\le k}\|\tilde u_j-u_j\|$ (for SFM,
$\ell_E\le(\ell_\psi/u_{\rm ref})(1+\|B_d\|/(1-e^{-\alpha_\xi dt}))$, Lemma~\ref{lem:m2-sfm}).
\end{enumerate}
Write $\ell_z,\ell_c$ for the sup of $\|\partial_z\hat T^{+}\|$, $\|\partial_c\hat T^{+}\|$ over the
segments a given episode traverses (defined in Proposition~\ref{prop:m3-iss}), and
$\ell^{\rm g}_z,\ell^{\rm g}_c$ for the global constants (finite for tanh/SiLU networks on a bounded
context set, for every conditioner; for FiLM, LoRA and hyper the available bound is a product of layer norms that is
orders of magnitude larger, Proposition~\ref{prop:m1-where}(b)). The latent deviations are linear,
input-free and decoder-independent, so they are fixed before any Lipschitz constant is used.

\begin{definition}[nominal error, held-out quantile]\label{def:m3-fit}
For an episode $\omega\sim\mathcal D$ (initial state and input path from the training episode
distribution) let $e_{\rm nom}(\omega):=\sup_{k\in S}\|\hat x^{\rm nom}_k-x(t_k)\|$, $S$ the steady
window, and let $\varepsilon^\beta$ be its $(1-\beta)$-quantile. For $n$ i.i.d.\ held-out episodes,
$\hat\varepsilon:=\max_ie_{\rm nom}(\omega_i)$ satisfies
\begin{equation}\label{eq:m3-quantile}
\mathbb P\big(\hat\varepsilon\ge\varepsilon^\beta\big)\ge1-(1-\beta)^n\qquad(n=64,\ \beta=0.05:\ 0.96),
\end{equation}
because all $n$ draws fall below the quantile with probability at most $(1-\beta)^n$.
\end{definition}

$\varepsilon^\beta$ contains everything the nominal observer gets wrong: approximation, window
truncation (by Lemma~\ref{lem:m2-tail}(a) and Theorem~\ref{thm:m2-window}(a) at most $\tfrac{\ell}{2}\varepsilon(L)$
for a sufficient context), nominal noise and the discretisation (Lemma~\ref{lem:D}(b) only relates the
deployed latent to the continuous map, $\|\bar z_k-z(t_k)\|\le e_{\rm disc}=O(dt^2)$). It is a
quantile with distribution-free confidence, not a supremum, and only for $\mathcal D$. It is defined on the
\emph{true} plant's episodes (phase~2 regresses true-plant state labels on the deployed latent): no model of $f$
enters it, which is the point of Proposition~\ref{prop:m3-mismatch}.

% -------------------------------------------------------------------------------------
\begin{proposition}[decoder-side observer: ISS relative to the nominal observer]\label{prop:m3-iss}
[(i), (ii), (iv) proven and exact; (iii) proven as a decomposition, see the paragraph after the proof]
Assume (A1)--(A4), (L). Let $\Delta z_k:=\hat z_k-\hat z^{\rm nom}_k$, $\Delta c_k:=\hat c_k-c_k$.
\begin{enumerate}[label=(\roman*)]
\item (latent deviation, exact) $\Delta z_k=A_d^k\Delta z_0+\sum_{i\le k}G^{(k)}_iv_i$ with
$G^{(k)}_i=H_{k-i}$ ($i\ge1$), $G^{(k)}_0=A_d^{k-1}W_0$; hence
$\|\Delta z_k\|\le\kappa e^{-\alpha t_k}\|\Delta z_0\|+\frac{\kappa\|B\|}{\alpha}\sup_{i\le k}\|v_i\|$ for every $k$,
and $\limsup_k\|\Delta z_k\|\le\gamma_y\limsup_k\|v_k\|$.
\item (path identity, exact)
\begin{equation}\label{eq:m3-identity}
\hat x_k-x(t_k)=\big(\hat x^{\rm nom}_k-x(t_k)\big)+\bar J_{z,k}\Delta z_k+\bar J_{c,k}\Delta c_k,
\end{equation}
$\bar J_{z,k}=\int_0^1\partial_z\hat T^{+}(\hat z^{\rm nom}_k+s\Delta z_k;\hat c_k)ds$,
$\bar J_{c,k}=\int_0^1\partial_c\hat T^{+}(\hat z^{\rm nom}_k;c_k+s\Delta c_k)ds$.
\item (ISS bound) If $\|v_k\|\le\delta_y$ and $\|\tilde u_k-u_k\|\le\delta_u$, then for every $k$
\begin{equation}\label{eq:m3-iss}
\|\hat x_k-x(t_k)\|\le\|\hat x^{\rm nom}_k-x(t_k)\|+\ell_z\Big(\kappa e^{-\alpha t_k}\|\Delta z_0\|+\frac{\kappa\|B\|}{\alpha}\delta_y\Big)+\ell_c\,\ell_E\,\delta_u ,
\end{equation}
with $\ell_z,\ell_c$ the sup of the Jacobian norms on the traversed segments
$[\hat z^{\rm nom}_k,\hat z_k]\times\{\hat c_k\}$ and $\{\hat z^{\rm nom}_k\}\times[c_k,\hat c_k]$; on the
steady window the first term is $\le\varepsilon^\beta$ for a $(1-\beta)$ fraction of episodes, and in
the limsup $\kappa\|B\|/\alpha$ may be replaced by the exact $\gamma_y$.
\item (stochastic deviations) For zero-mean i.i.d.\ $v_k$ with covariance $\sigma_v^2I$ (e.g.\ extra
Gaussian noise), $\limsup_k(\mathbb E\|\hat x_k-x(t_k)\|^2)^{1/2}\le\limsup_k(\mathbb E\|\hat x^{\rm nom}_k-x(t_k)\|^2)^{1/2}+\ell^{\rm g}_z\sigma_v\|H\|_2+\ell_c\ell_E\delta_u$,
$\|H\|_2^2=dt\,\mathrm{tr}W_c+o(dt)$, $AW_c+W_cA^\top+BB^\top=0$. The global constant is needed because
unbounded noise leaves every neighbourhood of the data with positive probability.
\end{enumerate}
\end{proposition}

\emph{Intuition.} The learned map never sees raw measurements: sensor deviations and
initialisation reach it through a known, input-free stable filter, a wrong input record moves
only the context, and each deviation is multiplied by the decoder's slope along the path it
actually travels.

\begin{proof}
(i) $\hat z$ and $\hat z^{\rm nom}$ obey the same linear recursion with forcing difference
$W_0v_k+W_1v_{k+1}$; unroll it, and use $\|A_d^k\|\le\kappa e^{-\alpha t_k}$ and
$\sum_i\|G^{(k)}_i\|\le\kappa\|B\|/\alpha$ (Lemma~\ref{lem:D}(c)). (ii) Split
$\hat T^{+}(\hat z;\hat c)-\hat T^{+}(\hat z^{\rm nom};c)$ at $\hat T^{+}(\hat z^{\rm nom};\hat c)$ and apply the
fundamental theorem of calculus on each segment. (iii) Norms in \eqref{eq:m3-identity}, (i), (L),
Definition~\ref{def:m3-fit}. (iv) $L^2(\mathbb P)$ norms of the pointwise inequality (Minkowski) with
$\|\bar J_z\|\le\ell^{\rm g}_z$; $\mathbb E\|\sum_iG^{(k)}_iv_i\|^2=\sigma_v^2\sum_i\|G^{(k)}_i\|_F^2\to\sigma_v^2\|H\|_2^2$,
a Riemann sum ($H_m=e^{At_m}B\,dt+O(dt^2)$) of $dt\int_0^\infty\|e^{As}B\|_F^2ds$.
\end{proof}

When the nominal sensor is noise-free, \eqref{eq:m3-iss} is the classical decomposition: fit
error plus decoder slope times the ISS gain $\kappa\|B\|/\alpha$ of the latent filter
\citep[cf.][Lemma~17, Theorem~19]{niazi2025kkl}, here for the sampled observer and with an
input-corruption term. The latent error equation of \citet[Proposition~1]{shaaban2026hyperkkl}
assumes no injection; with injection $\Phi\neq0$ the latent deviation is no longer input-free and
the margin of Proposition~\ref{prop:m1-where} is needed (sampled form: Lemma~\ref{lem:D}(d)).
The bound has no amplitude-dependent condition, but its constants $\ell_z,\ell_c$ may depend on the
context.

\emph{(iii) is a decomposition with measured constants, not a certificate.} With $\ell_z,\ell_c$ the suprema over the
traversed segments, (iii) is (ii) with norms taken, and the constants are only known after the fact (lower bounds of the
suprema when estimated from samples). Evaluated \emph{a priori}, with the phase-1 proxy \texttt{l\_dec\_q95} for the global
constant $\ell^{\rm g}_z$ and the nominal sensor level $\delta_y=0.03\,y_{\rm std}$, the second term alone is
$\ell_z\kappa\|B\|\delta_y/\alpha=66$ (battery; $\kappa\|B\|/\alpha=2.68$, $\ell=816$), $15$ (quad\_tank; $37.4$, $13.5$),
$6.7$ (duffing), $5.9$ (bioreactor), $1.1$ (vdp), $0.40$ (lorenz), against measured nominal NRMSE $0.016$--$0.052$
(C1-where, lora, \texttt{id}): $15\times$ to $1400\times$ too large (\texttt{check\_apriori\_vacuity}). The bound is therefore
vacuous as a guarantee on every benchmark and is \emph{not} presented as one. What it offers is the exact attribution
(ii) and the exact stochastic form (iv), both of which are measurable (Definition~\ref{def:m3-measured}).

\begin{remark}[reading the bound]\label{rem:m3-knobs}
Three quantities set the error: the nominal error $\varepsilon^\beta$ (capacity, data, window),
the decoder slope $\ell_z$ times a design-only latent deviation ($\propto1/\alpha$), and the input
sensitivity $\ell_c\ell_E$, which is zero for the autonomous observer. The architecture enters only
through these constants: the bound says which one a conditioning mechanism must keep small, not
which mechanism wins.
\end{remark}

\emph{What would test it (eval-only; implemented by no current node).} Under \texttt{input\_noise/bias/delay} the
excess error is $\bar J_c\Delta c$ exactly: the autonomous observer is unaffected, and variants would order by their
realised input gain $\|\bar J_c\Delta c\|/\|\Delta c\|$ if that gain were measured on saved checkpoints
(Definition~\ref{def:m3-measured}; C3-faults has the conditions, no node records the gains). Under \texttt{poor\_init}
the excess is $\bar J_zA_d^k\Delta z_0$, gone after a few $1/\alpha$ (C3-faults records the convergence time, which is
the observable consequence). Upper bounds do not rank errors, so none of this is implied by the bound.

\emph{Check} (\texttt{m3\_checks.py::check\_noise\_iss}, \texttt{check\_input\_corruption},
\texttt{check\_identity}; linear plant with exact maps $z=Px+\nu$, decoder $P^{+}(z-c)$,
$\ell_z=\ell_c=12.48$, $\ell_E=0.048$, $n_z=5$, $\alpha=1.5$, $dt=0.0206$ as for vdp).
Worst-case bounded deviation: latent $0.0372\le\gamma_y\delta_y=0.0537\le\kappa\|B\|\delta_y/\alpha=0.0745$.
White deviation, RMS latent $0.00427$ (Monte Carlo) vs $\sigma_v\|H\|_2=0.00430$; the sup form at
$\delta_y=4\sigma_v$ is $30\times$ looser, so only (iv) is informative for Gaussian conditions. Input
bias $0.1$ / delay $5$ / bounded noise: state error $0.022\le0.061$, $0.197\le1.21$, $0.024\le0.18$.
Identity \eqref{eq:m3-identity} for a tanh-MLP decoder holds to $4\times10^{-10}$ (finite-difference Jacobians); its realised gain differs from the on-data
$\|\partial_z\hat T^{+}\|_2$ by a factor $0.01$--$1.38$ (median $0.32$), so on-data slopes neither
bound nor estimate it. Coverage of \eqref{eq:m3-quantile} at $n=64$, $\beta=0.05$: $0.961$ vs $0.963$.

% -------------------------------------------------------------------------------------
\begin{corollary}[sensor faults: bounded and forgotten, uniformly in the state]\label{cor:m3-faults}
For a decoder-side observer and any sensor deviation $v$, pathwise
$\|\hat x_k-x(t_k)\|\le\|\hat x^{\rm nom}_k-x(t_k)\|+\ell^{\rm g}_z\|\sum_iG^{(k)}_iv_i\|$. In particular:
an outlier of size $M$ at $k_0$ adds at most
$\ell^{\rm g}_z\kappa\|B\|M\,\big(w_1+e^{\alpha dt}w_0\big)e^{-\alpha(t_k-t_{k_0})}\le\ell^{\rm g}_z\kappa\|B\|M\,dt\,e^{\alpha dt}e^{-\alpha(t_k-t_{k_0})}$
($w_0,w_1$ of Lemma~\ref{lem:D}(c), $w_0+w_1\le dt$), with time integral at most
$\ell^{\rm g}_z\kappa\|B\|M\,dt/\alpha$, for every state, input and $M$; an outlier process (each
sample with probability $p$, size $M$, random sign) adds at most $\ell^{\rm g}_z\gamma_ypM$ in mean and
$\ell^{\rm g}_z\sqrt pM\|H\|_2$ in RMS; sample-and-hold dropout adds at most
$\ell^{\rm g}_z\gamma_yL_y\,dt\,p/(1-p)$ in mean; a drift $b$ adds at most $\ell^{\rm g}_z\gamma_y\sup\|b\|$,
and for slow drift $\Delta z_k\approx-A^{-1}Bb_k$, since the discrete DC gain is
$(I-A_d)^{-1}(W_0+W_1)=-A^{-1}B$ exactly; coloured deviations obey
$\mathbb E\|\Delta z\|^2=\sum_{m,m'}\mathrm{tr}(H_mR_v(m'-m)H_{m'}^\top)$.
\end{corollary}

\emph{Intuition.} There is no feedback from the estimate to the latent: a fault is a forcing of a
fixed stable filter, so the observer's own error can never amplify it.

\begin{proof}
Proposition~\ref{prop:m3-iss}(ii) with $\Delta c=0$, $\Delta z_0=0$, $\|\bar J_z\|\le\ell^{\rm g}_z$. Outlier:
$\Delta z_{k_0}=W_1Me$ and $\Delta z_k=(A_d^{k-k_0-1}W_0+A_d^{k-k_0}W_1)Me$, bounded with the merged kernels of
Lemma~\ref{lem:D}(c); the integral is $dt\sum_m\|H_m\|M\le dt\,\kappa\|B\|M/\alpha$. Process:
$\mathbb E\|v_i\|\le pM$ in (i), and $R_v(m)=pM^2\delta_m$. Dropout: a held sample differs from the
current one by at most $L_y\,dt\,N_k$, $N_k$ geometric with $\mathbb EN_k\le p/(1-p)$. DC gain:
$W_0+W_1=A^{-1}(A_d-I)B$ and $A,A_d$ commute; the quasi-static error is bounded by summation by parts,
$\|\Delta z_k+A^{-1}Bb_k\|\le\sum_{m\ge1}\|\sum_{j\ge m}H_j\|\,\sup_j\|b_j-b_{j-1}\|$.
\end{proof}

\emph{What this proves about filters, and what it does not.} The corollary is a \emph{qualitative} ISS guarantee: the
excess is bounded and decays at rate $\alpha$ for every state, input and fault size, given $\ell^{\rm g}_z$, with no
feedback from the estimate to the latent. The EKF's stochastic stability is local (small initial error and noise,
\citealp{reif1999stochastic}) and an outlier moves its linearisation point. A qualitative guarantee is not a comparative
advantage, and the data say so: on \emph{exact} models no C3-faults condition (sensor, input, plant-drift and outage)
reaches the learned-wins regime of $\ge4/6$ systems (learned-wins regime: \CthreeFaultsRegimes; the learned observer
beats the best filter on $0/6$ under every sensor-side condition and on at most $2/6$ against the best filter of either
kind under input bias, input delay or $10\%$ plant drift), and it wins on $0/6$ under colored Student-$t$ noise
(C3-noise-ct), as Remark~\ref{rem:m3-kf} predicts for linear plants and as the EKF's
\ConeLearnedOverEkfRangeSeq$\times$ (median \ConeLearnedOverEkfMedianSeq$\times$) smaller exact-model
\texttt{sequence\_ood} error on the dev systems (C1-where) makes expected. Worse, the best \emph{learned}
observer is itself above climatology (a failed estimator) under exact-model \texttt{outliers} (battery, duffing,
quad\_tank) and under $30\%$ post-training plant drift (bioreactor, quad\_tank): the corollary bounds the excess by
$\ell^{\rm g}_z\times$(filter response), and with $\ell^{\rm g}_z$ in the hundreds that bound is consistent with failure.
Where learned observers win is model incompleteness (C3-mismatch10/20/30/50, C3-partial, C3-realistic), and the
theoretical lever there is not robustness to sensor faults but the model term of Proposition~\ref{prop:m3-mismatch}.

\emph{Exact filter-level ratios (design-invariant).} For the Butterworth design with $dt=0.1/\|A\|$,
the \emph{latent} deviations of the stress conditions, relative to the nominal noise
($0.03\,y_{\rm std}$), do not depend on the system: RMS $\times3.4$ for \texttt{colored}
($\varrho=0.95$), $\times57$ for \texttt{outliers} ($p=0.02$, $M=12\,y_{\rm std}$), and $\times30$--$32$
at the quasi-static peak for \texttt{sensor\_drift} ($0.3\,y_{\rm std}$); \texttt{dropout} costs
$O(L_ydt)$ per held sample. These are exact properties of the known filter; the \emph{state}-level excess multiplies
them by the variant's realised $\|\bar J_z\|$, which no current node records (Definition~\ref{def:m3-measured}).
\emph{Check} \texttt{m3\_checks.py::check\_fault\_ratios} computes these ratios for vdp, duffing,
bioreactor and lorenz (all agree to two digits; drift $\times32.1,32.1,32.1,30.5$).
\texttt{check\_outliers}, \texttt{check\_dropout\_drift}: outlier peak $0.519\le0.536$, envelope never
violated, integral $0.266\le0.369$; process mean $0.174\le0.258$, RMS $0.241$ vs $0.243$; dropout
($p=0.1$) mean $0.00105\le0.00737$; DC identity to $10^{-16}$; random-walk drift peak $0.137\le0.322$,
quasi-static deviation $0.0052\le0.049$.

\begin{remark}[Kalman filter vs.\ learned KKL on linear plants]\label{rem:m3-kf}
For $\dot x=Fx+Gu$, $y=Cx$ the exact maps are $z=Px+\nu$ with $PF=AP+BC$, $\dot\nu=A\nu-PGu$ and
decoder $P^{+}(z-\nu)$, so the KKL estimate is an LTI filter of $(y,u)$ with sensor-to-error response
$\Gamma_m=P^{+}H_m$; for $n_z=n_x$ it is the Luenberger observer with gain $P^{-1}B$. The steady KF
has the same structure with gain $K$. (1) Gaussian noise with correct $(Q,R)$: the KF is the MMSE
estimator, so no KKL design beats it. (2) White non-Gaussian noise (Gaussian plus outliers): every
LTI filter has steady MSE $\mathrm{MSE}_w+(\sigma_v^2+pM^2)\|\Gamma\|^2_{H_2}$, the KF with
$R=\sigma_v^2+pM^2$ is the best linear filter \citep{andersonmoore1979}, and the KKL observer against a
KF tuned to the nominal $R$ is decided by $\|\Gamma\|_{H_2}$ alone. \emph{Check}
(\texttt{check\_kf\_comparison}; oscillator, $\sigma_v=0.03$, $p=0.02$, $M=12$): Gaussian MSE KF
$1.98\times10^{-3}<$ KKL $6.75\times10^{-3}$; with outliers KF(nominal $R$) $0.514$, KF($R_{\rm eff}$)
$0.0115$, KKL $2.40$, each within $2\%$ of the formula. The KKL loses to both here ($\|P^{+}\|=12.5$):
outlier robustness is not a general property of KKL observers, and on the benchmarks any advantage
over the EKF must be attributed to its mis-tuning, local stability and linearisation, or, as the data show,
to the filter's model being wrong (next).
\end{remark}

% -------------------------------------------------------------------------------------
\subsection{Model incompleteness: the term that decides C3}

\begin{proposition}[persistent model error: filter bias floor vs.\ learned fit error; linear plants]\label{prop:m3-mismatch}
[proven for linear plants] Let the true plant be $x_{k+1}=Fx_k+Gu_k+d_k+w_k$, $y_k=Cx_k+v_k$, where $d_k$ is the part
of the dynamics the modeller's model lacks (a wrong $F$ or a missing term: $d_k=\Delta F\,x_k+\Delta G\,u_k$), $w_k,v_k$
white with covariances $Q,R$. Any fixed-gain predict--correct observer built on the model $(F,G,C)$, in particular the
steady-state Kalman filter and every Luenberger observer, $\hat x_{k+1}=\hat x^-_{k+1}+K(y_{k+1}-C\hat x^-_{k+1})$,
$\hat x^-_{k+1}=F\hat x_k+Gu_k$, has error $e_k=x_k-\hat x_k$ obeying
\[
e_{k+1}=(I-KC)F\,e_k+(I-KC)(d_k+w_k)-Kv_{k+1}.
\]
\begin{enumerate}[label=(\roman*)]
\item (\emph{bias}) For a constant $d_k\equiv\bar d$ (e.g.\ a constant input and $\Delta F\bar x$ at the steady state), the
error converges to the non-decaying bias $\bar e=\big(I-(I-KC)F\big)^{-1}(I-KC)\,\bar d$; in continuous time, for
$\dot e=(F-KC)e+\Delta F\,x$, $\bar e=-(F-KC)^{-1}\Delta F\,\bar x$, the DC value of $(sI-(F-KC))^{-1}\Delta F\,x$.
\item (\emph{bias--variance floor}) Let $Q\succ0$, $R\succ0$. For a stationary $d_k$ independent of $(w,v)$ (e.g.\ $d$
evaluated along the noise-free trajectory, as in the check; with $d_k=\Delta F\,x_k$ along the noisy trajectory the part
$\Delta F\,x^{\rm sto}_k$ is an extra process noise correlated with $w$, and $\Sigma$ must then be taken from the joint
Lyapunov equation of $(x^{\rm sto},e)$) the steady mean-square error splits exactly as
$\mathrm{MSE}(K)=b^2(K)+\mathrm{tr}\,\Sigma(K)$, with $b^2$ the mean square of the deterministic response to $d$
through $\Phi_K=(I-KC)F$ and $\Sigma=\Phi_K\Sigma\Phi_K^\top+(I-KC)Q(I-KC)^\top+KRK^\top$. Inflating the process
covariance to $Q+\lambda I$ (the oracle filters' \texttt{q\_inflate}) moves $K(\lambda)$ along a one-parameter curve
that trades $b^2$ against $\mathrm{tr}\,\Sigma$ (in the check: $\|K\|$ and $\mathrm{tr}\,\Sigma$ increase, $b^2$ decreases,
monotonically); the best achievable error $\mathrm{MSE}^\star:=\inf_K\mathrm{MSE}(K)\ge\inf_K\mathrm{tr}\,\Sigma(K)>0$
depends on $(F,G,C,\Delta F,\Delta G,Q,R)$ and the input, and on no data. Only $n_y$ directions of $d$ can be cancelled
by $K$, so $b^2$ keeps a positive limit as $\lambda\to\infty$ whenever $d$ has a component $K$ cannot cancel (e.g.\
$C\bar d=0$, as for the missing stiffness term of the check).
\item (\emph{learned decoder-side observer}) Its error obeys Proposition~\ref{prop:m3-iss} with $\varepsilon^\beta$
measured on the true plant: $\|\hat x_k-x_k\|\le\varepsilon^\beta+\ell_z\|\Delta z_k\|+\ell_c\|\Delta c_k\|$ on a
$(1-\beta)$ fraction of episodes, and no term depends on the model. In the linear case, a decoder fitted on true-plant
data is $P_{\rm true}^{+}(z-\nu_{\rm true})$ with $P_{\rm true}F_{\rm true}=AP_{\rm true}+BC$ (continuous-time Sylvester
equation of the true plant; for the sampled plant its discrete counterpart): its error is driven by $w,v$ only (no bias),
with a mean-square error $\mathrm{MSE}_{\rm KKL}$ that is a property of the true plant and the noise; at a fixed true
plant it does not depend on the modeller's model at all. The model enters the
learned observer only through the phase-1 PDE regulariser of $T_0$ (and through \texttt{kkl\_mb}, which is a model-based
filter and belongs to the other class).
\end{enumerate}
\emph{Prediction.} The learned observer beats the best (oracle-tuned) filter iff $\mathrm{MSE}^\star>\varepsilon^2$,
where $\varepsilon$ is its fit error on the true plant. For a missing term $s\,d$ of fixed shape and scale $s$ (at a fixed
true plant), $\mathrm{MSE}^\star(s)=\inf_K[s^2b_1^2(K)+\mathrm{tr}\,\Sigma(K)]$ ($b_1^2$ the bias term at $s=1$; $d$ independent of $(w,v)$ as in (ii)) is nondecreasing in $s$ and $\varepsilon$ does
not depend on $s$, so there is at most one crossing: filters win below a threshold in the model error (possibly $0$) and
learned observers above it. The threshold depends on the plant and the noise.
\end{proposition}
\emph{Intuition.} A model-based filter corrects its wrong model only through the $n_y$ measured directions and only
at the cost of trusting the noisy measurement more; a decoder trained on the true plant has no model to be wrong about,
it pays its fit error instead.
\begin{proof}
Subtract the predict--correct recursion from the true one: $x_{k+1}-\hat x^-_{k+1}=Fe_k+d_k+w_k$ and
$e_{k+1}=(I-KC)(x_{k+1}-\hat x^-_{k+1})-Kv_{k+1}$. (i) Set $e_{k+1}=e_k$ with $w=v=0$; $(I-(I-KC)F)$ is invertible when
the filter is stable. Continuous time: set $\dot e=0$. (ii) Linearity: the response to $d$ (deterministic) and to
$(w,v)$ (zero mean, independent of $d$) are orthogonal in $L^2$, so the mean squares add; $\Sigma$ is the stationary
covariance of the noise response (discrete Lyapunov equation). $\mathrm{tr}\,\Sigma(K)$ is the error of the linear filter
$K$ under the nominal model and the true noises, so it is bounded below by the nominal-model MMSE, which is positive for
$Q\succ0$, $R\succ0$; hence $\mathrm{MSE}^\star>0$ for every $d$. As $\lambda\to\infty$ the gain can cancel $d$ only on the $n_y$
directions that $C$ measures, so $b^2$ has a positive limit when $(I-KC)\bar d$ cannot be made zero. (iii)
Proposition~\ref{prop:m3-iss} with
Definition~\ref{def:m3-fit}; in the linear case $z-P_{\rm true}x-\nu_{\rm true}$ obeys a stable linear recursion forced
by $W_0v_k+W_1v_{k+1}-P_{\rm true}w_k$ only.
\end{proof}
\emph{Check} (\texttt{m3\_checks.py::check\_model\_mismatch}; damped oscillator, missing stiffness term
$\Delta F=\epsilon\,[0\ 0;-1\ 0]$, $q=0.05$, $\sigma_v=0.03$, PRBS input). (i) the steady bias agrees with the formula to
$10^{-8}$ (discrete, $\lambda=0$ and $10^{-3}$) and $10^{-6}$ (continuous Luenberger). (ii) at $\epsilon=0.3$
(relative model error $|\Delta Fx|/|Fx|=0.29$), over $\lambda\in[0,10^{-1}]$: $\|K\|$ rises $0.25\to1.19$, $b^2$ falls
$5.24\to4.80\cdot10^{-2}$, $\mathrm{tr}\,\Sigma$ rises $2.0\to3.0\cdot10^{-3}$, $\mathrm{MSE}$ is U-shaped with
$\mathrm{MSE}^\star=5.09\cdot10^{-2}$; Monte Carlo agrees with $b^2+\mathrm{tr}\,\Sigma$ within $1\%$. (iii) the
true-plant KKL has $\mathrm{MSE}=6.6\cdot10^{-3}$ and mean error $10^{-3}$ at every $\epsilon$. Threshold: filters win at
relative model error $0$ and $0.05$ ($\mathrm{MSE}^\star=2.0,4.2\cdot10^{-3}$), learned wins from $0.10$ on
($10.0\cdot10^{-3}$ and above), one crossing.

{\sloppy\emph{What it predicts, and where it is tested} (verdicts from \texttt{paper/generated/RESULTS.md}, the mechanical
application of each node's pre-registered rule; model errors are the generated macros). With model error the C3 nodes
measure, the proposition predicts a single crossing in the relative model error: filters win below it, learned observers
above, per plant. On the exact-model side no condition reaches a learned-wins regime (C3-faults: \CthreeFaultsRegimes;
C3-noise-ct $0/6$). On the
incomplete-model side the learned observer beats the best \emph{working} filter (a filter above climatology counts as
failed) at every mismatch level, $10\%$ ($5/6$ \texttt{id}, $4/6$ \texttt{sequence\_ood}), $20\%$ ($5/6$, $5/6$), $30\%$
($5/6$, $5/6$), $50\%$ ($6/6$, $5/6$), under a missing term (C3-partial $5/6$, $5/6$) and in the realistic node
(\CthreeRealisticIdVsWorking, \CthreeRealisticSeqVsWorking); the pre-registered crossover hypothesis ``between $10\%$ and
$30\%$'' therefore \textbf{fails} in the direction of the learned observer (first winning level $10\%$; the crossing lies
below the smallest tested mismatch for the $\ge4/6$ level rule, although duffing and battery are non-wins at $10\%$).
Across the $36$ model-wrong cells (\texttt{sequence\_ood}, vs.\ the best working filter) the learned observer wins in
$30/36$. \emph{Within} each system the outcome is ordered by model error: the six non-wins are duffing under every
parameter mismatch (model error \ModelErrDuffingMismatchThirty\ at $30\%$), while duffing under the missing term
(\ModelErrDuffingPartial) is a win; battery at $10\%$ mismatch, while battery wins from $20\%$ (model error $0.064$); and
quad\_tank under the missing term (\ModelErrQuadTankPartial), while quad\_tank wins under parameter mismatch
(\ModelErrQuadTankMismatchThirty\ at $30\%$); the systems with the largest errors (bioreactor \ModelErrBioreactorPartial,
lorenz \ModelErrLorenzPartial, vdp \ModelErrVdpPartial) are won at every level. \emph{Across} systems the ordering does
\emph{not} hold: the smallest model error with a win ($0.064$, mismatch20/battery) is below the largest without one
($0.152$, partial/quad\_tank). This is what the proposition allows (the threshold depends on the plant and the noise;
the toy crosses at $0.05$--$0.10$): the data are consistent with one crossing per plant, at plant-dependent thresholds
that overlap in the band $0.06$--$0.15$, not with a single cross-system threshold. This reading is post hoc, not a
pre-registered test; TEST-C3 (frozen recipe, held-out plants) is \textbf{inconclusive} at the time of writing ($6/18$ jobs). The proposition is
proven for linear plants; on the nonlinear benchmarks ``learned wins iff the floor exceeds the fit error'' is the reading
of an \textbf{empirical} result. The ``no model term'' half of (iii) has one caveat: the phase-1 PDE term does use the
(wrong) model; whether that helps or hurts is \emph{not} established (C3-partial-modelfree: the partial model beats a
model-free phase~1 with a CI excluding zero on $0/6$ systems, so no claim that the partial model helps is made). Two
further comparisons are needed before C3 can be read as a statement about KKL rather than about ``data beats a wrong
model'': a Koopman/EDMDc Kalman observer and a grey-box (partial model plus learned residual) EKF trained on the same data
(\texttt{baselines\_data.py}; not yet in the reported nodes).\par}

% -------------------------------------------------------------------------------------
\subsection{Measurement and certification}

\begin{definition}[measured decomposition and predicted bound]\label{def:m3-measured}
For a trained model, a stress condition and the steady window, in the $\sigma$-normalised norm of
\texttt{metrics.py}, compute exactly: $\Delta z_k$ by running the known filter on $v=\tilde y-\tilde y^{\rm nom}$
(and $A_d^k\Delta z_0$ for the prior); $\Delta c_k$ by running $E$ on $\tilde u$ and $u$; and
$\bar J_{z,k},\bar J_{c,k}$ in \eqref{eq:m3-identity} by Gauss quadrature along the traversed segments.
Report the three terms of \eqref{eq:m3-identity} (an exact attribution of the excess error) and the
\emph{realised effective gains} $g_z=\mathrm{RMS}\|\bar J_z\Delta z\|/\mathrm{RMS}\|\Delta z\|$,
$g_c=\mathrm{RMS}\|\bar J_c\Delta c\|/\mathrm{RMS}\|\Delta c\|$. The \emph{predicted bound} is
$\widehat{\mathcal B}=\mathrm{NRMSE}_{\rm nom}+\hat g_z\mathrm{RMS}\|\Delta z\|+\hat g_c\mathrm{RMS}\|\Delta c\|$ with
$\hat g_z,\hat g_c$ measured on \emph{other} conditions (or the a priori $\kappa\|B\|\delta_y/\alpha$,
$\sigma_v\|H\|_2$, $\ell_E\delta_u$ in place of the RMS terms).
\end{definition}

With the gains of the same condition, $\widehat{\mathcal B}\ge\mathrm{NRMSE}$ is the triangle
inequality and tests nothing. With transferred gains it is a prediction, and the test is one-sided:
a point above the diagonal falsifies the transferred gain (the decoder is steeper along this
condition's paths); a point below the diagonal confirms nothing. \emph{Status.} This procedure is evaluation-only on
saved checkpoints (the filter, the encoder and the decoder are all available) and has not been run by any node in the
current tree; until it is, the transferred-gain prediction is untested, and this section claims only (i), (ii), (iv) and
the exact filter-level ratios above.

\begin{remark}[what can be certified]\label{rem:m3-cert}
The result is a decomposition into measurable terms, not a certificate. \emph{Exact} (no learning):
$\kappa,\alpha,\gamma_y,\|H\|_2$, the latent deviation of any recorded fault, the path Jacobians on
recorded data, and an upper bound on the SFM gain $\ell_E$ (product of layer norms,
Lemma~\ref{lem:m2-sfm}). \emph{Lower bounds only}: any supremum of Jacobian norms estimated from samples.
Exact Lipschitz constants are NP-hard to compute even for two-layer ReLU networks
\citep{virmaux2018lipschitz}; certified upper bounds (layer-norm products, LipSDP,
\citealp{fazlyab2019efficient}) are valid but loose. On-data Jacobians are the wrong set for sensor
faults: with $n_z>n_x$ the nominal latents lie near an $n_x$-dimensional manifold (times the context
fibre) and $\hat z=\hat z^{\rm nom}+\Delta z$ leaves it, where the decoder's slope is inductive bias; hence
the traversed-segment Jacobians of Definition~\ref{def:m3-measured}. The random-direction quotient
$l_{\rm dec}$ of \texttt{train.phase1\_diagnostics} estimates about $\|J\|_F/\sqrt{n_z}$ on the base
decoder's data (check: $q_{0.95}=0.81\|J\|_2$ for a linear decoder) and is not the constant of the
bound. \emph{Statistical}: the nominal error only through \eqref{eq:m3-quantile}.
\end{remark}

% -------------------------------------------------------------------------------------
% the single \appendix of the document
\appendix
\section{Lemma D: the sampled-data observer}

\begin{lemma}[Lemma D]\label{lem:D}
Let $\|e^{At}\|\le\kappa e^{-\alpha t}$, $\rho_0:=e^{-\alpha dt}$, and let input and context be held on
each $[t_k,t_{k+1})$. Let
$w_0:=\int_0^{dt}\tfrac{\tau}{dt}e^{-\alpha\tau}d\tau=\frac{1-\rho_0(1+\alpha dt)}{\alpha^2dt}$,
$w_1:=\int_0^{dt}(1-\tfrac\tau{dt})e^{-\alpha\tau}d\tau=\frac{\alpha dt-1+\rho_0}{\alpha^2dt}$, so
$w_0+w_1=(1-\rho_0)/\alpha$.
\begin{enumerate}[label=(\alph*)]
\item \emph{Exact step.} For $y$ affine and $g$ constant on the step, $\dot z=Az+By+g$ gives exactly
$z_{k+1}=A_dz_k+W_0y_k+W_1y_{k+1}+W_gg$, with $A_d=e^{Adt}$,
$W_0=\int_0^{dt}(1-\tfrac s{dt})e^{A(dt-s)}B\,ds$, $W_1=\int_0^{dt}\tfrac s{dt}e^{A(dt-s)}B\,ds$,
$W_g=A^{-1}(A_d-I)$ (\texttt{design.latent\_step}).
\item \emph{Interpolation.} With (A3) and the regularity clause, the clean iterate $\bar z_k$ from
$\bar z_0=z(0)$ satisfies
$\sup_k\|\bar z_k-z(t_k)\|\le e_{\rm disc}:=\frac{\kappa\|B\|M_2dt^3}{12(1-\rho_0)}\le\frac{\kappa\|B\|M_2dt^2e^{\alpha dt}}{12\alpha}$,
also for switching inputs, since switches sit on the grid.
\item \emph{Merged kernels.} $\|A_d^mW_0\|\le\kappa\|B\|\rho_0^mw_0$ and $\|A_d^mW_1\|\le\kappa\|B\|\rho_0^mw_1$. Hence
$\sum_i\|G^{(k)}_i\|\le\kappa\|B\|/\alpha$ for every $k$ and $\gamma_y\le\kappa\|B\|/\alpha$ at every $dt$, and
$\gamma_y\to\int_0^\infty\|e^{As}B\|ds$ as $dt\to0$. More generally every sampled constant is bounded by its
continuous counterpart: $\rho_0^M=e^{-\alpha L}$, $c_d\le1/\alpha$ in (e), $\gamma_y\le\kappa\|B\|/\alpha$.
\item \emph{Injection: step-size condition.} For the deployed exponential midpoint step (an
exponential integrator, \citealp{hochbruck2010exponential}, with an explicit Euler predictor)
$z_{k+1}=A_dz_k+W_0y_k+W_1y_{k+1}+W_g\Phi(z^h_k,c_k)$, $z^h_k=z_k+\tfrac{dt}2(Az_k+By_k+\Phi(z_k,c_k))$,
with $\Phi(\cdot,c)$ $\ell_\Phi$-Lipschitz, two solutions driven by the same $(y,c)$ satisfy
$\|z_k-z'_k\|\le\kappa\rho^k\|z_0-z'_0\|$,
\[
\rho:=\rho_0+\kappa(1-\rho_0)\frac{\ell_\Phi\beta}{\alpha},\qquad\beta:=1+\tfrac{dt}2(\|A\|+\ell_\Phi).
\]
So $\kappa\ell_\Phi\beta<\alpha$ suffices; it tends to $\kappa\ell_\Phi<\alpha$
(Corollary~\ref{cor:m1-amplitude}) as $dt\to0$ with $-\ln\rho/dt\to\alpha-\kappa\ell_\Phi$, and $\beta\le1.1$
under $\|A\|dt\le0.1$ (\texttt{design.choose\_dt}). If the stationary latent $\zeta_k=\hat T_0(x(t_k))$ satisfies
$\zeta_{k+1}=(\text{clean step at }\zeta_k)+\int_0^{dt}e^{A(dt-s)}R\,ds+q_k$, with $\|R\|\le\varepsilon_{\rm res}$
the PDE residual and $q_k$ the interpolation/quadrature defect with $\|A_d^mq_k\|\le\kappa\rho_0^m\bar q$, and the
step is fed samples with $\|v_k\|\le\delta_y$, then
\[
\|\hat z_k-\zeta_k\|\le\kappa\rho^k\|\hat z_0-\zeta_0\|+\frac{\kappa\big(\|B\|(1+\ell_\Phi\tfrac{dt}2)\delta_y+\varepsilon_{\rm res}\big)}{\alpha-\kappa\ell_\Phi\beta}+\frac{\kappa\,\alpha\,\bar q}{(1-\rho_0)(\alpha-\kappa\ell_\Phi\beta)},
\]
a single factor $\kappa$, with $\bar q=O(dt^3)$ so the last term is $O(dt^2)$.
\item \emph{Deployed fading memory.} If $u_j=u'_j$ for $j=k-M,\dots,k-1$, the deployed map
$T^d_k(x;u)=\sum_{j\ge0}A_d^j(W_0y_{k-1-j}+W_1y_{k-j})$ satisfies
$\|T^d_k(x;u)-T^d_k(x;u')\|\le2\kappa\|B\|\bar y\,\rho_0^Mc_d$, $c_d:=(w_0+\rho_0w_1)/(1-\rho_0)\le1/\alpha$.
\end{enumerate}
\end{lemma}

\begin{proof}
(a) Variation of constants over one step with the affine $y$ inserted.
(b) $\bar z_k-z(t_k)=-\int_0^{t_k}e^{A(t_k-s)}B(y-\ell)(s)ds$ with $\ell$ the piecewise-linear
interpolant. On $[t_j,t_{j+1}]$ the input is constant, so $y\in C^2$ there, and the Peano form gives
$\|y(t_j+\sigma)-\ell(t_j+\sigma)\|\le\tfrac{\sigma(dt-\sigma)}2M_2$ in any norm. Use
$\|e^{A(t_k-s)}\|\le\kappa\rho_0^{k-1-j}$ on that interval and $\int_0^{dt}\tfrac{\sigma(dt-\sigma)}2d\sigma=\tfrac{dt^3}{12}$,
sum the geometric series, and use $1-\rho_0\ge\alpha dt\,\rho_0$.
(c) $A_d^mW_0=\int_0^{dt}(1-\tfrac s{dt})e^{A(mdt+dt-s)}B\,ds$, so
$\|A_d^mW_0\|\le\kappa\|B\|\rho_0^m\int_0^{dt}(1-\tfrac s{dt})e^{-\alpha(dt-s)}ds=\kappa\|B\|\rho_0^mw_0$, and likewise
for $W_1$. Then $\|H_0\|+\sum_{m\ge1}\|H_m\|\le\kappa\|B\|\big(w_1+\sum_{m\ge1}(\rho_0^{m-1}w_0+\rho_0^mw_1)\big)=\kappa\|B\|\frac{w_0+w_1}{1-\rho_0}=\frac{\kappa\|B\|}\alpha$,
and the finite-$k$ coefficients are a sub-sum of the same terms. The limit holds because
$H_m=e^{At_m}B\,dt+O(dt^2)$ with summable tails.
(d) $\Delta_k=A_d^k\Delta_0+\sum_{j<k}A_d^{k-1-j}W_g(\Phi(z^h_j)-\Phi(z'^h_j))$ with
$\|z^h_j-z'^h_j\|\le\beta\|\Delta_j\|$ and the merged kernel
$\|A_d^{k-1-j}W_g\|=\|\int_0^{dt}e^{A((k-j)dt-s)}ds\|\le\kappa\rho_0^{k-1-j}(1-\rho_0)/\alpha$. With
$a_k=e^{\alpha t_k}\|\Delta_k\|$ and $g=\kappa(e^{\alpha dt}-1)\ell_\Phi\beta/\alpha$, $a_k\le\kappa a_0+\sum_{j<k}ga_j$;
the right side $S_k$ obeys $S_{k+1}\le(1+g)S_k$, so $\|\Delta_k\|\le\kappa(\rho_0(1+g))^k\|\Delta_0\|=\kappa\rho^k\|\Delta_0\|$.
If $\kappa\ell_\Phi<\alpha$ and $\|A\|dt\le0.1$, then $\ell_\Phi dt<\alpha dt\le\|A\|dt$, so $\beta\le1.1$.
For the forced bound, unroll $e_k=\hat z_k-\zeta_k$; every forcing term is merged with the kernel
that propagates it: $\|A_d^{m}(W_0v_j+W_1v_{j+1})\|\le\kappa\rho_0^{m}\|B\|(w_0+w_1)\delta_y$,
$\|A_d^m\int_0^{dt}e^{A(dt-s)}R\,ds\|\le\kappa\rho_0^m(1-\rho_0)\varepsilon_{\rm res}/\alpha$, the predictor mismatch
$\le\beta\|e_j\|+\tfrac{dt}2\|B\|\delta_y$ enters through $A_d^mW_g$, and $\|A_d^mq_j\|\le\kappa\rho_0^m\bar q$. So
$a_k\le\kappa a_0+\sum_{j<k}\big(ga_j+\kappa e^{\alpha(j+1)dt}D\big)$ with
$D=\frac{1-\rho_0}{\alpha}\big(\|B\|(1+\ell_\Phi\tfrac{dt}2)\delta_y+\varepsilon_{\rm res}\big)+\bar q$, hence
$S_{k+1}\le(1+g)S_k+\kappa e^{\alpha(k+1)dt}D$ and $\|e_k\|\le\kappa\rho^k\|e_0\|+\kappa D/(1-\rho)$; insert
$1-\rho=(1-\rho_0)(\alpha-\kappa\ell_\Phi\beta)/\alpha$.
(e) The samples $y_i$ coincide for $i\ge k-M$; bound by (c), summing from $j=M$ ($W_0$) and $j=M+1$
($W_1$), with $\|y_i-y'_i\|\le2\bar y$.
\end{proof}

\emph{Continuous-time interpretation.} The latent step is exact for the held input, the context and
any generated decoder weights are constant on each interval by construction, and the decoder runs
at $t_k$ only. The bounds are about the implemented observer, and the continuous formulas are their
$dt\to0$ limits; the sampled observer adds only $e_{\rm disc}$ and, with injection, the factor $\beta$.

{\sloppy\emph{Check} (\texttt{m3\_checks.py}). (b) \texttt{check\_interpolation}: PRBS input, four step
sizes, order $2.00$, bound/error $3.6$ (\texttt{check\_theory.lemma\_d} uses the same constant, bound/error $4.4$--$4.5$).
(c) \texttt{check\_design\_constants}: $\gamma_y=1.074\le\kappa\|B\|/\alpha=1.491$
($\int\|e^{As}B\|=1.075$); $\gamma_y\le\kappa\|B\|/\alpha$ at $\alpha dt\in\{0.03,0.3,1,3\}$ and for a
non-normal $A$ with $\kappa=1.46$. The $\kappa$ of a non-normal $A$ is estimated by \texttt{spectral\_bounds} on a
$400$-point grid, so those rows are consistency checks, not certificates; the deployed Butterworth $A$ is normal and has
$\kappa=1$ exactly. (d) \texttt{check\_midpoint}: for $\Phi=\ell z$ the per-step gain is
below $\rho$ (e.g.\ $0.99591\le0.99762$ at $\ell=1.3$); 200 random linear injections:
$\max\|\text{step}\|/\rho=0.999$; for a non-normal $A$ ($\kappa=1.99$) $\max_k\|\text{step}^k\|/(\kappa\rho^k)=0.92$, so a single $\kappa$ suffices; $\Phi=-400z$, more stable in continuous time, gives discrete spectral radius $26.4$ at $\ell dt=8.2$,
so the step-size part is needed; the numpy step equals \texttt{LatentStep} to $1.5\times10^{-7}$.
(e) \texttt{m2\_checks.py::lemma\_d\_constants}.\par}

% Bibliography entries (natbib keys):
% bernard2019luenberger: P. Bernard, V. Andrieu. Luenberger observers for nonautonomous nonlinear
%   systems. IEEE TAC 64(1):270-281, 2019. doi:10.1109/TAC.2018.2872202.
% niazi2025kkl: M. U. B. Niazi, J. Cao, M. Barreau, K. H. Johansson. KKL observer synthesis for
%   nonlinear systems via physics-informed learning. arXiv:2501.11655v2, 2025.
% shaaban2026hyperkkl: Y. S. Shaaban, A. S. Sayed, M. U. B. Niazi, K. H. Johansson. HyperKKL:
%   learning KKL observers for non-autonomous nonlinear systems via hypernetwork-based input
%   conditioning. arXiv:2603.29744, 2026.
% reif1999stochastic: K. Reif, S. Guenther, E. Yaz, R. Unbehauen. Stochastic stability of the
%   discrete-time extended Kalman filter. IEEE TAC 44(4):714-728, 1999.
% andersonmoore1979: B. D. O. Anderson, J. B. Moore. Optimal Filtering. Prentice-Hall, 1979.
% hochbruck2010exponential: M. Hochbruck, A. Ostermann. Exponential integrators. Acta Numerica
%   19:209-286, 2010.
% virmaux2018lipschitz: A. Virmaux, K. Scaman. Lipschitz regularity of deep neural networks:
%   analysis and efficient estimation. NeurIPS 2018.
% fazlyab2019efficient: M. Fazlyab, A. Robey, H. Hassani, M. Morari, G. J. Pappas. Efficient and
%   accurate estimation of Lipschitz constants for deep neural networks. NeurIPS 2019.

\bibliographystyle{plainnat}
\bibliography{refs}
\end{document}
